\documentclass[11pt]{amsart}

% ================================================================
% Revision v6 (October 8, 2026): rebuilt on the frozen v5 M3--M9 evidence
% (integrated constrained free-space workflow, repaired mass-weighted POD, matched
% benchmarks, assembly-convention audit, refinement, Fourier and financial studies).
% All numerical tables are \input from tables/, generated by
% scripts/v6_build_manuscript_assets.py from the raw CSVs listed in evidence_v6.json.
% Based on review revision v5 (October 5, 2026).
% ================================================================

\usepackage{amsmath,amssymb,mathtools}
\usepackage{booktabs}
\usepackage{array}
\usepackage{enumitem}
\usepackage{xcolor}
\usepackage[protrusion=true,expansion=false]{microtype}
\usepackage{geometry}
\usepackage{graphicx}
\graphicspath{{figures/}}
\usepackage{bm}
\usepackage{url}
\usepackage{textcomp}
\usepackage{hyperref}
\geometry{margin=1in}
\setlength{\emergencystretch}{1em}
\hypersetup{hidelinks}

\newtheorem{theorem}{Theorem}[section]
\newtheorem{proposition}[theorem]{Proposition}
\newtheorem{lemma}[theorem]{Lemma}
\newtheorem{corollary}[theorem]{Corollary}
\newtheorem{assumption}[theorem]{Assumption}
\newtheorem{definition}[theorem]{Definition}
\theoremstyle{remark}
\newtheorem{remark}[theorem]{Remark}

% ---------------------- provenance callouts -----------------------

% ----------------------------- notation --------------------------
\newcommand{\R}{\mathbb{R}}
\newcommand{\Pset}{\mathcal{P}}
\newcommand{\Sop}{\mathcal{S}}
\newcommand{\Vh}{V_h^0}
\newcommand{\Vmu}{V_{\mu}}
\newcommand{\Vp}[1]{V_{#1}}
\newcommand{\Mh}{M_h}
\newcommand{\Ah}{A_h}
\newcommand{\Sh}{S_h}
\newcommand{\norm}[1]{\left\lVert #1\right\rVert}
\newcommand{\ip}[2]{\left\langle #1,#2\right\rangle}
\newcommand{\dt}{\Delta\tau}
\newcommand{\etaSharp}{\eta_{\mathrm{sh}}}
\newcommand{\etaMass}{\eta_{\mathrm{ms}}}
\newcommand{\etaGuar}{\eta_{\mathrm{gu}}}
\newcommand{\ee}{\mathbf e}
\newcommand{\cc}{\mathbf c}
\newcommand{\rr}{\bm\rho}
\newcommand{\ff}{\mathbf f}
\newcommand{\vv}{\mathbf v}
\newcommand{\bb}{\mathbf b}
\newcommand{\zz}{\mathbf z}
\newcommand{\yy}{\mathbf y}
\newcommand{\ww}{\mathbf w}
\newcommand{\sres}{\mathbf s}
\newcommand{\dd}{\bm\delta}
\newcommand{\rf}{r}
\newcommand{\Hs}{H^{s_\mu}_{M,G}}
\newcommand{\Dm}{\mathcal D_G^-}
\newcommand{\Dp}{\mathcal D_M^+}
\newcommand{\Aop}{\mathcal A_\mu}
\DeclareMathOperator{\Real}{Re}

\title[Residual Certification for CGMY Pricing Trajectories]
      {Residual Certification of Neural and Reduced Surrogates for\\ Tempered-Fractional CGMY Pricing PIDEs}
\author{Davood Damircheli}
\address{Department of Mathematics, Louisiana State University,
Baton Rouge, LA 70803, USA}
\date{}

\subjclass[2020]{65M60, 65M15, 65T50, 68T07, 91G20, 60G51}
\keywords{CGMY process, neural operator, certified scientific machine
learning, B-spline Galerkin method, tempered fractional PIDE, a posteriori
error estimation, residual certification, operator learning, option pricing}

\begin{document}
\begin{abstract}
We study residual certification of learned and reduced European-put pricing
trajectories for the Carr--Geman--Madan--Yor (CGMY) model. A Crank--Nicolson
energy estimate bounds the maximum-in-time error in the mass norm by the
initial error plus an accumulated dual residual norm, and an upper Riesz
bracket preserves this algebraic bound under incomplete iterative solves.
All guarantees are exact-arithmetic statements about the production discrete
problem, evaluated in ordinary floating point. That problem uses quadratic
splines with 62 free and four prescribed coefficients and fractional blocks
assembled by a principal-period FFT rule; the latter differs from exact
Galerkin assembly by 33--41\% in operator norm, concentrated at high
frequencies. We integrate prediction, boundary projection, certification,
correction and recertification in the constrained free coordinates, with the
initial error included explicitly. In the free coordinates the positivity gate
passes at all 459 audited queries, of which the full-coefficient gate rejects 154.
Against a 40-digit reference, 11{,}760 Riesz-bracket comparisons show no
detected violation (8{,}065 resolved, 3{,}695 numerically unresolved). On 120
benchmark queries every corrected neural and reduced trajectory meets the
algebraic tolerance, but the neural workflow is $1.18\times$ (cold) and
$1.58\times$ (reusable setup) slower than a dense direct solve of the same
discrete problem, whereas a repaired mass-weighted POD model of rank 48 is
$0.78\times$ and $0.30\times$. Operator assembly and certification, not the
solve, dominate the cost of every method. Certificates bound the distance to the discrete
price; refinement and independent Fourier prices are reported separately as
continuum diagnostics. The contribution is scheme-matched certification with
explicit scope, not neural acceleration over an optimised deterministic solver.
\end{abstract}

\maketitle

% =================================================================
\section{Introduction}
\label{sec:introduction}
% =================================================================

The CGMY tempered-stable process is a flexible pure-jump L\'evy model whose
four parameters control overall jump intensity, asymmetric exponential
tempering of large jumps, and the fine structure of small jumps
\cite{CarrGemanMadanYor2002}.  In log-price coordinates the European pricing
problem is a nonlocal tempered-fractional parabolic equation built from left-
and right-sided tempered Riemann--Liouville operators of order $Y\in(0,2)$
\cite{CarteadelCastilloNegrete2007}.  In the infinite-variation regime
$Y\in(1,2)$ these operators are strongly nonlocal and the natural spatial
discretisation is dense.

A B-spline Galerkin discretisation of this problem
\cite{DamircheliCGMY2026} motivates the deterministic backbone. Under
interior-supported cardinal-basis and exact-assembly assumptions, the
fractional blocks have Toeplitz structure and FFT matrix actions cost
$O(N_h\log N_h)$ operations. This theoretical construction is distinguished
below from the numerical matrices used in the reported experiments, whose
connection coefficients are computed by a principal-period FFT rule rather
than by the exact Galerkin integral (Remark~\ref{rem:assembly-convention}).
All certificates in this paper are relative to that production discrete problem. The implementation audit identifies $66$ stored coefficients, consisting of
$62$ homogeneous unknowns and four prescribed boundary coefficients.
We derive the corresponding constrained algebraic system below and separate
ideal variational assembly from the numerical matrices used in experiments.

That solver answers the question \emph{how does one solve one CGMY pricing
PIDE accurately?}  The present paper addresses a different question:
\emph{how does one answer thousands or millions of related CGMY queries
quickly without giving up a mathematically meaningful error indicator for each
individual query?}  The distinction matters.  A reliable Galerkin method
remains the reference method; operator learning is introduced to amortise its
cost over a parameter and payoff family, not to replace the underlying
numerical analysis.  Calibration, uncertainty quantification, sensitivity
analysis, repeated pricing, and inverse problems all have exactly this
many-query structure.

Neural operators learn maps between function spaces rather than a single
finite-dimensional regression target.  DeepONet
\cite{LuJinPangZhangKarniadakis2021}, the Fourier Neural Operator (FNO)
\cite{LiKovachkiEtAl2021}, their physics-informed variants
\cite{LiZhengEtAl2021}, and the general framework of
\cite{KovachkiEtAl2023} provide expressive models for parameterised PDE
solution maps; Nonlocal Kernel Networks \cite{YouYuDEliaGaoSilling2022} show
that nonlocal structure can be built directly into a learned architecture.
For pricing, however, raw prediction accuracy is not the only concern.  A
surrogate placed inside an optimiser or a risk engine will be evaluated at
parameter points that the training distribution may represent poorly, and an
aggregate test-set error is not a guarantee for the query in front of the
user.

This motivates a certified viewpoint.  The learned output is placed in the
trusted Galerkin trial space, and the governing discrete variational equation
is used to evaluate a residual in a problem-compliant dual norm.  The
philosophy is that of certified reduced-basis methods
\cite{GreplPatera2005,RozzaHuynhPatera2008,HaasdonkOhlberger2008} and of
recent variationally correct operator learning with a posteriori error
estimation \cite{QiuDahmenChen2026}.  The CGMY family adds two application-specific features: the spatial operator is nonlocal and nonsymmetric, and the
natural fractional energy order $s_\mu=Y/2$ \emph{varies with the parameter}.

The design principle of the paper is
\[
  (\mu,\psi)
  \xrightarrow{\ \text{neural operator}\ }
  \{\cc_\theta^n\}_{n=0}^{N_T}
  \xrightarrow{\ \text{Galerkin residual}\ }
  \eta_\theta
  \xrightarrow{\ \text{tolerance}\ }
  \text{accept or correct},
\]
where $\mu$ collects the CGMY and market parameters and $\psi$ is the payoff.
We write $\psi$ for the payoff throughout and reserve $g$ for the nonlocal
exterior lifting, so that the notation of the operator-learning literature and
the notation of the finite-domain formulation do not collide.

\subsection*{Contributions and scope of the evidence}

The contribution is an application and combination of established
variational certification tools for a nonlocal, nonsymmetric evolution family,
rather than a new general principle of neural-operator certification.
\begin{enumerate}[leftmargin=2em]
\item A compact CGMY family with a parameter-dependent energy order and
explicit parameter-uniform continuity and coercivity constants. The constants
establish uniformity but are too conservative for quantitative online use.
\item A scheme-matched algebraic residual bound for coefficient trajectories,
including the initial error, and a computable bracket for inexact Riesz
evaluation, stated in the constrained free coordinates of the production basis.
\item An integrated, tolerance-checked predict--project--certify--correct
workflow whose acceptance uses the augmented bound
$(\beta_0^2+\eta_{\rm res}^2)^{1/2}$, certifies the exact returned coefficients,
records solver failures, and never accepts on a heuristic indicator.
\item Matched evidence: a 459-query positivity audit, nontrivial-residual
bracket validation, a repaired mass-weighted POD baseline, a 120-query
benchmark with isolated cold and reusable setup, resolution, assembly,
refinement, Fourier and financial studies, each scoped to its population.
The benchmark shows no neural acceleration over a direct solve of the same
discrete problem; the repaired POD baseline is faster at rank 48.
\end{enumerate}

\subsection*{Relation to existing work}

The paper deliberately does \emph{not} claim novelty for neural operators, for
DeepONet or FNO, for residual-based training, or for certified reduced-basis
ideas.  The claim is at the intersection.  Deterministic numerical methods for
L\'evy PIDEs are well developed, including finite differences
\cite{ContVoltchkova2005}, wavelet and finite-element Galerkin methods
\cite{MatacheVonPetersdorffSchwab2004,HilberReichmannSchwabWinter2013}, and
transform methods \cite{FangOosterlee2008}; the companion paper
\cite{DamircheliCGMY2026} supplies the specific B-spline discretisation used
here.  Reduced-basis methods have long attached rigorous residual certificates
to parametric parabolic problems
\cite{GreplPatera2005,RozzaHuynhPatera2008,HaasdonkOhlberger2008}, and
\cite{QiuDahmenChen2026} makes the point---central to this paper---that a
small \emph{arbitrary} PDE residual need not imply a small solution error, so
the residual norm must be inherited from a stable variational formulation
rather than chosen for convenience.  Nonintrusive neural surrogates for
reduced spaces are studied in
\cite{HesthavenUbbiali2018,BhattacharyaEtAl2021}.  The application studied here combines (i) a nonlocal, nonsymmetric, tempered-fractional operator
family whose energy order $Y/2$ itself varies with the parameter, (ii) a
neural output that is conforming in the trusted trial space by construction,
(iii) a dual-norm certificate built from the \emph{discrete} Riesz map already
available inside the solver, and (iv) a workflow in which a computable bound decides each query.
These features must also be compared with certified reduced-basis pricing
methods, including \cite{ContLantosPironneau2011}.

\subsection*{Scope}

The ideal theory and the reported primary experiments concern European pricing in the
infinite-variation regime $Y\in[Y_{\min},Y_{\max}]\Subset(1,2)$ with positive
rates bounded away from zero, which is the regime in which
\cite{DamircheliCGMY2026} proves coercivity directly;
Section~\ref{ssec:shifted} removes the positive-rate restriction at the cost
of an $e^{\sigma T}$ factor.  American options are not treated here; the
obstacle formulation requires primal-feasibility and complementarity residuals
in addition to the PIDE residual and is better handled in a separate study.
The paper does not depend on a new generic neural-operator approximation
theorem or on a statistical generalisation bound.


The remainder is organised as follows.  Section~\ref{sec:model} introduces the
parameterised CGMY problem, freezes the admissible family, and proves uniform
stability.  Section~\ref{sec:galerkin} summarises the deterministic B-spline
backbone and records the two structural facts the certificate needs.
Section~\ref{sec:operator} defines the neural solution operator and the
training losses.  Section~\ref{sec:certificate} is the theoretical
centrepiece: the family of Galerkin residual certificates, their guaranteed
evaluation, their efficiency, and discrete derivative bounds.
Section~\ref{sec:pcc} gives the predict--certify--correct algorithm with its
soundness, contraction, and cost analysis.  Section~\ref{sec:data} describes
dataset generation and implementation, Section~\ref{sec:numerics} reports the available experiments and their limitations, and Sections~\ref{sec:discussion}--\ref{sec:conclusion}
discuss limitations and conclusions.  Appendix~\ref{app:proofs} contains the
longer proofs, Appendix~\ref{app:historical-figures} the superseded historical
figures, and Appendix~\ref{app:evidence} the run identifiers.

\subsection*{Notation}

Table~\ref{tab:notation} collects the symbols used throughout.

\begin{table}[ht]
\centering
\caption{Principal notation.}
\label{tab:notation}
\small
\begin{tabular}{ll}
\toprule
Symbol & Meaning\\
\midrule
$x=\log S$, $\tau=T-t$ & log price and backward time\\
$\mu=(C,G,M,Y,r)$ & CGMY/operator coefficient parameter vector\\
$T$ & maturity; sampled as a learning input and setting the time horizon\\
$\Pset$ & compact admissible coefficient-parameter box, \eqref{eq:parameter-set}\\
$\psi$ & payoff (functional input to the operator)\\
$g$ & prescribed smooth exterior lifting\\
$V=g+U$ & truncated price and homogeneous unknown\\
$s_\mu=Y/2$ & parameter-dependent energy order\\
$\Vmu=\Hs(\Omega)$ & zero-extension energy space\\
$a_\mu(\cdot,\cdot)$ & CGMY bilinear form on $\Omega$\\
$C_{\rm cont},c_1$ & uniform continuity / coercivity constants on $\Pset$\\
$\Vh$ & interior-supported B-spline space, \eqref{eq:Vh}\\
$\Mh,\Ah(\mu),\Sh(\mu)$ & mass, stiffness, symmetric stiffness\\
$\cc_h^n,\cc_\theta^n$ & reference and neural coefficient vectors\\
$\ee^n=\cc_h^n-\cc_\theta^n$ & coefficient error\\
$\rr_\theta^{n+1/2}$ & Crank--Nicolson Galerkin residual, \eqref{eq:cn-residual}\\
$\etaSharp,\etaMass,\etaGuar$ & sharp certificate, mass diagnostic/bound, guaranteed certificate\\
$I_{\rm eff}$ & effectivity index, \eqref{eq:effectivity}\\
\bottomrule
\end{tabular}
\end{table}

% =================================================================
\section{The parameterised CGMY pricing problem}
\label{sec:model}
% =================================================================

\subsection{CGMY backward operator}
\label{ssec:operator}

Let $S_t=e^{X_t}$ under a risk-neutral measure and let $V(x,\tau)$ be the
option value in the log-price variable $x=\log S$ and backward time
$\tau=T-t$.  Following the convention of the companion paper
\cite{DamircheliCGMY2026}, the CGMY parameters satisfy
\begin{equation}
  C>0,\qquad G>0,\qquad M>1,\qquad Y\in(0,2),\qquad Y\ne1,
  \label{eq:cgmy-parameters}
\end{equation}
where $M>1$ is the exponential-moment condition that makes the stock
martingale correction
\begin{equation}
  \omega_S(\mu)
  = C\Gamma(-Y)\bigl[(M-1)^Y-M^Y+(G+1)^Y-G^Y\bigr]
  \label{eq:martingale-correction}
\end{equation}
real valued.  Set
\begin{equation}
  \lambda_0(\mu) := \rf + C\Gamma(-Y)\bigl(M^Y+G^Y\bigr).
  \label{eq:lambda0}
\end{equation}
For $Y\in(1,2)$ one has $\Gamma(-Y)>0$ and hence $\lambda_0>\rf$.  The
tempered left- and right-sided Riemann--Liouville operators are
\begin{equation}
  \Dm u=e^{-Gx}\,{}_{-\infty}D_x^Y\bigl(e^{Gx}u\bigr),
  \qquad
  \Dp u=e^{Mx}\,{}_{x}D_\infty^Y\bigl(e^{-Mx}u\bigr),
  \label{eq:tempered-ops}
\end{equation}
and the single backward-time operator used for both the theory and the
discretisation is
\begin{equation}
  \Aop u
  =-(\rf-\omega_S)u_x
   -C\Gamma(-Y)\Dp u
   -C\Gamma(-Y)\Dm u
   +\lambda_0u.
  \label{eq:cgmy-operator}
\end{equation}
The European pricing problem is
\begin{equation}
  \partial_\tau V+\Aop V=0,
  \qquad
  V(\cdot,0)=\psi.
  \label{eq:pricing-pide}
\end{equation}
On the full line the Fourier symbols are
$\widehat{\Dm u}(\xi)=(G+i\xi)^Y\widehat u(\xi)$ and
$\widehat{\Dp u}(\xi)=(M-i\xi)^Y\widehat u(\xi)$ on the principal branch, so
that
\begin{equation}
  A_\mu(\xi)
  =-(\rf-\omega_S)i\xi
   -C\Gamma(-Y)\bigl[(M-i\xi)^Y+(G+i\xi)^Y\bigr]
   +\lambda_0 .
  \label{eq:symbol}
\end{equation}

Equations \eqref{eq:cgmy-operator}--\eqref{eq:symbol}, the
finite-domain lifting formulation recalled in
Section~\ref{ssec:finite-domain}, and the fixed-parameter coercivity statement
are established in \cite{DamircheliCGMY2026} and are used here without
reproof.

\subsection{The admissible family}
\label{ssec:family}

For the certified operator-learning study we freeze a compact set
\begin{equation}
\begin{aligned}
 \Pset=\bigl\{\mu=(C,G,M,Y,\rf):\;&
 C\in[C_{\min},C_{\max}],\quad C_{\min}>0,\\
 &G\in[G_{\min},G_{\max}],\quad G_{\min}>0,\\
 &M\in[M_{\min},M_{\max}],\quad M_{\min}>1,\\
 &Y\in[Y_{\min},Y_{\max}]\Subset(1,2),\\
 &\rf\in[r_{\min},r_{\max}],\quad r_{\min}>0\bigr\}.
\end{aligned}
 \label{eq:parameter-set}
\end{equation}
The maturity $T$ is the terminal point of the time grid and is sampled as an
additional learning input in the reported data, but it is kept outside
$\mu$ because it changes the time horizon rather than the spatial CGMY
operator coefficients.  The payoff $\psi$ is treated separately as a
functional input.  Two structural requirements on
$\Pset$ deserve emphasis, because they are what make the certificate constants
uniform rather than merely finite for each $\mu$.

First, $M_{\min}>1$ selects an interior exponential-tempering regime.
It is not needed merely to bound the displayed martingale correction:
$(M-1)^Y\to0$ as $M\downarrow1$, so that correction remains finite
and continuous at $M=1$ for the stated compact $Y$ interval.  Second, $Y$ must stay away from both endpoints of
$(1,2)$: as $Y\downarrow1$ the coercivity mechanism used in
Theorem~\ref{thm:uniform-coercivity} degenerates because the high-frequency
phase $Y\arctan(|\xi|/a)$ no longer exceeds $\pi/2$ uniformly, and $\Gamma(-Y)$ blows up at both $Y\downarrow1$ and $Y\uparrow2$.  The present analysis is restricted to this interior infinite-variation
regime; no claim of coverage of all calibrated CGMY models is made.

A single truncated domain $\Omega=(x_{\min},x_{\max})$ must be admissible for
every $\mu\in\Pset$.  Since the tail heuristic of \cite{DamircheliCGMY2026}
suggests left and right buffers of order $10/G$ and $10/M$, the domain is
dictated by $G_{\min}$ and $M_{\min}$; widening the $G$ range therefore
enlarges $\Omega$ and, at fixed $h$, increases $N_h$.  This coupling between the parameter box and the mesh is therefore reported
explicitly in the experimental setup.

\begin{table}[ht]
\centering
\caption{Parameter ranges reported for the fixed-strike study. The earlier box was
narrowed after a matrix-stability audit; the cause of the gate failures has
not been established from the supplied source and figures.}
\label{tab:proposed-box}
\small
\begin{tabular}{llll}
\toprule
Parameter & Production range & Sampling & Comment\\
\midrule
$C$    & $[0.05,1.0]$ & Sobol & narrowed from $0.01$ after the SPD audit\\
$G$    & $[1.0,5.0]$  & Sobol & unchanged\\
$M$    & $[2.0,9.0]$  & Sobol & narrowed from $20$ after the SPD audit\\
$Y$    & $[1.10,1.80]$& Sobol & bounded away from $1$ and $2$\\
$\rf$  & $[0.01,0.10]$& Sobol & positive-rate regime\\
$T$    & $[0.1,1.0]$  & Sobol & years\\
$K$    & $100$ & fixed in the main study & varying-strike study uses $[70,150]$\\
\bottomrule
\end{tabular}
\end{table}

The reported mesh uses $\Omega=(\log 100-8,\log 100+8)$ and $64$ elements,
with $66$ coefficient degrees of freedom. The original experiment description
calls the splines cubic, but its relation to the order-$p$ cardinal convention
in Section~\ref{ssec:spline-space} is unresolved. The varying-strike study
uses the same domain. Separation of the payoff kink from the boundary is
not itself a domain-truncation error estimate; domain and mesh refinement
against an independent price reference are required.

\subsection{Finite-domain exterior-value formulation}
\label{ssec:finite-domain}

A nonlocal operator cannot be truncated by prescribing pointwise Dirichlet
values at $x_{\min}$ and $x_{\max}$, so we retain the exterior lifting
formulation of \cite{DamircheliCGMY2026}.  On $\Omega=(x_{\min},x_{\max})$
write
\begin{equation}
  V(x,\tau;\mu,\psi)=g(x,\tau;\mu,\psi)+U(x,\tau;\mu,\psi),
  \label{eq:lifting-split}
\end{equation}
where $g$ is a prescribed smooth function carrying the far-field data and $U$
is the homogeneous unknown, extended by zero outside $\Omega$.  For a
European put the companion paper uses the explicit lifting
\begin{equation}
  g(x,\tau)=\bigl(Ke^{-\rf\tau}-e^{x}\bigr)\zeta(x)
  \label{eq:lifting-put}
\end{equation}
with $\zeta$ a fixed smooth cutoff equal to one on $(-\infty,x_{\min}]$ and to
zero on $[x_{\max},\infty)$.  The term $-e^x$ is required for compatibility
with the payoff at $\tau=0$.  Crucially, $g$ is prescribed rather than
computed, so the discrete and continuous liftings coincide and the two lifting
discrepancies $\Delta_g^{\rm sol}$ and $\Delta_g^{\rm frc}$ of
\cite{DamircheliCGMY2026} vanish identically.

With $s_\mu:=Y/2\in(1/2,1)$ the energy space is the zero-extension space
\begin{equation}
  \Vmu=\Hs(\Omega)
  =\bigl\{v\in L^2(\Omega):\widetilde v\in H^{s_\mu}(\R)\bigr\},
  \qquad
  \norm{v}_{\Vmu}:=\norm{\widetilde v}_{H^{s_\mu}_{M,G}(\R)},
  \label{eq:energy-space}
\end{equation}
where $\widetilde v$ is extension by zero and the full-line norm is defined by
\begin{equation}
  \norm{u}_{H^{q}_{M,G}(\R)}^2
  :=\frac1{2\pi}\int_\R
  \Bigl(1+|M-i\xi|^{2q}+|G+i\xi|^{2q}\Bigr)|\widehat u(\xi)|^2\,d\xi .
  \label{eq:HMG-norm}
\end{equation}
The finite-domain bilinear form is
$a_\mu(u,w):=\langle\Aop\widetilde u,\widetilde w\rangle$, and the truncated
weak problem is: find $U\in L^2(0,T;\Vmu)$ with
$\partial_\tau U\in L^2(0,T;\Vmu')$, $U(\cdot,0)=\psi-g(\cdot,0)$, and
\begin{equation}
  (\partial_\tau U,w)_{L^2(\Omega)}+a_\mu(U,w)=\ell^g(w;\tau)
  \qquad\forall\,w\in\Vmu,
  \label{eq:weak-problem}
\end{equation}
where the lifting functional is
\begin{equation}
  \ell^{g}(w;\tau)
  :=-\int_\Omega\partial_\tau g\,w\,dx-\int_\Omega(\Aop g)\,w\,dx .
  \label{eq:lifting-functional}
\end{equation}
All exterior data enter only through $\Aop g$ evaluated on $\Omega$.

The experimental payoff family is the European-put family
\[
  \Psi_{\rm put}=\{\psi_K(x)=(K-e^x)^+ : K\in[70,150]\}.
\]
The main fixed-strike comparison uses the single member $K=100$ in order to
isolate variation in $(C,G,M,Y,r,T)$, while the operator-generalisation study
samples $K$ jointly with the CGMY parameters.  For each strike, the exterior
lifting is the corresponding put far-field lifting used by the deterministic
solver.  Thus the certification argument is applied after projection to the
same finite-dimensional conforming space; we do not require or claim a
uniform continuum regularity theorem for arbitrary payoff families.

\subsection{Parameter-uniform continuity and coercivity}
\label{ssec:uniform}

The companion paper proves, for each fixed admissible $\mu$ with $Y\in(1,2)$,
continuity of $a_\mu$ and---when $\rf>0$---coercivity, together with the
pointwise symbol bound $\Real A_\mu(\xi)\ge \rf$.  For an operator-learning
paper the fixed-parameter statement is not enough: the constants must be
controlled over the whole family, both because a single trained model is used
across $\Pset$ and because the efficiency constant of
Theorem~\ref{thm:efficiency} is exactly the ratio of these constants.  The following theorem supplies the required family-level bound.

\begin{theorem}[Parameter-uniform continuity and coercivity]
\label{thm:uniform-coercivity}
Let $\Pset$ be as in \eqref{eq:parameter-set}.  Define
\begin{equation}
  a_\ast:=\max\{M_{\max},G_{\max}\},
  \qquad
  \delta_\ast:=\tfrac{\pi}{4}\bigl(Y_{\min}-1\bigr)>0,
  \qquad
  R_\ast:=a_\ast\tan\!\Bigl(\frac{\pi/2+\delta_\ast}{Y_{\min}}\Bigr).
  \label{eq:uniform-constants-aux}
\end{equation}
Then there exist constants $C_{\rm cont}<\infty$ and $c_1>0$, depending only
on $\Pset$, such that for every $\mu\in\Pset$ and all $v,w\in\Vmu$,
\begin{equation}
  |a_\mu(v,w)|\le C_{\rm cont}\norm{v}_{\Vmu}\norm{w}_{\Vmu},
  \qquad
  a_\mu(v,v)\ge c_1\norm{v}_{\Vmu}^2 .
  \label{eq:uniform-coercivity}
\end{equation}
Explicit admissible choices are
\begin{equation}
  c_1
  =\frac{1}{1+2\Theta}
  \min\left\{
    \frac{r_{\min}}{(1+R_\ast^2)^{Y_{\max}/2}},\;
    2C_{\min}\gamma_{\min}\sin(\delta_\ast)
    \Bigl(\frac{R_\ast^2}{1+R_\ast^2}\Bigr)^{Y_{\max}/2}
  \right\},
  \label{eq:c1-explicit}
\end{equation}
and
\begin{equation}
  C_{\rm cont}
  =2\Bigl(
    \Omega_\ast
    +2C_{\max}\gamma_{\max}\max\{1,a_\ast^2\}^{Y_{\max}/2}
    +\lambda_\ast
  \Bigr),
  \label{eq:Ccont-explicit}
\end{equation}
where
\begin{equation}
\begin{gathered}
\gamma_{\min}:=\!\!\min_{Y\in[Y_{\min},Y_{\max}]}\!\!\Gamma(-Y)>0,
\qquad
\gamma_{\max}:=\!\!\max_{Y\in[Y_{\min},Y_{\max}]}\!\!\Gamma(-Y)<\infty,
\\
\Theta:=\max\{1,M_{\max}^2,G_{\max}^2\}^{Y_{\max}/2},
\qquad
\Omega_\ast:=\sup_{\mu\in\Pset}|\rf-\omega_S(\mu)|,
\qquad
\lambda_\ast:=\sup_{\mu\in\Pset}\lambda_0(\mu),
\end{gathered}
\label{eq:uniform-aux-constants}
\end{equation}
the last two being finite.
\end{theorem}

\begin{proof}
See Appendix~\ref{app:proof-uniform}.
\end{proof}

\begin{remark}[Why $\Pset$ must be compact and interior]
\label{rem:why-compact}
Each hypothesis in \eqref{eq:parameter-set} is used exactly once in the proof.
$Y_{\min}>1$ produces $\delta_\ast>0$ and hence a strictly negative
high-frequency cosine; $Y_{\min}>1$ and $Y_{\max}<2$ keep $\gamma_{\max}$ finite; $C_{\min}>0$
keeps the fractional dissipation from vanishing; $r_{\min}>0$ controls the
zero mode, where the fractional part contributes nothing; $M_{\min}>1$ selects the interior model regime, while boundedness of
$\Omega_\ast$ follows from continuity on the compact box; and $G_{\max},M_{\max}<\infty$ bound the norm
equivalence.  Dropping $r_{\min}>0$ costs coercivity but not the theory: see
Section~\ref{ssec:shifted}.
\end{remark}

\begin{remark}[Quantitative size of the explicit constants]
\label{rem:constants-computable}
The extrema in \eqref{eq:uniform-aux-constants} are finite-dimensional
optimisation problems on compact sets. Numerical evaluation on the reported
box gives $\gamma_{\min}\approx2.3024$, $\gamma_{\max}\approx9.7148$,
$R_\ast\approx125.8364$, $\Theta\approx52.1959$,
$\Omega_\ast\approx23.9770$ and $\lambda_\ast\approx224.2716$.
Consequently,
\[
 c_1\approx1.576\times10^{-8},\qquad
 C_{\rm cont}\approx2.525\times10^3,\qquad
 2(C_{\rm cont}/c_1)^2\approx5.134\times10^{22}.
\]
These approximate numerical evaluations illustrate the conservatism of the
analytic choices; they are not validated enclosures of the extrema. The
theorem establishes qualitative uniformity. Quantitative discrete efficiency
is assessed using $\kappa_h$, not this large upper bound.
\end{remark}

\subsection{The solution operator and its parameter regularity}
\label{ssec:solution-operator}

The neural operator approximates
\begin{equation}
  \Sop:(\mu,T,\psi)\longmapsto U|_{[0,T]}(\cdot,\cdot;\mu,\psi),
  \label{eq:solution-operator}
\end{equation}
and any statement about how well it can be approximated presupposes that
$\Sop$ is regular.  The difficulty specific to CGMY is that the energy space
$\Vmu$ itself moves with $Y$, so a naive Lipschitz statement in ``the'' energy
norm is not even well posed.  We separate two claims of very different status.

The first is discrete, needs no new analysis, and is proved by exactly the
same energy argument that produces the certificate.  It is the statement that
matters for training, because the network is trained on discrete
trajectories.

\begin{proposition}[Discrete parameter stability of the solution map]
\label{prop:discrete-lipschitz}
Fix a mesh, a time grid, and a payoff $\psi$. Assume that the first
assembled symmetric matrix $\Sh(\mu_1)$ is positive definite, and let
$\mu_1,\mu_2\in\Pset$
with $\cc_h(\mu_1),\cc_h(\mu_2)$ the corresponding Crank--Nicolson
trajectories of \eqref{eq:cn-reference}.  Then, for every
$m=1,\dots,N_T$,
\begin{equation}
\begin{aligned}
  \norm{\cc_h^m(\mu_1)-\cc_h^m(\mu_2)}_{\Mh}^2
  \le{}& \norm{\cc_h^0(\mu_1)-\cc_h^0(\mu_2)}_{\Mh}^2\\
  &+\dt\sum_{n=0}^{m-1}
  \bigl\lVert\bm\varrho^{\,n+1/2}\bigr\rVert_{\Sh(\mu_1)^{-1}}^2 ,
\end{aligned}
  \label{eq:discrete-lipschitz}
\end{equation}
where the parameter-perturbation residual is
\begin{equation}
  \bm\varrho^{\,n+1/2}
  :=\bigl(\Ah(\mu_1)-\Ah(\mu_2)\bigr)\overline\cc_h^{\,n+1/2}(\mu_2)
   -\bigl(\overline\ff^{\,n+1/2}(\mu_1)-\overline\ff^{\,n+1/2}(\mu_2)\bigr).
  \label{eq:param-residual}
\end{equation}
Under exact assembly of the conforming variational matrices,
Theorem~\ref{thm:uniform-coercivity} additionally gives
\begin{equation}
  \norm{\bigl(\Ah(\mu_1)-\Ah(\mu_2)\bigr)\vv}_{\Sh(\mu_1)^{-1}}
  \le
  \frac{1}{c_1}\,
  \sup_{0\ne v_h,w_h\in\Vh}
  \frac{|a_{\mu_1}(v_h,w_h)-a_{\mu_2}(v_h,w_h)|}
       {\norm{v_h}_{\Vp{\mu_1}}\norm{w_h}_{\Vp{\mu_1}}}\;
  \norm{\vv}_{\Sh(\mu_1)},
  \label{eq:matrix-difference}
\end{equation}
so a bound on the difference of the two bilinear forms in a \emph{single}
norm---that of $\mu_1$---suffices, and no common energy space is needed.
\end{proposition}

\begin{proof}
The difference $\ee^n:=\cc_h^n(\mu_1)-\cc_h^n(\mu_2)$ satisfies the
Crank--Nicolson error equation \eqref{eq:discrete-error-equation} with
$\rr_\theta^{n+1/2}$ replaced by $\bm\varrho^{\,n+1/2}$ of
\eqref{eq:param-residual}, since $\cc_h(\mu_2)$ solves the $\mu_2$ scheme
exactly.  Apply Theorem~\ref{thm:fully-discrete-certificate} with
$\Sh=\Sh(\mu_1)$.  Estimate \eqref{eq:matrix-difference} follows from
Lemma~\ref{lem:kappa} applied to the bilinear form
$a_{\mu_1}-a_{\mu_2}$.
\end{proof}

\begin{remark}
Proposition~\ref{prop:discrete-lipschitz} is a small observation with a useful
consequence: the certificate machinery of Section~\ref{sec:certificate} is not
only an a posteriori device for neural predictions, it is also the natural
tool for quantifying how the discrete solution manifold varies over $\Pset$,
and therefore for diagnosing which regions of $\Pset$ are intrinsically hard
to learn.
\end{remark}

We do not assert a continuum Lipschitz theorem with respect to $Y$.  Such a
statement is delicate because the energy space itself varies with
$s_\mu=Y/2$ and differentiation of the symbol introduces logarithmic growth.
The discrete stability result above is sufficient for the learning and
certification arguments used in this paper; a common-space continuum
regularity theory is therefore left as a separate analytical question.

% =================================================================
\section{The trusted B-spline Galerkin backbone}
\label{sec:galerkin}
% =================================================================

This section records the ideal conforming construction motivated by
\cite{DamircheliCGMY2026}. Its exact Galerkin form and Toeplitz action
are stated explicitly. The reported numerical basis has not yet been
identified with this construction; the algebraic certification results are
therefore kept separate from that identification.

\subsection{Conforming spline space and Galerkin matrices}
\label{ssec:spline-space}

The algebraic results use any linearly independent basis
$\{\phi_{h,k}\}_{k\in J_h}$, its mass matrix and the numerical evolution
matrix. Reading them as conforming Galerkin results additionally requires
$V_h^0\subset V_\mu$ and exact restriction of $a_\mu$.

For the ideal cardinal construction, let $N_p$ denote a cardinal B-spline
of \emph{order} $p$, hence degree $p-1$ and support length $p$ in reference
coordinates. Set $h=(x_{\max}-x_{\min})/N$ and
$\phi_{h,k}(x)=h^{-1/2}N_p((x-x_{\min})/h-k)$. Then
\begin{equation}
 V_h^0=\operatorname{span}\{\phi_{h,k}:k=0,\ldots,N-p\},
 \qquad N_h=N-p+1.
 \label{eq:Vh}
\end{equation}
These are the translates supported within the closed domain; their zero
extensions have the required regularity. All exterior data enter through the
lifting. For $N=64$, order $p=3$ gives $62$ quadratic-spline coefficients;
cubic splines in this convention have order $p=4$ and give $61$ coefficients.
The production implementation stores the boundary-crossing translates as
well: $N+p-1=66$ raw coefficients, of which $2(p-1)=4$ are prescribed.
Thus its free dimension is exactly $N-p+1=62$.

\begin{remark}[Identification of the production implementation]
\label{rem:basis-identification}
The basis audit measures support width three and confirms that the code's
$p=3$ denotes order, hence degree two. The active indices are
$k=-(p-1),\ldots,N-1$. Prescribing the first and last $p-1$ coefficients
leaves precisely the interior-supported translates in \eqref{eq:Vh}.
These free splines have continuous, piecewise-polynomial zero extensions
in $H^1(\mathbb R)$ and therefore in $H^{Y/2}(\mathbb R)$ for $Y/2<1$.
Boundary-crossing translates represent the prescribed lifting, rather than
homogeneous trial or test functions. The numerical full coefficient arrays
must consequently be distinguished from the free-space matrices in the
variational analysis.
\end{remark}

\begin{lemma}[Conformity and exactness; {\cite[Lem.~4.1]{DamircheliCGMY2026}}]
\label{lem:conformity}
Let $p\ge2$ and $s_\mu=Y/2\in(1/2,1)$.  Then $\Vh\subset\Vmu$, every element
of $\Vh$ coincides with its own zero extension, and for all $v_h,w_h\in\Vh$
\begin{equation}
  (\Aop v_h,w_h)_{L^2(\Omega)}=a_\mu(v_h,w_h)=a_\R(\widetilde v_h,\widetilde w_h).
  \label{eq:conformity}
\end{equation}
Consequently the homogeneous fractional stiffness blocks are exactly the
principal submatrices of the full-line Toeplitz connection-coefficient
matrices, and no boundary corrections arise.
\end{lemma}

Lemma~\ref{lem:conformity} connects the ideal algebraic matrices with the
continuous bilinear form. The algebraic certificates themselves need only
a positive-definite symmetric part and a common mass matrix; identification
with the continuous CGMY problem additionally needs this assembly connection.  With
$U_h(x,\tau)=\sum_kc_k(\tau)\phi_{h,k}(x)$ set
\begin{equation}
 (\Mh)_{kk'}=(\phi_{h,k'},\phi_{h,k})_{L^2(\Omega)},
 \qquad
 (\Ah(\mu))_{kk'}=a_\mu(\phi_{h,k'},\phi_{h,k}),
 \label{eq:matrices}
\end{equation}
so that
\begin{equation}
 \Ah(\mu)
 =-(\rf-\omega_S)D_h
  -C\Gamma(-Y)K_M^+
  -C\Gamma(-Y)K_G^-
  +\lambda_0\Mh,
 \label{eq:Ah}
\end{equation}
where $D_h$ is the first-derivative matrix and the fractional blocks are
Toeplitz with the $h^{-Y}$ scaling
\begin{equation}
 (K_G^-)_{kk'}=h^{-Y}\Lambda_G^-(k-k';Gh),
 \qquad
 (K_M^+)_{kk'}=h^{-Y}\Lambda_M^+(k-k';Mh),
 \label{eq:toeplitz-blocks}
\end{equation}
the tempered connection coefficients $\Lambda_G^-,\Lambda_M^+$ being given by
absolutely convergent Fourier-symbol integrals.  The matrices $\Mh$ and $D_h$
are banded; the Toeplitz blocks are applied by circulant embedding in
$O(N_h\log N_h)$ operations with no compression error.  The semidiscrete
system is
\begin{equation}
  \Mh\dot\cc(\tau)+\Ah(\mu)\cc(\tau)=\ff(\tau;\mu,\psi),
  \qquad
  f_k(\tau)=\ell^{g}(\phi_{h,k};\tau).
  \label{eq:semidiscrete-system}
\end{equation}

The connection coefficients, their closed forms in the untempered
limit, the one-sided support and tempered algebraic decay, the Toeplitz
structure, and the FFT realisation are established in
\cite{DamircheliCGMY2026}; only \eqref{eq:toeplitz-blocks} and the
$O(N_h\log N_h)$ application cost are used below.

\begin{remark}[Production assembly convention]
\label{rem:assembly-convention}
By Plancherel, the exact Galerkin connection coefficient is the real-line
integral
$\Lambda(m;\delta)=\frac1{2\pi}\int_{\R}|\widehat N_p(\vartheta)|^2
S(\vartheta;\delta)e^{\mathrm i\vartheta m}\,d\vartheta$.
The production solver instead samples the same integrand at $L=2N_h$ points of
the principal period $[-\pi,\pi)$ and applies an inverse FFT, i.e.\ a
Riemann sum of the principal-period integral. We keep this production
assembly as the \emph{definition} of the discrete problem: every residual,
certificate, correction and benchmark in Section~\ref{sec:numerics} uses it.
Two separate measurements quantify it (Table~\ref{tab:assembly}). Against
exact 40-digit quadrature of the same principal-period integral, with
independently re-derived $M_h$ and $D_h$, the implementation of the stated
convention is confirmed: entrywise differences decrease from about
$10^{-3}$ at $N=64$ to $10^{-5}$ at $N=512$, and $\lambda_{\min}(S_I)$
agrees to $4\times10^{-7}$ or better. Against the exact real-line Galerkin
coefficients the deviation is $33$--$41\%$ in operator 2-norm and about
$9\%$ entrywise, does not decrease under refinement, and is concentrated at
high frequencies; its effect on $\lambda_{\min}(S_I)$ decreases from
$9\times10^{-4}$ (relative) at $N=64$ to $10^{-4}$ at $N=512$. Hence
statements below that assume exact Galerkin assembly
(Lemmas~\ref{lem:Sh-spd} and~\ref{lem:mass-lower}, the $h$-uniform bound in
Lemma~\ref{lem:kappa}, Corollary~\ref{cor:total-error}) describe the ideal
scheme, not the assembled matrices. The algebraic certificates do not need
exact assembly: they hold for the assembled matrices whenever the symmetric
part of the free block is positive definite. The observed convergence of the
production prices to independent Fourier prices under refinement
(Section~\ref{ssec:refinement}) is empirical evidence that the
frequency-truncated scheme is consistent, not a proof.
\end{remark}

\subsection{Prescribed coefficients and the algebraic trial space}
\label{ssec:constrained-system}

Write $\bar M,\bar A$ for the raw $(N+p-1)$-dimensional production matrices,
and let $E$ inject the $N-p+1$ free coefficients into the raw vector.
Let $b^n$ vanish on the free entries and carry the prescribed entries on the
boundary. Every admissible trajectory has the form
\[
 c^n=Ez^n+b^n,\qquad
 M_I=E^\top\bar M E,\quad A_I=E^\top\bar A E,\quad
 S_I=\tfrac12(A_I+A_I^\top)=E^\top\bar S E,
 \quad \bar S=\tfrac12(\bar A+\bar A^\top).
\]
For the implementation, the left prescribed coefficients equal
$K e^{-r\tau_n}\sqrt h$ and the right ones equal zero. These are
coefficient constraints; they are not four extra PIDE equations.
Restricting the CN equations to the free rows gives
\begin{equation}
 M_I\frac{z^{n+1}-z^n}{\Delta\tau}
 +A_I\frac{z^{n+1}+z^n}{2}=f_I^{n+1/2},\qquad
 f_I^{n+1/2}=E^\top\!\left[\bar f^{n+1/2}
 -\bar M\frac{b^{n+1}-b^n}{\Delta\tau}
 -\bar A\frac{b^{n+1}+b^n}{2}\right].
 \label{eq:constrained-cn}
\end{equation}
Even when the stored raw forcing $\bar f$ is zero, the effective free forcing
generally is not. For a candidate with the identical $b^n$, its residual is
\begin{equation}
 \rho_I^{n+1/2}=E^\top\!\left[\bar M\frac{c_\theta^{n+1}-c_\theta^n}{\Delta\tau}
 +\bar A\frac{c_\theta^{n+1}+c_\theta^n}{2}-\bar f^{n+1/2}\right].
 \label{eq:constrained-residual}
\end{equation}
Thus extracting the free rows of the raw residual is exactly equivalent to
forming the lifted free-space residual. Boundary discrepancies must first be
removed by projection or separately bounded; zeroing their residual entries
alone does not control them. The error then satisfies $e^n=Ee_I^n$ and
$(e^n)^\top\bar M e^n=(e_I^n)^\top M_I e_I^n$.
The reliability, efficiency and inexact-Riesz results below apply directly to
$(M_I,A_I,S_I,\rho_I)$ when $S_I\succ0$. The integrated implementation used
for every result in Section~\ref{sec:numerics} works in exactly these
coordinates: it projects a candidate onto the prescribed coefficients, forms
$\rho_I$ with the stored forcing and lifting, factors the principal block
$M_I$ itself (not a restriction of $\bar M^{-1}$), and applies $A_I$, $S_I$,
the preconditioners and the positivity gate through inject--apply--restrict
adapters of the single production operator.

\begin{proposition}[Conservatism of the legacy full-space certificate]
\label{prop:legacy-conservative}
Suppose $\bar S\succ0$, and the reference and candidate have identical
prescribed coefficients. Then $S_I\succ0$ and
\begin{equation}
 \rho_I^\top S_I^{-1}\rho_I
 \le (E\rho_I)^\top\bar S^{-1}(E\rho_I).
 \label{eq:legacy-conservative}
\end{equation}
Consequently a valid upper Riesz bracket for the zero-padded free residual
using $\bar S$ is also a conservative reliability bound for the constrained
problem, subject to the same initial-error and arithmetic assumptions.
\end{proposition}
\begin{proof}
The dual quadratic form equals
$\sup_v(2r^\top v-v^\top S v)$. Restricting the full-space supremum to
$v=Ew$ yields the free-space supremum, proving
\eqref{eq:legacy-conservative}. Positive definiteness restricts to every
principal subspace.
\end{proof}

The converse positivity implication fails: $S_I$ may be positive while
$\bar S$ is indefinite on prescribed directions. The legacy gate can
therefore reject a mathematically certifiable constrained query.
Replacing only that gate while retaining $\bar S$ in the Riesz solve is
invalid; the integrated workflow replaces the Riesz action, preconditioners,
spectral bound, residual coordinates and graph-norm operators together.
The legacy full-space pipeline is retained only for the historical results of
Section~\ref{ssec:historical} and for comparison in Table~\ref{tab:free-spd-audit}.
A direct computation with exact real-line coefficients shows that the full
$\bar S$ remains indefinite on the two legacy-failure queries examined, while
$\lambda_{\min}(S_I)$ changes by less than $10^{-3}$ relative; the legacy
failures are therefore not an artefact of the assembly convention of
Remark~\ref{rem:assembly-convention}.

\subsection{Reference time integrators}
\label{ssec:reference-time}

Let $\tau_n=n\dt$, $n=0,\dots,N_T$, and write
$\overline\vv^{\,n+1/2}=\tfrac12(\vv^{n+1}+\vv^{n})$.  The Crank--Nicolson
reference trajectory solves
\begin{equation}
 \Mh\frac{\cc_h^{n+1}-\cc_h^n}{\dt}
 +\Ah(\mu)\overline\cc_h^{\,n+1/2}
 =\overline\ff^{\,n+1/2}.
 \label{eq:cn-reference}
\end{equation}
The production solver uses pure CN from the initial payoff, with no startup
smoothing. As a theoretical extension, Rannacher start-up
\cite{Rannacher1984} may be used: the first $2\nu$ steps are taken with backward Euler at
step size $\dt/2$ (equivalently, $\nu$ pairs of half steps), after which
\eqref{eq:cn-reference} is used.  A backward-Euler step reads
\begin{equation}
 \Mh\frac{\cc_h^{n+1}-\cc_h^n}{\dt_n}
 +\Ah(\mu)\cc_h^{n+1}
 =\ff^{n+1}.
 \label{eq:be-reference}
\end{equation}

\begin{remark}[Certify the scheme you actually ran]
\label{rem:certify-what-you-ran}
A residual formed from \eqref{eq:cn-reference} does not certify a trajectory
produced by \eqref{eq:be-reference} and vice versa.  This is not a technicality:
the two residuals differ at $O(\dt)$, which is the same order as the
Rannacher damping itself.  Section~\ref{ssec:be-certificate} therefore derives
the backward-Euler certificate explicitly, and
Corollary~\ref{cor:rannacher} concatenates the two so that a Rannacher-smoothed
reference trajectory is certified end to end.
\end{remark}

\subsection{The deterministic estimate used in the total-error split}
\label{ssec:deterministic-error}

For the truncated physical price $V=g+U$, under the regularity, lifting, and
boundary-buffer hypotheses of \cite{DamircheliCGMY2026}, the fully discrete
Crank--Nicolson estimate reads
\begin{equation}
 \max_{0\le n\le N_T}\norm{V(\cdot,\tau_n)-V_h^n}_{L^2(\Omega)}
 \le C_\mu\bigl(
 h^{p-Y/2}+\dt^2+\varepsilon_b+\Delta_g^{\rm frc}+\Delta_g^{\rm sol}
 \bigr),
 \label{eq:deterministic-error}
\end{equation}
and for the prescribed smooth lifting \eqref{eq:lifting-put} one has
$\Delta_g^{\rm frc}=\Delta_g^{\rm sol}=0$.

We use \eqref{eq:deterministic-error} only to separate deterministic
discretisation error from surrogate error.  All numerical certificates
reported below are relative to the trusted Galerkin trajectory; no fully
computable continuum-price certificate is claimed.

\subsection{The lifting cancels}
\label{ssec:lifting-cancels}

One elementary observation is used repeatedly and is worth isolating, because
it is what allows a certificate for the homogeneous unknown to be quoted as a
certificate for the price.

\begin{lemma}[Exact cancellation of the lifting]
\label{lem:lifting-cancels}
Let $V_h^n=g(\cdot,\tau_n)+U_h^n$ and
$V_{\theta,h}^n=g(\cdot,\tau_n)+U_{\theta,h}^n$ use the \emph{same} prescribed
lifting.  Then, for every $n$ and every $q\in[1,\infty]$,
\begin{equation}
  V_h^n-V_{\theta,h}^n=U_h^n-U_{\theta,h}^n,
  \qquad
  \partial_x^{(j)}\bigl(V_h^n-V_{\theta,h}^n\bigr)
  =\partial_x^{(j)}\bigl(U_h^n-U_{\theta,h}^n\bigr),
  \quad j=1,2 .
  \label{eq:lifting-cancels}
\end{equation}
In particular the surrogate error in the price and in every spatial derivative
is exactly the surrogate error in the homogeneous unknown, with no lifting
contribution.
\end{lemma}

\begin{proof}
Immediate from \eqref{eq:lifting-split} and linearity of differentiation; the
lifting is prescribed and identical in both objects, so it subtracts out.
\end{proof}

Lemma~\ref{lem:lifting-cancels} is specific to the design choice that the
neural operator predicts the \emph{homogeneous} coefficients and inherits the
lifting analytically.  A surrogate trained to predict the full price directly
would not enjoy it, and its certificate would inherit an additional lifting
consistency term.

% =================================================================
\section{The conforming neural solution operator}
\label{sec:operator}
% =================================================================

\subsection{Operator-learning target}
\label{ssec:target}

The learning target maps a parameter vector and a payoff \emph{function} to
the Galerkin coefficient trajectory:
\begin{equation}
 \boxed{
 (\mu,T,\psi,\tau)
 \longmapsto
 \cc_\theta(\tau;\mu,T,\psi)\in\R^{N_h}
 \longmapsto
 U_{\theta,h}(\tau,\cdot;\mu,\psi)=\sum_{k}c_{\theta,k}(\tau)\phi_{h,k}\in\Vh,}
 \label{eq:operator-map}
\end{equation}
with physical prediction $V_{\theta,h}=g+U_{\theta,h}$.  The network never has
to learn the far-field asymptotics, which are known analytically.  A simpler
prototype fixes the payoff family and learns $(\mu,T,\tau)\mapsto\cc_\theta$;
the varying-strike experiment of Section~\ref{ssec:operator-generalization}
adds variation within a one-parameter payoff family. It is not evidence
for generalisation to arbitrary payoff shapes.

When the output basis is conforming and the matrices represent that basis,
three properties follow from \eqref{eq:operator-map}, and together
they are the reason for choosing a coefficient-valued output over pointwise
values on a grid:
\begin{enumerate}[leftmargin=2em]
\item \emph{Conformity.}  $U_{\theta,h}(\tau)\in\Vh\subset\Vmu$ for every
$\theta$, so the prediction is an admissible competitor in the same
variational problem as $U_h$.
\item \emph{Insertability.}  $\Ah(\mu)\cc_\theta$ is meaningful and is
computed with the solver's own FFT routine, so the residual is a residual of
the governing equation, not of a surrogate model of it.
\item \emph{Exact differentiation.}  $\partial_xU_{\theta,h}$ and
$\partial_{xx}U_{\theta,h}$ are obtained from the B-spline derivative
recursions, not from finite differences of noisy pointwise output.  This is
what makes Corollary~\ref{cor:greeks} possible.
\end{enumerate}

\subsection{Exact embedding of the initial condition}
\label{ssec:exact-ic}

The discrete initial target is the one the trusted solver uses. It forms the
raw $L^2(\Omega)$ projection of the payoff onto all $N+p-1$ stored splines,
$\bar M\bar\cc^0=\bar\bb^0$ with $\bar b^0_k=(\psi,\phi_{h,k})_{L^2(\Omega)}$,
and then overwrites the four prescribed coefficients by their boundary values
($K\sqrt h$ on the left, $0$ on the right). We denote the result by
$\cc_h^0(\mu,\psi)$ and its free part by $\zz_h^0=E^\top\cc_h^0$. This is
\emph{not} the constrained projection
$M_I^{-1}E^\top(\bar\bb^0-\bar M b^0)$; the two differ, and every
certificate in this paper is relative to the specified target $\zz_h^0$,
because it is the initial value of every stored reference trajectory and of
every training trajectory. In the implementation the lifting is the
prescribed coefficient vector $b(\tau)$ of the boundary-crossing splines,
which plays the role of $g$ in Lemma~\ref{lem:lifting-cancels}.
We impose the initial condition architecturally rather than by penalty:
\begin{equation}
 \boxed{
 \cc_\theta(\tau;\mu,T,\psi)
 =\cc_h^0(\mu,\psi)+\tau\,\mathcal N_\theta(\tau,\mu,T,\psi),}
 \label{eq:exact-ic}
\end{equation}
so that $\cc_\theta(0;\mu,\psi)=\cc_h^0(\mu,\psi)$ and hence
\begin{equation}
  \ee^0=\cc_h^0-\cc_\theta^0=0
  \label{eq:zero-initial-error}
\end{equation}
in exact arithmetic. Every certificate in Section~\ref{sec:certificate} contains
$\beta_0^2=\norm{\zz_\theta^0-\zz_h^0}_{M_I}^2$. This term is computable
from the known target without the reference trajectory, and the integrated
workflow always computes it rather than assuming it vanishes: for the
DeepONet checkpoints it evaluates to exactly zero, whereas a finite-rank POD
model has $\beta_0>0$, which enters its acceptance. Omitting $\beta_0$ would
be unsafe: in the 48 validation cases of Table~\ref{tab:augmented}, trajectories
following the exact discrete dynamics from a perturbed initial value have
rounding-level residuals, and the residual term alone underestimates their
error in every case.

\begin{remark}[If the initial condition is not embedded exactly]
\label{rem:inexact-ic}
Should a future variant relax \eqref{eq:exact-ic}, the initial term remains
computable provided the projection $\cc_h^0$ is computed online, which costs a
single banded solve.  The certificate then reads
$\etaSharp^2+\norm{\cc_h^0-\cc_\theta^0}_{\Mh}^2$.  Dropping this term merely because it is numerically small would invalidate
the bound.
\end{remark}

The neural checkpoints use this embedding. The repaired POD baseline of
Section~\ref{ssec:pod-results} projects the target initial value onto its
subspace in the $M_I$ inner product and carries the resulting $\beta_0$ in its
certificate. Prescribed coefficients are always overwritten by the specified
lifting before certification, and the returned trajectory is the projected one.

\subsection{Architectures}
\label{ssec:architectures}

The main architecture is a DeepONet-style coefficient operator with three
encoders and a decoder,
\begin{equation}
  z_\psi=E_\psi(\psi),\quad
  z_\mu=E_\mu(\mu,T),\quad
  z_\tau=E_\tau(\tau),
  \qquad
  \mathcal N_\theta=D_\theta(z_\psi,z_\mu,z_\tau)\in\R^{N_h}.
  \label{eq:coefficient-deeponet}
\end{equation}
Maturity is shown explicitly because, at fixed $N_T$, the discrete time
step depends on $T$ even at a common physical evaluation time $\tau$.
The payoff encoder may consume B-spline coefficients of $\psi$, sensor values
$\{\psi(x_j)\}$, or coefficients in a separate payoff basis; the first is the
most natural because the projection is already computed for
\eqref{eq:exact-ic}.  Architectural novelty is explicitly not a contribution
of this paper.  The comparison set is
\begin{enumerate}[leftmargin=2em]
\item a plain MLP mapping $(\mu,T,\psi,\tau)$ to coefficients (debugging
baseline; if the sophisticated models cannot beat it, that is itself a
result);
\item the DeepONet-style coefficient operator \eqref{eq:coefficient-deeponet};
\item an FNO on a grid representation of payoff and parameters;
\item a projected FNO, $\Pi_h\circ\mathrm{FNO}$, where $\Pi_h$ is the $L^2$
projection onto $\Vh$, so that certification is applied after projection;
\item a POD/reduced-basis surrogate and a POD-plus-neural coefficient map.
\end{enumerate}

\begin{remark}[What certification demands of a baseline]
\label{rem:fno-caveat}
An unconstrained grid FNO does not produce an element of $\Vh$, so
Theorem~\ref{thm:fully-discrete-certificate} does not apply to it directly.
Two honest options exist: report the FNO as an accuracy-only baseline, or
certify $\Pi_h\circ\mathrm{FNO}$ and state clearly that the certified object
is the projection, whose error relative to the raw FNO output must then also
be reported.  The reported study takes the second route. Projection costs and the
difference between raw and projected outputs must be included in a matched
comparison; the supplied tables do not separately identify that difference.
\end{remark}

The historical experiments compare an MLP with three hidden layers of width
$128$, a coefficient DeepONet with $64$-dimensional latent space, a
$32$-mode POD--Galerkin reduced model (superseded, see
Section~\ref{ssec:pod-results}) and a projected FNO with width $32$, $16$
Fourier modes and four spectral layers, all on the same discrete space. The
v5 experiments use the DeepONet and the repaired mass-weighted POD model at
ranks $8$--$62$.
Their reported parameter counts and errors are listed in
Section~\ref{sec:numerics}. A reproducible comparison also needs exact
input encodings, activations, optimiser, training schedule, normalisation
and checkpoint-selection rules.

\subsection{Galerkin-consistent losses}
\label{ssec:losses}

Coefficient-space Euclidean MSE has no variational meaning, since the basis is
not orthonormal in any relevant inner product.  We use instead the
mass-weighted data loss
\begin{equation}
 \mathcal L_{\rm data}
 =\frac{(\cc_\theta-\cc_h)^\top\Mh(\cc_\theta-\cc_h)}
        {\cc_h^\top\Mh\cc_h+\epsilon},
 \label{eq:data-loss}
\end{equation}
which is the relative squared $L^2(\Omega)$ error of the homogeneous
spline field, with regularisation $\epsilon$,
and, where an energy-weighted variant is wanted,
\begin{equation}
 \mathcal L_{\rm energy}
 =\frac{(\cc_\theta-\cc_h)^\top\Sh(\mu)(\cc_\theta-\cc_h)}
        {\cc_h^\top\Sh(\mu)\cc_h+\epsilon}.
 \label{eq:energy-loss}
\end{equation}
The residual loss uses the \emph{same} dual norm as the certificate,
\begin{equation}
 \mathcal L_{\rm res}
 =\frac{1}{N_T}\sum_{n=0}^{N_T-1}
 (\rr_\theta^{n+1/2})^\top\Sh(\mu)^{-1}\rr_\theta^{n+1/2},
 \label{eq:residual-loss}
\end{equation}
with $\rr_\theta^{n+1/2}$ given by \eqref{eq:cn-residual}, and the hybrid
objective is
\begin{equation}
 \mathcal L=\mathcal L_{\rm data}+\lambda_R\mathcal L_{\rm res}.
 \label{eq:hybrid-loss}
\end{equation}
The conceptual difference from a pointwise PINN residual is that
\eqref{eq:residual-loss} is not a heuristic penalty: by
Theorem~\ref{thm:fully-discrete-certificate} the corresponding time-weighted accumulation bounds the absolute
maximum-in-time mass error, with its initial term. The efficiency estimate
uses a different, stronger graph norm; neither result makes the unnormalised
residual loss equivalent to the normalised data loss. Because
$\etaSharp^2=\dt N_T\mathcal L_{\rm res}=T\mathcal L_{\rm res}$ for a
uniform grid, comparing residual weights across maturities also requires
consistent scaling.

Residual regularisation is tested through a dedicated
$\lambda_R$ sweep in Section~\ref{ssec:learning-ablation}.  The observed
benefit is confined to a very small weight ($\lambda_R\approx10^{-9}$); larger
weights degrade both predictive error and the certificate, so residual
regularisation is not used as an unqualified accuracy claim.

% =================================================================
\section{Galerkin residual certification}
\label{sec:certificate}
% =================================================================

This section is the mathematical centrepiece.  The organising idea is simple:
once the neural output is conforming, certifying it is not a machine-learning
question at all, but a stability estimate for the error equation between two
trajectories in the \emph{same} Galerkin space.  Everything that follows is an
elaboration of that observation.

\subsection{The discrete energy Riesz map}
\label{ssec:riesz}

For real coefficient vectors define the symmetric part of the stiffness
matrix,
\begin{equation}
  \Sh(\mu):=\tfrac12\bigl(\Ah(\mu)+\Ah(\mu)^\top\bigr).
  \label{eq:Sh}
\end{equation}

\begin{lemma}[Discrete Riesz map under exact Galerkin assembly]
\label{lem:Sh-spd}
Let $\mu\in\Pset$ and suppose the assembled stiffness matrix represents the
bilinear form exactly, $(\Ah)_{kk'}=a_\mu(\phi_{h,k'},\phi_{h,k})$.  Then for
every $\zz\in\R^{N_h}\setminus\{0\}$, with
$z_h=\sum_kz_k\phi_{h,k}$,
\begin{equation}
  \zz^\top\Sh(\mu)\zz=\zz^\top\Ah(\mu)\zz
  =a_\mu(z_h,z_h)\ge c_1\norm{z_h}_{\Vmu}^2>0,
  \label{eq:Sh-spd}
\end{equation}
so $\Sh(\mu)$ is symmetric positive definite and induces the norms
\begin{equation}
 \norm{\zz}_{\Sh}^2=\zz^\top\Sh\zz,
 \qquad
 \norm{\bb}_{\Sh^{-1}}^2=\bb^\top\Sh^{-1}\bb
 =\sup_{\zz\ne0}\frac{(\bb^\top\zz)^2}{\zz^\top\Sh\zz}.
 \label{eq:discrete-energy-dual}
\end{equation}
Moreover $B_+:=\Mh+\tfrac{\dt}{2}\Ah(\mu)$ and
$\Mh+\dt\Ah(\mu)$ are nonsingular for every $\dt>0$.
\end{lemma}

\begin{proof}
The equality $\zz^\top\Sh\zz=\zz^\top\Ah\zz$ follows because the quadratic
form is scalar, while exact Galerkin assembly gives
$\zz^\top\Ah\zz=a_\mu(z_h,z_h)$.  The inequality is
Theorem~\ref{thm:uniform-coercivity}.  The symmetric parts of $B_+$ and
$\Mh+\dt\Ah$ are respectively $\Mh+\tfrac{\dt}{2}\Sh$ and
$\Mh+\dt\Sh$, hence positive definite; both matrices are therefore
nonsingular.
\end{proof}

\begin{remark}[Scope of the implemented positivity gate]
\label{rem:spd-gate}
The exact conforming result concerns the free trial space, and the integrated
workflow gates the free block $S_I$. The historical pipeline gated the raw
$66$-dimensional $\bar S$, including four prescribed directions; its failures
do not establish loss of free-space coercivity
(Proposition~\ref{prop:legacy-conservative}). When the free gate fails, no
certificate is issued and the status \textsc{uncertifiable} is returned; on the
populations of Section~\ref{sec:numerics} the free gate never failed.
\end{remark}

The choice \eqref{eq:Sh} instantiates the abstract Riesz matrix with an object
that the solver already possesses: no auxiliary fractional Gram matrix is required for the algebraic bound.
Materialising the matrix and establishing its spectral lower bound can
nevertheless introduce substantial per-query setup work.  This is a practical
point of some importance, since assembling an independent $H^{Y/2}$ Gram
matrix for every $\mu$ would cost as much as the stiffness matrix itself.

\begin{lemma}[Positive-rate mass lower bound for exact assembly]
\label{lem:mass-lower}
Under the exact-assembly hypothesis of Lemma~\ref{lem:Sh-spd}, for
$\mu\in\Pset$ (so $\rf>0$),
\begin{equation}
  \Sh(\mu)\succeq \rf\,\Mh
  \label{eq:S-ge-rM}
\end{equation}
in the Loewner order, and consequently
\begin{equation}
  \bb^\top\Sh(\mu)^{-1}\bb\le \rf^{-1}\,\bb^\top\Mh^{-1}\bb
  \qquad\text{for every }\bb\in\R^{N_h}.
  \label{eq:mass-upper}
\end{equation}
\end{lemma}

\begin{proof}
By Lemma~\ref{lem:conformity} and Plancherel's theorem,
\[
  \zz^\top\Sh\zz=a_\mu(z_h,z_h)=a_\R(\widetilde z_h,\widetilde z_h)
  =\frac1{2\pi}\int_\R\Real A_\mu(\xi)\,\bigl|\widehat{\widetilde z_h}(\xi)\bigr|^2\,d\xi .
\]
The exact symbol identity of \cite{DamircheliCGMY2026}, namely
\begin{equation}
\Real A_\mu(\xi)
=\rf+C\Gamma(-Y)\Bigl[\bigl(M^Y-\Real(M-i\xi)^Y\bigr)
+\bigl(G^Y-\Real(G+i\xi)^Y\bigr)\Bigr],
\label{eq:symbol-real-exact}
\end{equation}
together with $\Gamma(-Y)>0$ for $Y\in(1,2)$ and the monotonicity
$\Real(a\pm i\xi)^Y\le a^Y$, gives $\Real A_\mu(\xi)\ge \rf$ for all $\xi$.  Hence
\[
  \zz^\top\Sh\zz\ \ge\ \rf\norm{\widetilde z_h}^2_{L^2(\R)}
  =\rf\norm{z_h}^2_{L^2(\Omega)}=\rf\,\zz^\top\Mh\zz .
\]
Inversion of the Loewner
inequality between positive definite matrices reverses it, giving
\eqref{eq:mass-upper}.
\end{proof}

\begin{remark}
Bound \eqref{eq:mass-upper} is attractive because $\Mh$ is banded, parameter
independent, and factorised once for the whole study.  It is, however,
proportional to $\rf^{-1}$ and therefore degrades as $\rf\downarrow0$.  This
is the precise reason for reporting both the sharp $\Sh^{-1}$ certificate and
the cheap theoretical mass bound.  The implemented online rule instead uses the guaranteed bracket of Section~\ref{ssec:two-stage}.
\end{remark}

The last structural quantity needed is the discrete continuity constant of the
stiffness matrix measured between the $\Sh$ and $\Sh^{-1}$ norms.

\begin{lemma}[Discrete operator norm of the stiffness]
\label{lem:kappa}
Whenever the assembled matrix $\Sh(\mu)$ is symmetric positive definite,
define
\begin{equation}
  \kappa_h(\mu):=\sup_{\vv\ne0}
  \frac{\norm{\Ah(\mu)\vv}_{\Sh(\mu)^{-1}}}{\norm{\vv}_{\Sh(\mu)}}
  =\bigl\lVert \Sh^{-1/2}\Ah\,\Sh^{-1/2}\bigr\rVert_2 .
  \label{eq:kappa-def}
\end{equation}
Then $\kappa_h(\mu)\ge1$ for the actually assembled matrices.  If, in addition,
the exact-assembly hypothesis of Lemma~\ref{lem:Sh-spd} holds, then
\begin{equation}
  1\le\kappa_h(\mu)\le\frac{C_{\rm cont}}{c_1}
  \qquad\text{for every }\mu\in\Pset,
  \label{eq:kappa-bound}
\end{equation}
uniformly in $h$.  Under the same exact-assembly hypothesis the upper bound
also holds with $\Ah$ replaced by $\Ah^\top$.
\end{lemma}

\begin{proof}
For any assembled $\Ah$ with SPD symmetric part $\Sh$ and any nonzero $\vv$,
duality in the $\Sh$ pairing gives
\[
 \norm{\Ah\vv}_{\Sh^{-1}}
 \ge \frac{\vv^\top\Ah\vv}{\norm{\vv}_{\Sh}}
 =\frac{\vv^\top\Sh\vv}{\norm{\vv}_{\Sh}}
 =\norm{\vv}_{\Sh},
\]
because the skew-symmetric part of $\Ah$ vanishes in the quadratic form.
Hence $\kappa_h\ge1$ without invoking exact assembly.

Under exact Galerkin assembly,
$\ww^\top\Ah\vv=a_\mu(v_h,w_h)$.  Continuity and
$\norm{z_h}_{\Vmu}\le c_1^{-1/2}\norm{\zz}_{\Sh}$ from
\eqref{eq:Sh-spd} yield
\[
  \norm{\Ah\vv}_{\Sh^{-1}}
  =\sup_{\ww\ne0}\frac{a_\mu(v_h,w_h)}{\norm{\ww}_{\Sh}}
  \le\frac{C_{\rm cont}}{c_1}\norm{\vv}_{\Sh},
\]
which proves the upper bound.  Interchanging the trial and test arguments
gives the corresponding estimate for $\Ah^\top$.
\end{proof}

\begin{remark}[$\kappa_h$ is directly computable]
\label{rem:kappa-computable}
The matrix-level constant in \eqref{eq:kappa-def} is available whether or not
the numerical assembly is identified with the exact variational matrix:
$\kappa_h(\mu)^2$ is the largest generalised eigenvalue of
\[
  \Ah^\top\Sh^{-1}\Ah\vv=\kappa_h^2\Sh\vv .
\]
At the reported sizes used here it is evaluated by a stable dense
factorisation/eigenvalue calculation; a matrix-free extremal-eigenvalue route
is also available at larger dimension.  It is important not to confuse
$\kappa_h$ with the spectral condition number of $\Sh$.  The numerical
efficiency test in Section~\ref{ssec:certificate-results} uses
\eqref{eq:kappa-def} exactly.
\end{remark}

\subsection{Semidiscrete certificate}
\label{ssec:semidiscrete}

Let $\cc_h(\tau)$ solve \eqref{eq:semidiscrete-system} and let
$\cc_\theta(\tau)$ be any differentiable conforming prediction.  Define the
semidiscrete residual and coefficient error
\begin{equation}
  \rr_\theta(\tau)=\Mh\dot\cc_\theta+\Ah\cc_\theta-\ff,
  \qquad
  \ee=\cc_h-\cc_\theta,
  \qquad\text{so that}\qquad
  \Mh\dot\ee+\Ah\ee=-\rr_\theta.
  \label{eq:semidiscrete-error-equation}
\end{equation}

\begin{theorem}[Sharp semidiscrete certificate]
\label{thm:semidiscrete-certificate}
Let $\mu\in\Pset$ and assume that the assembled $\Sh(\mu)$ is symmetric
positive definite.  For every $t\in[0,T]$,
\begin{equation}
 \norm{\ee(t)}_{\Mh}^2
 +\int_0^t\norm{\ee(\varsigma)}_{\Sh}^2\,d\varsigma
 \le
 \norm{\ee(0)}_{\Mh}^2
 +\int_0^t\norm{\rr_\theta(\varsigma)}_{\Sh^{-1}}^2\,d\varsigma .
 \label{eq:semidiscrete-certificate}
\end{equation}
If in addition $\ee(0)=0$, then
$\norm{U_h(t)-U_{\theta,h}(t)}_{L^2(\Omega)}
 \le\bigl(\int_0^t\rr_\theta^\top\Sh^{-1}\rr_\theta\bigr)^{1/2}$.
\end{theorem}

\begin{proof}
Multiply \eqref{eq:semidiscrete-error-equation} by $\ee^\top$.  Since $\Mh$ is
symmetric and time independent and
$\ee^\top\Ah\ee=\norm{\ee}_{\Sh}^2$,
\[
 \tfrac12\tfrac{d}{dt}\norm{\ee}_{\Mh}^2+\norm{\ee}_{\Sh}^2=-\ee^\top\rr_\theta .
\]
Cauchy--Schwarz and Young's inequality give
\[
 |\ee^\top\rr_\theta|\le\norm{\ee}_{\Sh}\norm{\rr_\theta}_{\Sh^{-1}}
 \le\tfrac12\norm{\ee}_{\Sh}^2+\tfrac12\norm{\rr_\theta}_{\Sh^{-1}}^2.
\]
Multiplying by two and integrating gives
\eqref{eq:semidiscrete-certificate}.  The identity
$\norm{\ee}_{\Mh}=\norm{U_h-U_{\theta,h}}_{L^2(\Omega)}$ is the definition of
$\Mh$.
\end{proof}

The semidiscrete estimate is the cleanest statement of the mechanism, but it
is not the one we implement: it requires $\dot\cc_\theta$, hence either
automatic differentiation in a continuous time variable or a time-continuous
architecture, and it certifies a semidiscrete reference that the solver never
computes.  We therefore turn to the fully discrete case.

\subsection{Fully discrete Crank--Nicolson certificate}
\label{ssec:cn-certificate}

For a neural trajectory $\{\cc_\theta^n\}_{n=0}^{N_T}$ define the
Crank--Nicolson Galerkin residual
\begin{equation}
 \boxed{
 \rr_\theta^{n+1/2}
 :=\Mh\frac{\cc_\theta^{n+1}-\cc_\theta^n}{\dt}
  +\Ah(\mu)\overline\cc_\theta^{\,n+1/2}
  -\overline\ff^{\,n+1/2},}
 \label{eq:cn-residual}
\end{equation}
and the coefficient error $\ee^n=\cc_h^n-\cc_\theta^n$.  Subtracting
\eqref{eq:cn-residual} from \eqref{eq:cn-reference},
\begin{equation}
 \Mh\frac{\ee^{n+1}-\ee^n}{\dt}+\Ah(\mu)\overline\ee^{\,n+1/2}
 =-\rr_\theta^{n+1/2}.
 \label{eq:discrete-error-equation}
\end{equation}

\begin{theorem}[Fully discrete sharp certificate]
\label{thm:fully-discrete-certificate}
Let $\mu\in\Pset$ and assume that the assembled matrix $\Sh(\mu)$ is
symmetric positive definite.  For every $m=1,\dots,N_T$,
\begin{equation}
 \boxed{
 \norm{\ee^m}_{\Mh}^2
 +\dt\sum_{n=0}^{m-1}\norm{\overline\ee^{\,n+1/2}}_{\Sh}^2
 \le
 \norm{\ee^0}_{\Mh}^2
 +\dt\sum_{n=0}^{m-1}\norm{\rr_\theta^{n+1/2}}_{\Sh^{-1}}^2 .}
 \label{eq:fully-discrete-certificate}
\end{equation}
Define the \emph{sharp certificate}
\begin{equation}
 \boxed{
 \etaSharp^2(m;\mu,\psi)
 :=\dt\sum_{n=0}^{m-1}
 (\rr_\theta^{n+1/2})^\top\Sh(\mu)^{-1}\rr_\theta^{n+1/2}.}
 \label{eq:sharp-estimator}
\end{equation}
If the payoff is embedded exactly, so that \eqref{eq:zero-initial-error}
holds, then
\begin{equation}
  \max_{1\le m\le N_T}\ \Bigl[
  \norm{U_h^m-U_{\theta,h}^m}_{L^2(\Omega)}
  \Bigr]
  \le\etaSharp(N_T;\mu,\psi),
  \label{eq:sharp-L2}
\end{equation}
and, by Lemma~\ref{lem:lifting-cancels}, the same bound holds for the price,
$\norm{V_h^m-V_{\theta,h}^m}_{L^2(\Omega)}\le\etaSharp(N_T;\mu,\psi)$.
\end{theorem}

\begin{proof}
Multiply \eqref{eq:discrete-error-equation} on the left by
$(\overline\ee^{\,n+1/2})^\top$.  Because $\Mh$ is symmetric,
\[
 (\overline\ee^{\,n+1/2})^\top\Mh\frac{\ee^{n+1}-\ee^n}{\dt}
 =\frac{1}{2\dt}\bigl(\norm{\ee^{n+1}}_{\Mh}^2-\norm{\ee^n}_{\Mh}^2\bigr),
\]
and, by the definition \eqref{eq:Sh},
\[
 (\overline\ee^{\,n+1/2})^\top\Ah\overline\ee^{\,n+1/2}
 =\norm{\overline\ee^{\,n+1/2}}_{\Sh}^2 .
\]
The residual term obeys
\[
 |(\overline\ee^{\,n+1/2})^\top\rr_\theta^{n+1/2}|
 \le\tfrac12\norm{\overline\ee^{\,n+1/2}}_{\Sh}^2
 +\tfrac12\norm{\rr_\theta^{n+1/2}}_{\Sh^{-1}}^2 .
\]
Multiplying by $2\dt$ and summing over $n=0,\dots,m-1$ telescopes the mass
terms and gives \eqref{eq:fully-discrete-certificate}.  Since the right-hand
side of \eqref{eq:fully-discrete-certificate} is nondecreasing in $m$, the
maximum over $m$ in \eqref{eq:sharp-L2} is bounded by its value at $m=N_T$.
\end{proof}

\begin{corollary}[Cheap mass certificate under exact assembly]
\label{cor:mass-certificate}
Under the hypotheses of Theorem~\ref{thm:fully-discrete-certificate},
$\ee^0=0$, and the exact-assembly hypothesis of
Lemma~\ref{lem:mass-lower}, define
\begin{equation}
 \boxed{
 \etaMass^2(m;\mu,\psi)
 :=\frac{\dt}{\rf}\sum_{n=0}^{m-1}
 (\rr_\theta^{n+1/2})^\top\Mh^{-1}\rr_\theta^{n+1/2}.}
 \label{eq:mass-estimator}
\end{equation}
Then
\begin{equation}
 \norm{U_h^m-U_{\theta,h}^m}_{L^2(\Omega)}
 \le\etaSharp(N_T)\le\etaMass(N_T).
 \label{eq:sharp-cheap-chain}
\end{equation}
\end{corollary}

\begin{proof}
Apply \eqref{eq:mass-upper} termwise in \eqref{eq:sharp-estimator}.
\end{proof}

The next lemma records that the estimator does not merely bound the error but
vanishes exactly when it should.  It is the property that rules out a
certificate which is simultaneously reliable and uninformative.

\begin{lemma}[Consistency of the estimator]
\label{lem:consistency}
Assume $\Sh\succ0$ and \eqref{eq:zero-initial-error}. Then
$\etaSharp(N_T;\mu,\psi)=0$ if and only if $\cc_\theta^n=\cc_h^n$ for every
$n=0,\dots,N_T$.
\end{lemma}

\begin{proof}
If $\cc_\theta^n=\cc_h^n$ for all $n$ then every residual
\eqref{eq:cn-residual} vanishes by \eqref{eq:cn-reference}, so
$\etaSharp=0$.  Conversely, if $\etaSharp=0$ then, since $\Sh^{-1}$ is
positive definite, $\rr_\theta^{n+1/2}=0$ for every $n$, and
\eqref{eq:discrete-error-equation} becomes $B_+\ee^{n+1}=B_-\ee^n$ with
$\ee^0=0$.  The symmetric part of $B_+$ is
$\Mh+\tfrac{\dt}{2}\Sh\succ0$, so $B_+$ is nonsingular and
$\ee^{n}=0$ for all $n$ by induction.
\end{proof}

\begin{remark}[The discrete target and evaluation discrepancies]
\label{rem:scope-of-certificate}
The bound refers to the exact algebraic trajectory on the selected mesh and
time scheme with the same matrices, forcing and initial vector. It does not
certify continuum discretisation error or automatically absorb a different
quadrature rule. If the reference uses $(A_h,f_{\rm ref})$ and the residual
uses $(A_{\rm ev},f_{\rm ev})$, the residual required in the error equation is
\[
 \rho_{\rm ref}^{n+1/2}
 =\rho_{\rm ev}^{n+1/2}
  +(A_h-A_{\rm ev})\overline c_\theta^{\,n+1/2}
  +(f_{\rm ev}^{n+1/2}-f_{\rm ref}^{n+1/2}).
\]
The discrepancy can cancel part of the evaluated residual; it cannot be
assumed merely to inflate it. Use identical assembly, or bound the discrepancy
in the same dual norm and add its bound before time accumulation.
A computed reference trajectory also has a solver residual. Comparisons with
that trajectory must either include its error allowance or demonstrate that
it is negligible relative to the tested certificate margins.
\end{remark}

\subsection{Backward-Euler and Rannacher-smoothed reference trajectories}
\label{ssec:be-certificate}

If the reference solver damps the payoff kink by backward-Euler start-up, the
certificate must follow suit, for the reason given in
Remark~\ref{rem:certify-what-you-ran}.  For a step of size $\dt_n$ define
\begin{equation}
 \rr_\theta^{n+1}
 :=\Mh\frac{\cc_\theta^{n+1}-\cc_\theta^n}{\dt_n}
  +\Ah(\mu)\cc_\theta^{n+1}-\ff^{n+1}.
 \label{eq:be-residual}
\end{equation}

\begin{theorem}[Backward-Euler certificate]
\label{thm:be-certificate}
Let $\mu\in\Pset$, assume the assembled $\Sh(\mu)$ is symmetric positive
definite, and let $\cc_h^{n+1}$ solve \eqref{eq:be-reference}.  Then
\begin{equation}
 \norm{\ee^{n+1}}_{\Mh}^2
 +\norm{\ee^{n+1}-\ee^{n}}_{\Mh}^2
 +\dt_n\norm{\ee^{n+1}}_{\Sh}^2
 \le
 \norm{\ee^{n}}_{\Mh}^2
 +\dt_n\norm{\rr_\theta^{n+1}}_{\Sh^{-1}}^2 .
 \label{eq:be-certificate}
\end{equation}
\end{theorem}

\begin{proof}
Subtracting \eqref{eq:be-residual} from \eqref{eq:be-reference} gives
$\Mh(\ee^{n+1}-\ee^{n})/\dt_n+\Ah\ee^{n+1}=-\rr_\theta^{n+1}$.  Multiply on
the left by $(\ee^{n+1})^\top$ and use the elementary identity
$2\,\yy^\top\Mh(\yy-\zz)=\norm{\yy}_{\Mh}^2-\norm{\zz}_{\Mh}^2
+\norm{\yy-\zz}_{\Mh}^2$, valid for symmetric $\Mh$ and any $\yy,\zz$,
together with
$(\ee^{n+1})^\top\Ah\ee^{n+1}=\norm{\ee^{n+1}}_{\Sh}^2$ and
$|(\ee^{n+1})^\top\rr_\theta^{n+1}|
\le\tfrac12\norm{\ee^{n+1}}_{\Sh}^2+\tfrac12\norm{\rr_\theta^{n+1}}_{\Sh^{-1}}^2$.
Multiplying by $\dt_n$ and rearranging gives \eqref{eq:be-certificate}.
\end{proof}

\begin{corollary}[Rannacher-smoothed certificate]
\label{cor:rannacher}
Assume $\Sh\succ0$. Let the reference use $n_{\rm BE}$ backward-Euler steps of sizes
$\dt_0,\dots,\dt_{n_{\rm BE}-1}$ followed by Crank--Nicolson steps of size
$\dt$, and suppose \eqref{eq:zero-initial-error} holds.  Define
\begin{equation}
 \etaSharp^2(m)
 :=\sum_{n=0}^{n_{\rm BE}-1}\dt_n\norm{\rr_\theta^{n+1}}_{\Sh^{-1}}^2
  +\dt\!\!\sum_{n=n_{\rm BE}}^{m-1}\!\!
   \norm{\rr_\theta^{n+1/2}}_{\Sh^{-1}}^2 ,
 \label{eq:rannacher-estimator}
\end{equation}
Use \eqref{eq:be-residual} on the start-up steps and
\eqref{eq:cn-residual} thereafter. For every $m$,
\[
 \norm{U_h^m-U_{\theta,h}^m}_{L^2(\Omega)}\le\etaSharp(N_T).
\]
\end{corollary}

\begin{proof}
Discard the nonnegative terms $\norm{\ee^{n+1}-\ee^n}_{\Mh}^2$ and
$\dt_n\norm{\ee^{n+1}}_{\Sh}^2$ in \eqref{eq:be-certificate} and sum over the
start-up steps; this bounds $\norm{\ee^{n_{\rm BE}}}_{\Mh}^2$ by the first sum
in \eqref{eq:rannacher-estimator}, since $\ee^0=0$.  Then apply
Theorem~\ref{thm:fully-discrete-certificate} on the remaining steps with
initial error $\ee^{n_{\rm BE}}$.
\end{proof}

Corollary~\ref{cor:rannacher} covers the practically important case of
Rannacher start-up for the kinked European-put payoff; applying a pure
Crank--Nicolson residual to those start-up steps would be inconsistent with
the reference trajectory.

\subsection{A shifted certificate for nonpositive rates}
\label{ssec:shifted}

The positive-rate restriction $\rf\ge r_{\min}>0$ in \eqref{eq:parameter-set}
buys the clean constant-free certificate, but it excludes maturity-specific
calibration rates that can be nonpositive.  The exponential shift already
present in the deterministic analysis removes the restriction at the cost of a
factor $e^{\sigma T}$.

Let $c_2\ge0$ be the G\aa rding constant of \cite{DamircheliCGMY2026}, so that
$a_\mu(v,v)\ge c_1\norm{v}_{\Vmu}^2-c_2\norm{v}_{L^2(\Omega)}^2$, and let
$\sigma>c_2$.  Put
\begin{equation}
  A_{h,\sigma}:=\Ah+\sigma\Mh,
  \qquad
  S_{h,\sigma}:=\tfrac12\bigl(A_{h,\sigma}+A_{h,\sigma}^\top\bigr)=\Sh+\sigma\Mh .
  \label{eq:shifted-matrices}
\end{equation}
Define a shifted reference by advancing $\mathbf d_h^n$ and reconstructing
$\cc_{h,\sigma}^n=e^{\sigma\tau_n}\mathbf d_h^n$ after the solve.
Write $U_{h,\sigma}^n=\sum_k c_{h,\sigma,k}^n\phi_{h,k}$.
The shifted scheme is
\begin{equation}
  \Mh\frac{\mathbf d_h^{n+1}-\mathbf d_h^{n}}{\dt}
  +A_{h,\sigma}\overline{\mathbf d}_h^{\,n+1/2}
  =\overline{\mathbf g}^{\,n+1/2},
  \qquad
  \mathbf g^n:=e^{-\sigma\tau_n}\ff^n .
  \label{eq:shifted-cn}
\end{equation}

\begin{theorem}[Shifted certificate]
\label{thm:shifted-certificate}
Assume exact Galerkin assembly, let $\sigma>c_2$, and set $\mathbf d_\theta^n:=e^{-\sigma\tau_n}\cc_\theta^n$, and
define the shifted residual
\begin{equation}
  \rr_{\theta,\sigma}^{n+1/2}
  :=\Mh\frac{\mathbf d_\theta^{n+1}-\mathbf d_\theta^{n}}{\dt}
   +A_{h,\sigma}\overline{\mathbf d}_\theta^{\,n+1/2}
   -\overline{\mathbf g}^{\,n+1/2}.
  \label{eq:shifted-residual}
\end{equation}
Then $S_{h,\sigma}$ is symmetric positive definite with
$S_{h,\sigma}\succeq(\sigma-c_2)\Mh$, and if
$\cc_\theta^0=\cc_{h,\sigma}^0$ then, for every $m$,
\begin{equation}
  \norm{U_{h,\sigma}^m-U_{\theta,h}^m}_{L^2(\Omega)}
  \le e^{\sigma T}\,
  \Bigl(
  \dt\sum_{n=0}^{N_T-1}
  (\rr_{\theta,\sigma}^{n+1/2})^\top S_{h,\sigma}^{-1}\rr_{\theta,\sigma}^{n+1/2}
  \Bigr)^{1/2}
  =:e^{\sigma T}\,\eta_{\sigma}(N_T;\mu,\psi).
  \label{eq:shifted-certificate}
\end{equation}
The mass bound \eqref{eq:mass-upper} holds with $\rf$ replaced by
$\sigma-c_2$.
\end{theorem}

\begin{proof}
By the G\aa rding inequality and Lemma~\ref{lem:conformity},
$\zz^\top S_{h,\sigma}\zz=a_\mu(z_h,z_h)+\sigma\norm{z_h}^2_{L^2(\Omega)}
\ge c_1\norm{z_h}_{\Vmu}^2+(\sigma-c_2)\norm{z_h}^2_{L^2(\Omega)}$, which is
positive for $\zz\ne0$ and dominates $(\sigma-c_2)\zz^\top\Mh\zz$.  Applying
Theorem~\ref{thm:fully-discrete-certificate} verbatim to the pair
\eqref{eq:shifted-cn}, \eqref{eq:shifted-residual} with $\Sh$ replaced by
$S_{h,\sigma}$ bounds
$\norm{\mathbf d_h^m-\mathbf d_\theta^m}_{\Mh}$ by $\eta_\sigma$, the initial
term vanishing because $\mathbf d^0=\cc^0$.  Undoing the shift,
$\cc_{h,\sigma}^m-\cc_\theta^m=e^{\sigma\tau_m}(\mathbf d_h^m-\mathbf d_\theta^m)$ and
$\tau_m\le T$.
\end{proof}

\begin{remark}
The shifted discrete target is not the original unshifted
Crank--Nicolson trajectory: exponential transformation and CN discretisation
do not commute exactly. The theorem certifies the reconstructed shifted
trajectory only. It is stated for completeness and for the
calibration regime; it is not needed for the main experiments, which stay
inside the positive-rate box.  The factor $e^{\sigma T}$ is the price of abandoning coercivity, so we do
not present the shifted estimator as equally sharp.
\end{remark}

\subsection{A guaranteed certificate under inexact Riesz solves}
\label{ssec:guaranteed}

Theorem~\ref{thm:fully-discrete-certificate} contains the exact quantity
$\rr^\top\Sh^{-1}\rr$.  In practice the Riesz problem is solved iteratively,
so the online method needs an upper bound that remains valid before the
iteration has converged.  The implementation first computes a
query-wise lower-bound estimate $\alpha_{\rm lb}(\mu)$ intended to satisfy
\begin{equation}
  0<\alpha_{\rm lb}(\mu)\le \lambda_{\min}(\Sh(\mu)).
  \label{eq:alpha-lb}
\end{equation}
The implemented spectral estimate is
$\widehat\lambda_{\min}-\|\widehat V^\top S\widehat V-
\operatorname{diag}(\widehat\lambda)\|_F$.
For an exactly orthogonal $V$ and exact arithmetic this follows from Weyl's
inequality. A computed eigenvector matrix is only approximately orthogonal,
and the residual and its norm are themselves rounded. The submitted code
does not enclose those defects. Hence \eqref{eq:alpha-lb} remains an
explicit mathematical hypothesis; the double-precision estimate is not a
validated numerical proof of that hypothesis. A complete enclosure must
account for the orthogonality defect and arithmetic errors, or use another
verified lower-bound construction.

For exact conforming assembly of the free space one may use
$\alpha_{\rm lb}=r\lambda_{\min}(M_I)$, with a verified mass lower bound.
This route is not available for the production matrices, which are not
exact Galerkin assemblies (Remark~\ref{rem:assembly-convention}); the
reported results use the spectral estimate above. Against a 40-digit
eigenvalue of the same assembled $S_I$, the estimate satisfied
$\alpha_{\rm lb}\le\lambda_{\min}(S_I)$ on all 24 validated queries
(Table~\ref{tab:riesz}); this is finite-sample evidence, not an enclosure.
For an assembled perturbation $\widetilde S_h=S_h^{\rm exact}+E$ with a
proved $\|E\|_2\le\delta$, the corresponding lower bound is
$r\lambda_{\min}(M_h)-\delta$, provided it is positive. Alternatively a
validated factorisation or spectral enclosure may establish a query-wise
bound. These options differ in cost and sharpness and must be compared.

\begin{remark}[Algebraic guarantee versus validated floating-point output]
\label{rem:arithmetic-scope}
The theorems and brackets below are exact-arithmetic statements for the
selected numerical matrices. A validated implementation additionally needs
outward bounds for residual evaluation, matrix actions, bracket arithmetic,
initial projection and final accumulation. If a computed residual
$\widehat\rho_n$ differs from the required residual by at most $d_n$ in
$S_h^{-1}$ norm, and $u_n$ is a validated upper bound for
$\widehat\rho_n^\top S_h^{-1}\widehat\rho_n$, then
\[
 \widehat\eta^2=\beta_0^2+\sum_n\Delta\tau_n
                (\sqrt{u_n}+d_n)^2
\]
is valid, where $\beta_0$ bounds the initial discrepancy and accumulation
is rounded outward. The reported experiments demonstrate observed algebraic
non-violation in double precision; the supplied files do not establish all
of these validated-arithmetic conditions.
\end{remark}

\begin{lemma}[Computable two-sided bracket for the dual norm]
\label{lem:guaranteed}
Let $\Sh\succ0$, let $\rr\in\R^{N_h}$, and let
$0<\alpha_{\rm lb}\le\lambda_{\min}(\Sh)$.  For an arbitrary
$\vv\in\R^{N_h}$ set $\sres:=\rr-\Sh\vv$ and define
\begin{equation}
  \mathfrak l(\vv):=2\vv^\top\rr-\vv^\top\Sh\vv,
  \qquad
  \mathfrak u(\vv):=\mathfrak l(\vv)
    +\frac{\norm{\sres}_2^2}{\alpha_{\rm lb}}.
  \label{eq:riesz-bracket-def}
\end{equation}
Then, for every $\vv$,
\begin{equation}
 \boxed{
   \mathfrak l(\vv)
   \le \rr^\top\Sh^{-1}\rr
   \le \mathfrak u(\vv).}
 \label{eq:riesz-bracket}
\end{equation}
The bracket width is
\begin{equation}
  \mathfrak u(\vv)-\mathfrak l(\vv)
  =\frac{\norm{\sres}_2^2}{\alpha_{\rm lb}},
  \label{eq:bracket-width}
\end{equation}
so it closes as the Riesz residual tends to zero. The upper bracket need
not decrease monotonically with iteration count, because the Euclidean
residual norm need not do so.
\end{lemma}

\begin{proof}
With $\zz=\Sh^{-1}\rr$,
\[
 \rr^\top\Sh^{-1}\rr-\mathfrak l(\vv)
 = (\vv-\zz)^\top\Sh(\vv-\zz)
 = \sres^\top\Sh^{-1}\sres\ge0.
\]
Since $\Sh\succeq\alpha_{\rm lb}I$,
$\sres^\top\Sh^{-1}\sres\le
\alpha_{\rm lb}^{-1}\norm{\sres}_2^2$, which gives the upper bound and
\eqref{eq:bracket-width}.
\end{proof}

\begin{definition}[Guaranteed certificate]
\label{def:guaranteed}
Let $\vv^{n+1/2}$ be the approximate Riesz vector returned by PCG for
$\Sh\vv=\rr_\theta^{n+1/2}$.  Define
\begin{equation}
 \boxed{
 \etaGuar^2(m;\mu,\psi)
 :=\dt\sum_{n=0}^{m-1}\mathfrak u\bigl(\vv^{n+1/2}\bigr).}
 \label{eq:guaranteed-estimator}
\end{equation}
\end{definition}

\begin{corollary}
\label{cor:guaranteed}
Under the hypotheses of Theorem~\ref{thm:fully-discrete-certificate},
\eqref{eq:zero-initial-error}, and \eqref{eq:alpha-lb},
$\norm{U_h^m-U_{\theta,h}^m}_{L^2(\Omega)}\le\etaGuar(N_T;\mu,\psi)$ for every
$m$, irrespective of the PCG stopping error.  Moreover
$\etaGuar\to\etaSharp$ as the Riesz residuals tend to zero.
\end{corollary}

\begin{proof}
Apply Lemma~\ref{lem:guaranteed} at every time level, sum the upper bounds,
and use Theorem~\ref{thm:fully-discrete-certificate}.
\end{proof}

Thus PCG effort controls sharpness but not soundness.  With the zero iterate,
the upper bracket is the spectral bound
$\norm{\rr}_2^2/\alpha_{\rm lb}$; with a converged iterate it is the sharp
Riesz value.  The analytic mass bound of Corollary~\ref{cor:mass-certificate}
is a separate theoretical comparator.  The mass quantity stored by the
reported implementation is only a heuristic diagnostic and is never used for
acceptance.

\subsection{Efficiency}
\label{ssec:efficiency}

Reliability alone permits a certificate that is valid but useless.  We now
show that $\etaSharp$ is also bounded \emph{by} the error in the natural
discrete space--time norm.  For an arbitrary SPD-gated assembled query the
constant is the directly computable matrix quantity $\kappa_h$; under exact
Galerkin assembly the variational estimates bound $\kappa_h$ uniformly in
$h$ and $\mu\in\Pset$.  Define
\begin{equation}
  |||\ee|||^2
  :=\dt\sum_{n=0}^{N_T-1}
  \Bigl(
   \bigl\lVert \Mh\,\partial\ee^{\,n+1/2}\bigr\rVert_{\Sh^{-1}}^2
   +\norm{\overline\ee^{\,n+1/2}}_{\Sh}^2
  \Bigr),
  \qquad
  \partial\ee^{\,n+1/2}:=\frac{\ee^{n+1}-\ee^{n}}{\dt}.
  \label{eq:triple-norm}
\end{equation}
The two terms are the discrete analogues of
$\norm{\partial_\tau e}_{L^2(0,T;\Vmu')}$ and
$\norm{e}_{L^2(0,T;\Vmu)}$, i.e.\ exactly the norm in which the parabolic
problem is well posed.

\begin{theorem}[Efficiency of the sharp certificate]
\label{thm:efficiency}
Let an assembled query satisfy $\Sh(\mu)\succ0$, and let $\kappa_h(\mu)$ be
defined by \eqref{eq:kappa-def}.  Then
\begin{equation}
  \etaSharp^2(N_T;\mu,\psi)
  \le 2\max\bigl\{1,\kappa_h(\mu)^2\bigr\}\,|||\ee|||^2 .
  \label{eq:efficiency}
\end{equation}
If, in addition, the matrices satisfy the exact-Galerkin-assembly hypothesis
of Lemma~\ref{lem:Sh-spd}, then
\begin{equation}
  \etaSharp^2(N_T;\mu,\psi)
  \le 2\max\Bigl\{1,\bigl(C_{\rm cont}/c_1\bigr)^2\Bigr\}\,|||\ee|||^2 ,
  \label{eq:efficiency-uniform}
\end{equation}
and the latter constant is independent of $h$, of $\dt$, and of
$\mu\in\Pset$.
\end{theorem}

\begin{proof}
By \eqref{eq:discrete-error-equation} and the triangle inequality in the
$\Sh^{-1}$ norm,
\[
  \norm{\rr_\theta^{n+1/2}}_{\Sh^{-1}}
  \le\bigl\lVert \Mh\,\partial\ee^{\,n+1/2}\bigr\rVert_{\Sh^{-1}}
   +\bigl\lVert \Ah\overline\ee^{\,n+1/2}\bigr\rVert_{\Sh^{-1}}
  \le\bigl\lVert \Mh\,\partial\ee^{\,n+1/2}\bigr\rVert_{\Sh^{-1}}
   +\kappa_h(\mu)\norm{\overline\ee^{\,n+1/2}}_{\Sh}.
\]
Squaring, using $(a+b)^2\le2a^2+2b^2$, multiplying by $\dt$ and summing gives
\eqref{eq:efficiency}.  Under exact assembly,
Lemma~\ref{lem:kappa} gives $\kappa_h\le C_{\rm cont}/c_1$, which yields
\eqref{eq:efficiency-uniform}.
\end{proof}

\begin{remark}[Reading Theorem~\ref{thm:efficiency} correctly]
\label{rem:efficiency-reading}
Theorems~\ref{thm:fully-discrete-certificate} and \ref{thm:efficiency}
together give
\[
  \norm{\ee^m}_{\Mh}^2+\dt\sum\norm{\overline\ee^{\,n+1/2}}_{\Sh}^2
  \ \le\ \etaSharp^2\ \le\ 2\max\{1,\kappa_h^2\}\,|||\ee|||^2 ,
\]
These displayed bounds compare two different norms: the efficiency
bound contains the discrete derivative term, whereas the stated reliability
bound contains the terminal mass error and midpoint energy error.  For $\ee^0=0$, a graph-norm bound in the other direction follows as well.
The error equation and $(a+b)^2\le2a^2+2b^2$ give
\[
 \dt\sum_n\|\Mh\partial\ee^{n+1/2}\|_{\Sh^{-1}}^2
 \le2\etaSharp^2+2\kappa_h^2\dt\sum_n\|\overline\ee^{n+1/2}\|_{\Sh}^2.
\]
Reliability then implies
\[
 |||\ee|||^2\le (3+2\kappa_h^2)\etaSharp^2.
\]
Thus graph-norm equivalence can be obtained with matrix-dependent constants;
this is distinct from equivalence with maximum-in-time spatial error.
Near-unit observed ratios are informative numerical observations, not
automatic identities for the nonsymmetric operator.
\end{remark}

\begin{remark}[Historical graph-norm diagnostics]
The archived efficiency routine uses the full matrices, whereas its
evolution residual has prescribed rows removed. These conventions must not
be substituted into a proof that assumes the complete error equation on
one space. The efficiency theorem here applies to the consistent free-space
equation \eqref{eq:constrained-cn}; the reported legacy ratios are numerical
diagnostics, and require recomputation with $M_I,A_I,S_I$ for direct
theorem validation.
\end{remark}

\subsection{Preconditioned certificates}
\label{ssec:preconditioned}

Between the mass and sharp certificates lies a family of preconditioned
bounds.  Suppose $P_h(\mu)$ is symmetric positive definite and a spectral
lower bound
\begin{equation}
  \Sh(\mu)\succeq\gamma_P(\mu)P_h(\mu),\qquad\gamma_P(\mu)>0,
  \label{eq:preconditioner-lower}
\end{equation}
is available.  Then
$\rr^\top\Sh^{-1}\rr\le\gamma_P(\mu)^{-1}\rr^\top P_h^{-1}\rr$.  Natural
candidates are a circulant approximation of the Toeplitz blocks or a banded
truncation of $\Sh$.

\begin{remark}
Lemma~\ref{lem:guaranteed} makes \eqref{eq:preconditioner-lower} largely
unnecessary as a \emph{certification} device: a preconditioner is better used
to accelerate the Riesz solve inside $\etaGuar$, where no spectral lower bound
is needed for rigour, than to produce a separate certified bound requiring a
computable $\gamma_P$.  This separates the role of preconditioning (iteration reduction) from the
role of certification (the guaranteed upper bracket).
\end{remark}

The reported preconditioner comparison is discussed once in
Section~\ref{ssec:runtime-scaling}. A preconditioner reduces iteration work;
it does not establish the spectral enclosure or remove per-query assembly.

\subsection{Discrete derivative and pointwise error bounds}
\label{ssec:greeks-cert}

Because the prediction is a spline in the trusted space, the certificate
propagates to spatial derivatives through discrete inverse inequalities whose
constants are computable once per mesh and are independent of $\mu$.  Let
\begin{equation}
  (R_h^{(1)})_{kk'}=(\phi'_{h,k'},\phi'_{h,k})_{L^2(\Omega)},
  \qquad
  (R_h^{(2)})_{kk'}=(\phi''_{h,k'},\phi''_{h,k})_{L^2(\Omega)},
  \label{eq:derivative-gram}
\end{equation}
both banded, and define
\begin{equation}
  \Lambda_j(h):=\lambda_{\max}\bigl(\Mh^{-1}R_h^{(j)}\bigr),
  \qquad j=1,2 ,
  \label{eq:inverse-constants}
\end{equation}
the largest eigenvalues of banded symmetric generalised eigenvalue problems.
By scaling, $\Lambda_1=O(h^{-2})$ and $\Lambda_2=O(h^{-4})$; $R_h^{(2)}$ requires $H^2$ regularity of the output space;
for the cardinal order convention this holds for $p\ge3$.

\begin{corollary}[Certified Delta and Gamma]
\label{cor:greeks}
Assume the hypotheses of Theorem~\ref{thm:fully-discrete-certificate} and
\eqref{eq:zero-initial-error}, and assume $V_h^0\subset H^2(\Omega)$ (order $p\ge3$ for the
cardinal construction). Then for every $m$,
\begin{equation}
  \norm{\partial_xV_h^m-\partial_xV_{\theta,h}^m}_{L^2(\Omega)}
  \le\sqrt{\Lambda_1(h)}\;\etaSharp(N_T),
  \qquad
  \norm{\partial_{xx}V_h^m-\partial_{xx}V_{\theta,h}^m}_{L^2(\Omega)}
  \le\sqrt{\Lambda_2(h)}\;\etaSharp(N_T).
  \label{eq:greek-L2}
\end{equation}
Consequently, with $\Delta=e^{-x}\partial_xV$ and
$\Gamma=e^{-2x}(\partial_{xx}V-\partial_xV)$,
\begin{equation}
  \norm{\Delta_h^m-\Delta_{\theta,h}^m}_{L^2(\Omega)}
  \le e^{-x_{\min}}\sqrt{\Lambda_1}\;\etaSharp,
  \qquad
  \norm{\Gamma_h^m-\Gamma_{\theta,h}^m}_{L^2(\Omega)}
  \le e^{-2x_{\min}}\bigl(\sqrt{\Lambda_2}+\sqrt{\Lambda_1}\bigr)\etaSharp .
  \label{eq:greek-financial}
\end{equation}
Finally, write $\Phi_h(x)=(\phi_{h,k}(x))_{k\in J_h}$ and define
\begin{equation}
 \gamma_\infty(h)=\sup_{x\in\Omega}
       \sqrt{\Phi_h(x)^\top M_h^{-1}\Phi_h(x)}.
 \label{eq:gamma-infty}
\end{equation}
Then the discrete price error obeys
\begin{equation}
 \|V_h^m-V_{\theta,h}^m\|_{L^\infty(\Omega)}
 \le\gamma_\infty(h)\,\etaSharp(N_T).
 \label{eq:greek-Linf}
\end{equation}
At a specified spot $S_0=e^{x_0}$, the smaller evaluation factor
$\sqrt{\Phi_h(x_0)^\top M_h^{-1}\Phi_h(x_0)}$ suffices. For a general
trajectory replace $\etaSharp$ by the bound including initial discrepancy;
for inexact Riesz evaluation use the corresponding upper bracket.
\end{corollary}

\begin{proof}
By Lemma~\ref{lem:lifting-cancels} the lifting cancels, so
$\partial_x^{(j)}(V_h^m-V_{\theta,h}^m)$ is the $j$-th derivative of the
spline with coefficient vector $\ee^m$.  Hence
$\norm{\partial_x^{(j)}(V_h^m-V_{\theta,h}^m)}_{L^2(\Omega)}^2
=(\ee^m)^\top R_h^{(j)}\ee^m\le\Lambda_j(\ee^m)^\top\Mh\ee^m
=\Lambda_j\norm{\ee^m}_{\Mh}^2$, and
$\norm{\ee^m}_{\Mh}\le\etaSharp$ by
Theorem~\ref{thm:fully-discrete-certificate}.  The financial forms follow from
$|e^{-x}|\le e^{-x_{\min}}$ on $\Omega$ and the triangle inequality.  For \eqref{eq:greek-Linf}, Cauchy--Schwarz in the mass pairing gives
$|\Phi_h(x)^\top e|\le
\sqrt{\Phi_h(x)^\top M_h^{-1}\Phi_h(x)}\,\|e\|_{M_h}$.
This expression applies to the actual basis rather than assuming cardinal
support overlap. For a polynomial spline basis its squared factor is a
piecewise polynomial; cellwise extremum calculations give the supremum,
and validated root isolation or interval bounds are needed for rigorous
numerical constants.
\end{proof}

\begin{remark}
The factors $\sqrt{\Lambda_1}=O(h^{-1})$ and $\sqrt{\Lambda_2}=O(h^{-2})$ are
the honest price of differentiating a certificate stated in $L^2$.  They mean
that certified Gamma is meaningful only when $\etaSharp\ll h^2$, which is a
demanding but checkable condition and precisely the kind of statement a
practitioner needs.  The global financial weights can be very conservative on a wide log-price
domain. Local derivative or evaluation functionals should be bounded in the
actual discrete geometry when practical Greek accuracy is required. No
useful numerical Greek-acceptance rate is established here.
\end{remark}

\subsection{Total error relative to the continuum problem}
\label{ssec:total-error}

Because the prediction lies in the same space as the trusted solution,
\begin{equation}
  V(\tau_m)-V_{\theta,h}^m
  =\bigl(V(\tau_m)-V_h^m\bigr)+\bigl(V_h^m-V_{\theta,h}^m\bigr),
  \label{eq:total-error-split}
\end{equation}
where the second term is certified and the first is the deterministic
discretisation error.

\begin{corollary}[Discretisation plus surrogate error]
\label{cor:total-error}
Under the hypotheses of \eqref{eq:deterministic-error}, exact conforming
assembly, Theorem~\ref{thm:fully-discrete-certificate}, and $\ee^0=0$,
with the same prescribed lifting
in both models,
\begin{equation}
 \norm{V(\tau_m)-V_{\theta,h}^m}_{L^2(\Omega)}
 \le C_\mu\bigl(h^{p-Y/2}+\dt^2+\varepsilon_b\bigr)
   +\etaSharp(N_T;\mu,\psi).
 \label{eq:total-error-bound}
\end{equation}
\end{corollary}

\begin{remark}[Terminology discipline]
\label{rem:terminology}
The surrogate term $\etaSharp$ in \eqref{eq:total-error-bound} is computable,
and the online implementation replaces it by the fully computable guaranteed
upper bound $\etaGuar$.  The first term in \eqref{eq:total-error-bound}
inherits the status of the deterministic analysis, including the constant
$C_\mu$ and the buffer quantity $\varepsilon_b$.  We therefore use the phrase
\emph{fully computable certificate for the surrogate error relative to the
trusted Galerkin discretisation}, not \emph{certified continuum price}.
\end{remark}

\subsection{Effectivity}
\label{ssec:effectivity}

For test queries where the reference is available, define
\begin{equation}
  E_m:=\norm{U_h^m-U_{\theta,h}^m}_{L^2(\Omega)},
  \qquad
  I_{\rm eff}:=\frac{\etaSharp(N_T)}{\max_mE_m},
  \label{eq:effectivity}
\end{equation}
with the analogous definitions for $\etaMass$ and $\etaGuar$.  Reliability
gives $I_{\rm eff}\ge1$ up to floating-point error, and
Theorem~\ref{thm:efficiency} bounds $I_{\rm eff}$ from above whenever the
discrete time-derivative term of \eqref{eq:triple-norm} is comparable to the
error itself.

Section~\ref{ssec:certificate-results} reports the full effectivity
statistics.  The guaranteed certificate has zero observed violations on the
certifiable main test set; the mass quantity plotted in the experiments is
explicitly labelled as a heuristic indicator and is never used to accept a
query.

% =================================================================
\section{Predict--certify--correct}
\label{sec:pcc}
% =================================================================

A certificate that only appears in a scatter plot has not been used.  In this
section the certificate becomes the online decision variable.

\subsection{The algorithm}
\label{ssec:algorithm}

The reported implementation uses a short Riesz-bracket evaluation as
a screen. Its numerical cost and validated-arithmetic status must be
measured separately.  A mass-based residual indicator may be logged for diagnostics, but it
never authorises acceptance.

\begin{figure}[ht]
\centering
\begin{minipage}{0.96\linewidth}
\textbf{Tolerance-checked predict--certify--correct rule.}
\begin{enumerate}[leftmargin=2em]
\item Construct the initial target, predict the coefficient trajectory, and
overwrite the prescribed coefficients by the specified lifting at every time
level (the pre-projection discrepancy is kept as a diagnostic). Compute
$\beta_0=\|\zz_\theta^0-\zz_h^0\|_{M_I}$.
\item Validate the time grid (uniform, pure CN in this version; other schemes
are rejected), and form free residuals with the same operator, stored forcing
and lifting as the reference.
\item Establish a positive spectral lower bound. If it cannot be
established, return status \textsc{uncertifiable} and invoke a separately implemented
deterministic fallback if a price is required; do not report a neural certificate.
\item Compute the augmented upper bound $\eta_{\rm aug}=(\beta_0^2+\eta_{\rm res}^2)^{1/2}$
with a short PCG budget. If it is at most $\varepsilon$, return
\textsc{direct accept} with the hash of the returned coefficients.
\item Otherwise correct by a free CN march from $\zz_h^0$, warm-started by the
candidate, record per-step convergence and linear residuals, and recertify the
corrected trajectory. If its bound passes, return \textsc{corrected accept}
(qualified when some linear solve did not converge).
\item If correction fails or its bound does not pass, run the recorded direct
fallback and recertify it (\textsc{qualified fallback}), or return
\textsc{over budget} or \textsc{numerical failure}. No output is delivered
without a passing certificate of that exact output.
\end{enumerate}
\end{minipage}
\caption{Tolerance-checked online rule. The status distinguishes a returned
approximation from a returned approximation certified to the requested budget.}
\label{alg:pcc}
\end{figure}

The historical experiments use $k_q=5$ and $k_{\max}=200$; the integrated
workflow uses $k_q=10$, $k_{\max}=200$, no preconditioner and
$\varepsilon=10^{-3}$, with GMRES tolerance $10^{-10}$ for correction.  The reported acceptance logic uses the upper bracket of
Lemma~\ref{lem:guaranteed}. Under its exact-arithmetic hypotheses,
incomplete PCG iterations affect sharpness rather than validity.
Remark~\ref{rem:arithmetic-scope} states the additional requirements for
validated floating-point acceptance.

\subsection{Soundness of the online rule}
\label{ssec:two-stage}

\begin{proposition}[Guaranteed acceptance]
\label{prop:two-stage}
Assume the hypotheses of Theorem~\ref{thm:fully-discrete-certificate},
$\ee^0=0$, a valid positive lower bound, and exact residual evaluation.
If the rule in Figure~\ref{alg:pcc} returns
\textup{\textsc{accepted}} with reported value $\etaGuar\le\varepsilon$, then
\begin{equation}
  \max_{1\le m\le N_T}\norm{V_h^m-V_{\theta,h}^m}_{L^2(\Omega)}
  \le \etaGuar\le\varepsilon .
  \label{eq:soundness}
\end{equation}
The conclusion is independent of the PCG stopping error. The same proof
applies to a corrected trajectory whose recomputed bound passes; nonzero
initial and bounded evaluation discrepancies use the augmented bound of
Remark~\ref{rem:arithmetic-scope}.
\end{proposition}

\begin{proof}
At every time step Lemma~\ref{lem:guaranteed} bounds
$\rr^\top\Sh^{-1}\rr$ from above by the computable upper bracket used in
$\etaGuar$.  Summing these bounds and applying
Theorem~\ref{thm:fully-discrete-certificate} proves
\eqref{eq:soundness}.  The mass-based indicator is absent from the proof and
from the acceptance logic.
\end{proof}

\subsection{Correction}
\label{ssec:correction}

Let $\dd^n$ be a correction trajectory, $\cc^n_{\rm corr}=\cc_\theta^n+\dd^n$.
The exact correction recovering the reference trajectory solves
\begin{equation}
  \Mh\frac{\dd^{n+1}-\dd^n}{\dt}+\Ah\overline\dd^{\,n+1/2}
  =-\rr_\theta^{n+1/2},
  \qquad \dd^0=0,
  \label{eq:correction-equation}
\end{equation}
that is, with $B_\pm=\Mh\pm\tfrac{\dt}{2}\Ah$,
\begin{equation}
  B_+\dd^{n+1}=B_-\dd^n-\dt\,\rr_\theta^{n+1/2}.
  \label{eq:correction-step}
\end{equation}
Solving \eqref{eq:correction-equation} exactly is equivalent to recomputing
the reference solution, so the practical aim is to apply a small number of
preconditioned or Krylov iterations warm-started at $\dd^n=0$, i.e.\ at the
neural prediction.

\begin{lemma}[Exact correction is exact]
\label{lem:exact-correction}
The solution of \eqref{eq:correction-equation} satisfies
$\cc_\theta^n+\dd^n=\cc_h^n$ for every $n$, and the corrected trajectory has
zero residual and zero certificate.
\end{lemma}

\begin{proof}
Both $\{\cc_h^n\}$ and $\{\cc_\theta^n+\dd^n\}$ satisfy
\eqref{eq:cn-reference} with the same initial value.  Under the certification
hypothesis $\Sh\succ0$, the symmetric part
$\Mh+\tfrac{\dt}{2}\Sh$ of $B_+$ is positive definite, hence $B_+$ is
nonsingular.  Uniqueness gives the claim, and Lemma~\ref{lem:consistency} the
rest.
\end{proof}

\begin{proposition}[Contraction of a preconditioned correction]
\label{prop:correction-contraction}
Fix a time step and write the step equation as $B_+\zz=\bb$.  Let $P\succ0$ be
a preconditioner and consider the Richardson iteration
$\zz^{(i+1)}=\zz^{(i)}+P^{-1}(\bb-B_+\zz^{(i)})$ with residuals
$\sres^{(i)}=\bb-B_+\zz^{(i)}$.  Then
\begin{equation}
  \norm{\sres^{(i+1)}}_{P^{-1}}\le\rho_P\,\norm{\sres^{(i)}}_{P^{-1}},
  \qquad
  \rho_P:=\bigl\lVert I-P^{-1/2}B_+P^{-1/2}\bigr\rVert_2 .
  \label{eq:contraction}
\end{equation}
If $P=\Mh+\tfrac{\dt}{2}\Sh$, the symmetric part of $B_+$, then
$I-P^{-1/2}B_+P^{-1/2}=-\tfrac{\dt}{2}P^{-1/2}N_hP^{-1/2}$ with
$N_h=\tfrac12(\Ah-\Ah^\top)$ skew symmetric, so $\rho_P$ is determined
entirely by the nonsymmetric part of the CGMY operator and satisfies
$\rho_P\le\tfrac{\dt}{2}\lVert P^{-1/2}N_hP^{-1/2}\rVert_2$.
\end{proposition}

\begin{proof}
$\sres^{(i+1)}=(I-B_+P^{-1})\sres^{(i)}$, so
$P^{-1/2}\sres^{(i+1)}=(I-P^{-1/2}B_+P^{-1/2})P^{-1/2}\sres^{(i)}$ and
\eqref{eq:contraction} follows by taking Euclidean norms.  For the second
claim, $B_+=P+\tfrac{\dt}{2}N_h$ by \eqref{eq:Sh} and the definition of $N_h$.
\end{proof}

\begin{remark}[Status of Proposition~\ref{prop:correction-contraction}]
\label{rem:contraction-status}
The proposition supplies a Richardson residual estimate. Contraction is
\emph{conditional}: $\rho_P<1$ is a
hypothesis, not a conclusion.  It is, however, a computable hypothesis, since an extreme singular-value calculation can estimate $\rho_P$;
a few iterations do not by themselves certify an upper bound. GMRES minimises
the residual in its prescribed norm over a growing Krylov space in exact
arithmetic; restarted or preconditioned implementations require their own
precise interpretation. No Richardson contraction or timing benefit is
inferred from the reported GMRES iteration counts.  We do not claim an unconditional contraction
proof.
\end{remark}

The correction is implemented with the same matrix-free Toeplitz/FFT
operator action as the reference solver. In the historical warm-start study
the neural initialisation lowered the mean total GMRES iteration count by
$7.4\%$ without a wall-clock gain. The matched benchmark of
Section~\ref{ssec:benchmark} confirms this at the workflow level: the
DeepONet-warm-started free march is $1.10\times$ (cold) and $1.33\times$
(reusable setup) the time of a free dense direct solve of the same problem.

\subsection{Cost model}
\label{ssec:cost}

Let $T_{\rm shared}$ include per-query operator assembly, forcing,
initial projection and shared solver setup. Let $k_B$ be the average Krylov
work per time step and $k_S$ the Riesz work. An explicit cost decomposition is
\begin{equation}
\begin{aligned}
 T_{\rm ref}&=T_{\rm shared}+T_{\rm ref,setup}+T_{\rm advance},\\
 T_{\rm accepted}&=T_{\rm shared}+T_{\rm infer}+T_{\rm residual}
                   +T_{\alpha}+T_{\rm Riesz},\\
 T_{\rm corrected}&=T_{\rm accepted}+T_{\rm correction}
                         +T_{\rm recertification}.
\end{aligned}
\label{eq:cost-model}
\end{equation}
FFT iteration work is $O(k_BN_TN_h\log N_h)$ or
$O(k_SN_TN_h\log N_h)$ only for the corresponding matrix-action phases.
Dense assembly and spectral decomposition have different costs and are not
covered by this count. A dense LU reference factorises $B_+$ once per query
and reuses it over time; its cost must be measured, especially at small sizes.
An optimised cold solve using exactly the correction machinery is necessary
to isolate the value of neural initialisation.

The benchmark of Section~\ref{ssec:benchmark} reports per-query ratios
$T_{{\rm method},i}/T_{{\rm baseline},i}$, summarised by their geometric
mean over the queries on which both methods meet the common target, with
bootstrap intervals over queries and coverage. For reusable setup the cost
of $Q$ queries is $C_0+QC_1$, with one-time mesh-level cost $C_0$ and online
cost $C_1$. A break-even count against a baseline,
\[
 Q_{\rm break-even}=\max\Bigl\{1,\Bigl\lceil\frac{C_0-C_0^{\rm base}}
 {C_1^{\rm base}-C_1}\Bigr\rceil\Bigr\},
\]
is reported only when the method's verified coverage is not lower than the
baseline's and $C_1<C_1^{\rm base}$ on the common verified queries with a
paired bootstrap interval of the difference excluding zero; offline data generation and training
are reported separately.

In the resolution study of Section~\ref{ssec:scaling-results}, operator
assembly and certification dominate the cost at every tested size, and the
time-stepping solve is at most a few percent of the total; the dense spectral
gate and the residual/PCG certificate grow faster than linearly at $N=512$.
The implementation therefore does not have a wholly matrix-free certification
path, regardless of its PCG matrix actions.

Certificate-driven data enrichment is a natural extension, but it is not part
of the evidence set used in this paper.  We therefore do not make an adaptive-
sampling claim here.

% =================================================================
\section{Dataset and implementation}
\label{sec:data}

The implementation wraps the trusted B-spline Galerkin solver rather than
reimplementing it.  Every stored trajectory contains the parameter vector,
payoff representation, projected initial coefficients, coefficient trajectory,
forcing data needed to reproduce the residual, and mesh/time-grid metadata.
The production numerical configuration uses quadratic cardinal splines
(order $p=3$), $N=64$ elements over
$\Omega=(\log100-8,\log100+8)$, 66 stored coefficients (62 free, four
prescribed), and 60 pure CN steps. The neural input includes $T\in[0.1,1]$;
the spatial-operator parameter remains $\mu=(C,G,M,Y,r)$.
Neural and reduced predictions are projected onto the prescribed coefficients
before residual evaluation, and all certification and correction use the
constrained free coordinates of Section~\ref{ssec:constrained-system}.


The reported dataset contains $800$ successfully generated trajectories,
sampled by a Sobol design from the frozen box in
Table~\ref{tab:proposed-box}; the fixed training seed is $0$.  Independent
Sobol seeds are used for test, stress and learning-curve experiments.  The
main interpolation experiment uses $100$ unseen queries, the main
certificate experiment uses $100$ sampled queries, the OOD experiment uses
$50$ queries in each of four tiers, and the resolution-transfer experiment
uses $50$ parameter points at each resolution.  Three independent network
initialisation seeds are used for the DeepONet variance study.

The candidate models are:
\begin{table}[ht]
\centering
\caption{Architectures used in the numerical comparison.  Parameter counts
exclude the non-trainable POD basis.}
\label{tab:architectures}
\small
\resizebox{\textwidth}{!}{%
\begin{tabular}{llll}
\toprule
Model & Representation & Main size & Trainable parameters\\
\midrule
MLP & $(\tau,\mu,T,\psi_h)\mapsto\Vh$ & $128$--$128$--$128$ & $51{,}010$\\
DeepONet & payoff/parameter branches + time trunk & latent dim. $64$ & $80{,}002$\\
POD--Galerkin (historical) & reduced dynamics without lifting & $32$ modes & --\\
Mass-weighted POD (v5) & free CN dynamics in $M_I$-POD space & $r=8$--$62$ & --\\
Projected FNO & grid FNO followed by $\Pi_h$ & width $32$, $16$ modes, $4$ layers & $136{,}673$\\
\bottomrule
\end{tabular}}
\end{table}
The MLP and DeepONet use
$\cc_\theta(\tau)=\cc_h^0+\tau N_\theta(\tau,\mu,\psi_h)$.
The historical POD and projected FNO used an analogous $c_h^0+[g(\tau)-g(0)]$
construction. The historical POD model also evolved incorrect reduced
equations (no lifting forcing); it is superseded by the repaired
mass-weighted model of Section~\ref{ssec:pod-results}, whose nonzero
finite-rank initial error $\beta_0$ enters its augmented bound.  Checkpoint compatibility checks reject a
model whose recorded mesh, time grid or parameter box differs from the loaded
configuration.

For certification, the sharp estimator is computed by direct factorisation at
moderate sizes only and is used as a reference diagnostic.  The online
quantity is $\etaGuar$, obtained from the upper Riesz bracket with PCG.  The
reported choice is no preconditioner and at most $200$ PCG iterations;
quick online screening uses ten iterations (five in the historical runs).  A separate mass-based residual
quantity is retained in plots only as a heuristic indicator and cannot accept
a query.  The implementation also computes a conservative spectral lower
bound used to decide whether the query is certifiable; failed gates are
reported rather than silently discarded.

The complete test suite (290 tests) passes on the evidence revision. The
earlier failing ill-conditioned $A=S$ diagnostic ($\kappa=1.0000000015738064$
against a $10^{-10}$ tolerance at $\mathrm{cond}(S)=10^8$) was reproduced and
explained: the stored test matrix is not exactly symmetric, so the represented
problem has exact $\kappa=1+2.9\times10^{-10}$ (40-digit computation on the same
entries), and the remaining deviation lies far inside the first-order forward
error $n\varepsilon_{\rm mach}\mathrm{cond}(S)$; Cholesky and triangular solves meet
their componentwise backward-error bounds. The test was replaced by a strict
well-conditioned identity test and an ill-conditioned test with a tolerance
derived from conditioning; no $A=S$ shortcut exists in the code. No neural
model was retrained.

Recovered training documentation records 800 training trajectories and a
separate 30-query validation set (seed 999999). MLP and DeepONet use Adam
with learning rate $2\times10^{-3}$ and cosine scheduling; projected FNO uses
60 epochs with trajectory batches of eight and grid-space MSE, while POD
uses a single snapshot SVD (historical model; the repaired basis is an SVD
of $L^\top Z$ over the same 800 trajectories, whose artefact hash is
recorded in each run manifest). The saved networks are final-epoch models,
without validation-based checkpoint selection. The documented training
times are 372.3 s for DeepONet, 295.2 s for projected FNO and 0.113 s for
the historical POD SVD, excluding data generation. These details are recovered from the
submitted documentation; its referenced training entry script is absent
from this archive, so complete training reproduction remains outstanding.
Stale textual configuration snapshots in checkpoints must not override their
separately recorded realised discretisation.

Every v5 run records a manifest (source and configuration hashes, git state,
checkpoints and bases, query-population hashes, software, hardware, thread
settings and seeds) and is single-threaded. The benchmark pins each timing
process to one performance core, uses one warmup on a query outside every
population and five timed repetitions with a reproducible per-query method
order, and isolates caches as described in Section~\ref{ssec:benchmark}.
Historical and v5 workloads have separate setup and eligibility rules and are
never combined.


% =================================================================
\section{Numerical experiments}
\label{sec:numerics}

The experiments fall into two evidence tiers that are never pooled.
Sections~\ref{ssec:audit-results}--\ref{ssec:fourier-financial} report the
v5 evidence: every run uses the integrated constrained free workflow
(method identifier \texttt{constrained\_free\_cn\_v1}), records a manifest,
and is summarised by scripts that read only its raw records; every table in
these sections is generated from those records
(Appendix~\ref{app:evidence} lists the runs). Section~\ref{ssec:historical}
retains the earlier neural-model studies, which used the legacy full-space
certificate and the defective historical POD model; their accuracy
observations are kept, with their original populations, and their speed
claims are withdrawn. All v5 runs use the production discretisation
$N=64$, $p=3$, $60$ CN steps, $K=100$ and the parameter box of
Table~\ref{tab:proposed-box}, unless a study states otherwise.

\paragraph{Populations.} Each study is scoped to its own predeclared,
hashed population, and statistics are not transferred between them:
\texttt{Q\_audit} (459 stored audit coordinates: 256 Sobol points of the
box, the 64 exact vertices and 139 deliberately selected legacy failures,
reported as separate strata); \texttt{Q\_bench120} (120 Sobol points used
only for the benchmark); \texttt{Q\_ref30} (30 pricing queries, five spots
$0.8K,\dots,1.2K$ each); the fixed 30-query POD validation set; and
\texttt{Q\_scale} (30 queries; a predeclared subset of six at $N=512$).
Statistics over spots are taken within a query first (maximum), and
bootstrap intervals resample queries, never repetitions or spots.

\paragraph{Target statuses.} A method's output is \emph{verified} when its
own certificate of the exact returned coefficients meets the common target;
\emph{empirically checked} when it is not verified but agrees with the
stored reference within the target after adding that reference's own
allowance; \emph{failed} or \emph{unavailable} otherwise. A delivered output
or a converged solver flag is never counted as verified.

% -----------------------------------------------------------------
\subsection{Positivity of the constrained operator}
\label{ssec:audit-results}

Table~\ref{tab:free-spd-audit} repeats the positivity gate on all 459 audit
coordinates in both coordinates. The free gate on $S_I$ passes everywhere;
the legacy gate on the full raw matrix $\bar S$ rejects 3 of 256 Sobol
points, 16 of 64 exact vertices and 135 of the 139 selected cases, and the
computed $\lambda_{\min}(\bar S)$ is negative on those queries. The free
operator is therefore the one the certificate requires
(Section~\ref{ssec:constrained-system}); the full-matrix failures are a
property of the prescribed-coefficient block, not of the constrained
problem. On the 305 queries where both gates pass, seeded fixed residuals
give $q_{\rm free}\le q_{\rm full}$ in every case, as the Schur-complement
argument predicts. The selected stratum is not a population sample, and the
sweep is a finite-sample floating-point diagnostic, not a uniform box
theorem.

% GENERATED by scripts/v6_build_manuscript_assets.py from raw CSVs -- do not edit by hand
\begin{table}[ht]
\centering\small
\caption{Positivity gates on the 459 stored audit coordinates at $N=64$ (run \texttt{20261007T175335Z\_m6\_audit459}). Gate = positive floating-point lower-bound estimate of $\lambda_{\min}$; minima are computed eigenvalues, not validated bounds.}
\label{tab:free-spd-audit}
\resizebox{\textwidth}{!}{%
\begin{tabular}{lrrrrrrr}
\toprule
Stratum & $n$ & free $S_I$ & full $\bar S$ & legacy-only rej. & $\min\lambda_{\min}(S_I)$ & $\min\lambda_{\min}(\bar S)$ & median $\kappa_I$\\
\midrule
Sobol box sample & 256 & 256/256 & 253/256 & 3 & $0.01579$ & $-0.0362$ & $1.45$\\
Exact box vertices & 64 & 64/64 & 48/64 & 16 & $0.01038$ & $-1.80$ & $1.78$\\
Selected legacy failures$^\ast$ & 139 & 139/139 & 4/139 & 135 & $0.01030$ & $-30.4$ & $1.52$\\
\bottomrule
\end{tabular}}
\par\smallskip\parbox{0.95\linewidth}{\footnotesize $^\ast$Deliberately selected cases; their legacy failure fraction is not a population failure rate. Fixed seeded residuals where both gates pass (305 queries): $q_{\rm free}/q_{\rm full}\in[0.167,0.999]$, Schur bound held in 305/305.}
\end{table}


% -----------------------------------------------------------------
\subsection{Riesz brackets and the augmented bound}
\label{ssec:bracket-results}

The PCG bracket of Lemma~\ref{lem:guaranteed} was compared with a 40-digit
evaluation of $r^\top S_I^{-1}r$ on the same assembled double-precision
matrix (Table~\ref{tab:riesz}), for residual norms from about $10^{-6}$
(stored and corrected trajectories) to about $10^{2}$ (raw DeepONet).
Among 11{,}760 comparisons there are no detected violations: $8{,}065$ are
resolved passes, and $3{,}695$ are numerically unresolved because a margin
lies within the comparison uncertainty $8n\varepsilon_{\rm mach}\,
\mathrm{cond}_2(S)\,q$ ($3{,}640$ on both sides, $55$ on the lower side
only; $3{,}621$ of them are converged iterates whose margins are zero up to
roundoff). Unresolved rows are not counted as confirmed passes. The
high-precision reference validates the bracket implementation for the
assembled matrix; it does not validate assembly accuracy and is not an
outward-rounded enclosure. The hypothesis $\alpha_{\rm lb}\le\lambda_{\min}(S_I)$
held on all 24 queries (finite-sample evidence).

Table~\ref{tab:augmented} tests the initial-error augmentation against the
true error to a dense direct solution. For exact dynamics started from a
perturbed initial vector, the residual is at roundoff and $\eta_{\rm res}$
alone underestimates the error in all 48 cases, whereas $\eta_{\rm aug}$
holds in all of them with effectivity one. For raw DeepONet and rank-16
POD candidates the bound holds in all cases with median effectivities 5.0
and 5.4.

% GENERATED by scripts/v6_build_manuscript_assets.py from raw CSVs -- do not edit by hand
\begin{table}[ht]
\centering\small
\caption{Riesz-bracket validation against a 40-digit reference on the same assembled $S_I$ (24 declared audit queries, steps $n\in\{0,30,59\}$, PCG budgets $\{0,\dots,200\}$, four preconditioners; manufactured: $n=62$, $\mathrm{cond}\le10^8$). Unresolved = a margin within the comparison uncertainty $8n\varepsilon_{\rm mach}\mathrm{cond}_2(S)q$ (both sides/lower only); unresolved rows are not confirmed passes.}
\label{tab:riesz}
\begin{tabular}{lrrrrr}
\toprule
Residual source & $\|r\|_2$ range & rows & resolved pass & unresolved & violations\\
\midrule
Stored reference & $7.1\times10^{-7}$--$1.5\times10^{-5}$ & 2304 & 1584 & 720 (720/0) & 0\\
Raw DeepONet candidate & $0.96$--$79$ & 2304 & 1553 & 751 (724/27) & 0\\
POD $r=16$ & $0.035$--$7.5$ & 2304 & 1569 & 735 (720/15) & 0\\
POD $r=32$ & $5.6\times10^{-4}$--$0.13$ & 2304 & 1571 & 733 (726/7) & 0\\
Corrected trajectory & $5.8\times10^{-7}$--$1.4\times10^{-5}$ & 2304 & 1579 & 725 (720/5) & 0\\
Manufactured SPD & $6.8$--$9.5$ & 240 & 209 & 31 (30/1) & 0\\
\textbf{All} &  & 11760 & 8065 & 3695 & 0\\
\bottomrule
\end{tabular}
\par\smallskip\parbox{0.95\linewidth}{\footnotesize $\alpha\le\lambda_{\min}(S_I)$ against the 40-digit eigenvalue of the same matrix held on 24/24 queries (finite-sample evidence). The reference validates the bracket implementation, not assembly accuracy or an outward-rounded enclosure.}
\end{table}

% GENERATED by scripts/v6_build_manuscript_assets.py from raw CSVs -- do not edit by hand
\begin{table}[ht]
\centering\small
\caption{Augmented bound $\eta_{\rm aug}=(\beta_0^2+\eta_{\rm res}^2)^{1/2}$ against the true maximum-in-time $M_I$ error to a dense direct solution (48 declared audit queries: first 16 per stratum).}
\label{tab:augmented}
\resizebox{\textwidth}{!}{%
\begin{tabular}{lrrrr}
\toprule
Candidate & $n$ & bound holds & effectivity med./max & $\eta_{\rm res}$ alone too small\\
\midrule
Exact dynamics, perturbed $z^0$ ($\|\delta\|_{M_I}=10^{-4}$) & 48 & 48/48 & $1.00$ / $1.00$ & 48/48\\
Raw DeepONet candidate & 48 & 48/48 & $4.98$ / $26.28$ & 0/48\\
POD $r=16$ candidate & 48 & 48/48 & $5.37$ / $10.65$ & 0/48\\
\bottomrule
\end{tabular}}
\end{table}


% -----------------------------------------------------------------
\subsection{Matched benchmark at the production resolution}
\label{ssec:benchmark}

The benchmark runs eight methods on all 120 \texttt{Q\_bench120} queries
in two workloads with one target each: the trajectory workload requires
$\eta_{\rm aug}\le10^{-3}$, and the terminal-price workload requires
$\max_s\beta_I(x_s)\eta_{\rm aug}\le\$0.01$ relative to the exact production
discrete price at the five spots. Every method is certified through the same
free operator, stored forcing and lifting. A comparison with the stored
iterative reference adds that reference's own certificate:
\begin{equation}
 \max_n\|\zz_\theta^n-\zz_{\rm ref}^n\|_{M_I}
 \le \eta_{{\rm aug},\theta}+\eta_{\rm ref}.
 \label{eq:reference-solve-allowance}
\end{equation}
Timings are single-threaded, with one warmup on a query outside every
population and five timed repetitions per query in a reproducible shuffled
method order. In the \emph{cold} mode each method, query and repetition
builds a fresh operator cache inside its timer. In the \emph{reuse} mode
only mesh-level objects (operator mesh data, the free mass factor and
point functionals) are reused; their one-time cost is reported separately
(Table~\ref{tab:amortised}). A time ratio is
$t_{\rm method}/t_{\rm baseline}$ per query (values below one are faster),
summarised by its geometric mean over the queries on which both methods are
verified, with a 95\% bootstrap interval over queries and the coverage of
each method.

% GENERATED by scripts/v6_build_manuscript_assets.py from raw CSVs -- do not edit by hand
\begin{table}[ht]
\centering\small
\caption{Trajectory workload on all 120 \texttt{Q\_bench120} queries, $N=64$, target $\eta_{\rm aug}\le10^{-3}$ (algebraic, relative to the exact production discrete solution). Times: median over queries of the median over 5 timed repetitions (single thread). Ratio $=t_{\rm method}/t_{\rm free\,LU}$ (values below 1 are faster): geometric mean over the common target-verified queries with a 95\% bootstrap interval over queries; coverage is the verified count. Cold: fresh operator cache per method, query and repetition, charged to the method. Reuse: only mesh-level objects reused (one-time cost excluded here, see Table~\ref{tab:amortised}).}
\label{tab:benchmark-trajectory}
\resizebox{\textwidth}{!}{%
\begin{tabular}{lrrrrrr}
\toprule
Method & verified & cold s & cold ratio [CI] & verified & reuse s & reuse ratio [CI]\\
\midrule
Free dense LU + certificate & 120/120 & $0.843$ & 1 & 120/120 & $0.271$ & 1\\
Free cold GMRES + certificate & 120/120 & $0.905$ & 1.08 [1.07, 1.08] & 120/120 & $0.335$ & 1.24 [1.23, 1.25]\\
Trusted reference solve + certificate & 120/120 & $1.487$ & 1.77 [1.77, 1.78] & 120/120 & $0.926$ & 3.45 [3.42, 3.48]\\
DeepONet, direct certificate only & 0/120 & $0.649$ & -- & 0/120 & $0.078$ & --\\
DeepONet-warm-started GMRES + certificate & 120/120 & $0.927$ & 1.10 [1.10, 1.10] & 120/120 & $0.358$ & 1.33 [1.32, 1.33]\\
DeepONet predict--certify--correct & 120/120 & $0.995$ & 1.18 [1.18, 1.19] & 120/120 & $0.425$ & 1.58 [1.57, 1.59]\\
Mass-weighted POD $r=32$ workflow & 120/120 & $0.950$ & 1.10 [1.08, 1.11] & 120/120 & $0.380$ & 1.25 [1.16, 1.34]\\
Mass-weighted POD $r=48$ workflow & 120/120 & $0.651$ & 0.78 [0.77, 0.78] & 120/120 & $0.080$ & 0.30 [0.30, 0.30]\\
\bottomrule
\end{tabular}}
\end{table}


Table~\ref{tab:benchmark-trajectory} and Figure~\ref{fig:benchmark-ratios}
show that every delivering method is verified on all 120 queries, except
direct certification of the raw DeepONet prediction, which is over budget
on all 120 ($\eta_{\rm aug}$ between about 2 and 13). The neural
predict--certify--correct workflow is slower than a free dense LU solve of
the same constrained problem: $1.18\times$ [1.18, 1.19] in the cold mode and
$1.58\times$ [1.57, 1.59] with reusable setup. Warm-starting GMRES with the
DeepONet prediction is also slower than dense LU ($1.10\times$ and
$1.33\times$), and slower than an unwarmed GMRES solve ($1.08\times$ and
$1.24\times$). Relative to
the trusted reference solve the neural workflow takes $0.67\times$ (cold)
and $0.46\times$ (reuse) of the time, but this measures the reference
implementation, not a neural acceleration.

The repaired mass-weighted POD model at rank 48 is the only method faster
than dense LU: $0.78\times$ [0.77, 0.78] cold and $0.30\times$ [0.30, 0.30]
with reuse, with direct acceptance on all 120 queries. At rank 32 the
candidate bound slightly exceeds the target on 111 of 120 queries, which
are corrected; the mixture of nine cheap direct accepts and 111 corrections
explains its wider reuse interval ($1.25\times$ [1.16, 1.34]). In the cold
mode operator assembly and certification, common to all methods, dominate
the time, which compresses every cold ratio towards one.

\begin{figure}[ht]
\centering
\includegraphics[width=0.95\linewidth]{figures/v6_benchmark_ratios.png}
\caption{Trajectory workload, $N=64$, 120 queries: geometric-mean time
ratio to free dense LU with 95\% bootstrap intervals over queries (log
scale). All plotted methods are verified on 120/120 queries; direct
certification of the raw DeepONet prediction is verified on none and has no
ratio. Replaces the historical speed figures (Appendix~\ref{app:historical-figures}).}
\label{fig:benchmark-ratios}
\end{figure}

% GENERATED by scripts/v6_build_manuscript_assets.py from raw CSVs -- do not edit by hand
\begin{table}[ht]
\centering\small
\caption{Reusable-setup mode: one-time mesh-level setup $C_0$ (operator mesh data, free mass factor, point functionals) and online cost per query $C_1$ (mean of query medians); amortised cost $C_0/Q+C_1$ at $Q=120$. Break-even against free dense LU is reported only when the method's verified coverage is not lower and its online advantage has a bootstrap interval excluding zero.}
\label{tab:amortised}
\resizebox{\textwidth}{!}{%
\begin{tabular}{lrrrr}
\toprule
Method & $C_0$ (s) & $C_1$ (s) & $Q=120$ (s) & break-even vs LU\\
\midrule
Free dense LU + certificate & $0.575$ & $0.268$ & $0.273$ & n/a\\
Free cold GMRES + certificate & $0.568$ & $0.332$ & $0.337$ & none\\
Trusted reference solve + certificate & $0.565$ & $0.924$ & $0.928$ & none\\
DeepONet, direct certificate only & $0.568$ & $0.078$ & $0.083$ & none (verified 0/120 vs 120/120)\\
DeepONet-warm-started GMRES + certificate & $0.566$ & $0.355$ & $0.360$ & none\\
DeepONet predict--certify--correct & $0.564$ & $0.423$ & $0.428$ & none\\
Mass-weighted POD $r=32$ workflow & $0.566$ & $0.355$ & $0.360$ & none\\
Mass-weighted POD $r=48$ workflow & $0.568$ & $0.080$ & $0.085$ & 1\\
\bottomrule
\end{tabular}}
\end{table}


Amortising the mesh-level setup over the 120 queries changes none of these
conclusions (Table~\ref{tab:amortised}): the shared setup $C_0\approx0.57$ s
is the same for every method, so only POD at rank 48 has a break-even point
against dense LU, at the first query.

% GENERATED by scripts/v6_build_manuscript_assets.py from raw CSVs -- do not edit by hand
\begin{table}[ht]
\centering\small
\caption{Terminal-price workload (cold mode, 120 queries, spots $0.8K,\dots,1.2K$ aggregated by the maximum within each query). Target: within \$0.01 of the exact production discrete price. Verified: $\max_s\beta_I(x_s)\eta_{\rm aug}\le0.01$; empirically checked: not verified but within budget of the stored reference including its own allowance; unavailable: no price delivered.}
\label{tab:benchmark-price}
\resizebox{\textwidth}{!}{%
\begin{tabular}{lrrrr}
\toprule
Method & verified & emp. checked & failed & unavailable\\
\midrule
Free dense LU + certificate & 120 & 0 & 0 & 0\\
Free cold GMRES + certificate & 120 & 0 & 0 & 0\\
Trusted reference solve + certificate & 120 & 0 & 0 & 0\\
DeepONet, direct certificate only & 0 & 0 & 0 & 120\\
DeepONet-warm-started GMRES + certificate & 120 & 0 & 0 & 0\\
DeepONet predict--certify--correct & 120 & 0 & 0 & 0\\
Mass-weighted POD $r=32$ workflow & 120 & 0 & 0 & 0\\
Mass-weighted POD $r=48$ workflow & 120 & 0 & 0 & 0\\
COS (independent continuum price) & 0 & 9 & 111 & 0\\
\bottomrule
\end{tabular}}
\par\smallskip\parbox{0.95\linewidth}{\footnotesize COS approximates the continuum price, not the discrete one: its gap to the stored $N=64$ discrete price has median \$0.031 and maximum \$0.31, a discretisation diagnostic (Section~\ref{ssec:refinement}), so it never qualifies under this target.}
\end{table}


In the terminal-price workload (Table~\ref{tab:benchmark-price}) the same
eight methods reach the same verified counts. The COS method is reported in
this workload only and approximates the continuum price, not the discrete
one: it is within \$0.01 of the stored discrete price on 9 of 120 queries,
with median gap \$0.031 and maximum \$0.31. That gap is a discretisation
diagnostic studied in Section~\ref{ssec:refinement}, not an error of either
method, and COS timings are therefore not compared with the discrete
methods under this target.

% -----------------------------------------------------------------
\subsection{The repaired POD baseline}
\label{ssec:pod-results}

The historical POD model reduced the full equations without the
prescribed-coefficient lifting and its forcing; its full-rank dynamics error
of about $9.46$ is an implementation defect, and the historical POD
comparisons are withdrawn. The repaired model reduces the free CN system
\eqref{eq:constrained-cn}, forcing included. With $M_I=LL^\top$ and the
uncentred free snapshot matrix $Z$, the basis is $V=L^{-\top}U_r$, where
$U_r$ holds the leading left singular vectors of $L^\top Z$; this is the
mass-weighted optimal POD subspace, and $V^\top M_IV=I$. The reduced CN
equations are
\begin{equation}
 \frac{\alpha^{n+1}-\alpha^n}{\dt}
 + (V^\top A_I V)\frac{\alpha^{n+1}+\alpha^n}{2}
 = V^\top f_I^{n+1/2},\qquad
 \alpha^0=V^\top M_I \zz_h^0,
 \label{eq:repaired-pod-cn}
\end{equation}
and the candidate $V\alpha^n$ is certified with its nonzero finite-rank
initial error $\beta_0$. A Euclidean-SVD basis, mass-orthonormalised, is
kept as a labelled diagnostic.

% GENERATED by scripts/v6_build_manuscript_assets.py from raw CSVs -- do not edit by hand
\begin{table}[ht]
\centering\small
\caption{Repaired free POD baseline on the fixed 30-query validation set (training: 800 production trajectories, 48{,}800 uncentred free snapshots). Errors in the $M_I$ norm against the stored reference: projection and $E_{\rm rel}$ are maxima over time, $\beta_0$ is the finite-rank initial error. Last column: direct / corrected acceptance at $10^{-3}$ through the production free workflow (all 30 delivered).}
\label{tab:repaired-pod}
\resizebox{\textwidth}{!}{%
\begin{tabular}{lrrrrrrr}
\toprule
basis & $r$ & projection & $\beta_0$ & $E_{\rm rel}$ med./max & price err. (\$) & cand. $\eta_{\rm aug}$ & dir./corr.\\
\midrule
mass-weighted & 8 & $0.332$ & $0.332$ & $1.33\times10^{-3}$ / $4.17\times10^{-3}$ & $0.0395$ & $2.96$ & 0/30\\
mass-weighted & 16 & $0.0235$ & $0.0235$ & $1.05\times10^{-4}$ / $4.37\times10^{-4}$ & $4.57\times10^{-4}$ & $0.252$ & 0/30\\
mass-weighted & 32 & $3.34\times10^{-4}$ & $1.44\times10^{-4}$ & $1.36\times10^{-6}$ / $3.40\times10^{-6}$ & $1.14\times10^{-6}$ & $1.58\times10^{-3}$ & 0/30\\
mass-weighted & 48 & $4.16\times10^{-6}$ & $7.65\times10^{-7}$ & $1.72\times10^{-8}$ / $5.57\times10^{-8}$ & $1.66\times10^{-8}$ & $1.63\times10^{-5}$ & 30/0\\
mass-weighted & 62 & $2.58\times10^{-13}$ & $2.02\times10^{-13}$ & $1.06\times10^{-9}$ / $6.13\times10^{-9}$ & $1.54\times10^{-8}$ & $3.40\times10^{-11}$ & 30/0\\
Euclidean (diag.) & 8 & $0.286$ & $0.286$ & $1.15\times10^{-3}$ / $4.45\times10^{-3}$ & $0.0331$ & $2.80$ & 0/30\\
Euclidean (diag.) & 16 & $0.0230$ & $0.0197$ & $1.01\times10^{-4}$ / $4.37\times10^{-4}$ & $3.83\times10^{-4}$ & $0.229$ & 0/30\\
Euclidean (diag.) & 32 & $3.41\times10^{-4}$ & $1.42\times10^{-4}$ & $1.40\times10^{-6}$ / $5.03\times10^{-6}$ & $7.38\times10^{-7}$ & $1.98\times10^{-3}$ & 4/26\\
Euclidean (diag.) & 48 & $3.98\times10^{-6}$ & $7.76\times10^{-7}$ & $1.61\times10^{-8}$ / $5.10\times10^{-8}$ & $1.67\times10^{-8}$ & $1.51\times10^{-5}$ & 30/0\\
Euclidean (diag.) & 62 & $1.90\times10^{-13}$ & $1.19\times10^{-13}$ & $1.06\times10^{-9}$ / $6.13\times10^{-9}$ & $1.54\times10^{-8}$ & $3.19\times10^{-11}$ & 30/0\\
\bottomrule
\end{tabular}}
\par\smallskip\parbox{0.95\linewidth}{\footnotesize Full-rank recovery against a dense direct solve of the same free CN system: 46/46 within the residual-derived allowance (30 validation queries, 8 exact vertices, 8 selected legacy failures; max error/allowance 0.011, max relative coefficient error $5.1\times10^{-14}$).}
\end{table}


On the fixed 30-query validation set (Table~\ref{tab:repaired-pod}) errors
fall by roughly two orders of magnitude per 16 ranks. At ranks 48 and 62
all 30 candidates are accepted directly at $10^{-3}$; at rank 32 or less,
$\beta_0$ dominates $\eta_{\rm aug}$ and the mass-weighted candidates need
correction. Mass weighting minimises the aggregate training-snapshot
projection error in the $M_I$ norm by construction; on the validation maxima
neither basis dominates, and the Euclidean basis accepts 4 of 30 directly at
rank 32. At full rank the model recovers a dense direct solve of the same
free CN system on all 46 tested queries (validation, exact vertices and
selected legacy failures) within an allowance derived from the reference
residual. These are properties of a correctly constrained reduced model at
$N=64$; the basis is mesh-specific and is not evaluated at other
resolutions.

% -----------------------------------------------------------------
\subsection{Solver and certificate costs under refinement}
\label{ssec:scaling-results}

Table~\ref{tab:scaling} and Figure~\ref{fig:scaling} measure the
deterministic solve and certification costs at $N=64,128,256$ on
\texttt{Q\_scale} and at $N=512$ on its predeclared six-query subset, cold
mode, 60 CN steps. The neural and POD models are $N=64$ artefacts with no
transfer strategy and are not run at other resolutions, so this study
measures solver and certificate costs, not neural or POD scaling. Total
time grows roughly linearly in the number of degrees of freedom up to
$N=256$; operator assembly and certification dominate at every size, and
the time-stepping solve is at most a few percent of the total for the free
direct and GMRES solvers. At $N=512$ certification grows faster (5.0 s
against 1.0 s at $N=256$): the dense spectral gate costs more and every
time step reaches the 200-iteration PCG cap. The upper bracket remains valid
at the cap and still meets the target on all 18 runs, but sharper or cheaper
Riesz evaluation would be needed on finer meshes.

% GENERATED by scripts/v6_build_manuscript_assets.py from raw CSVs -- do not edit by hand
\begin{table}[ht]
\centering\small
\caption{Solver and certificate cost under spatial refinement (cold mode, fixed 60 CN steps, single thread; \texttt{Q\_scale}: 30 queries at $N\le256$, a predeclared 6 at $N=512$; 3 timed repetitions). Neural and POD models are $N=64$ artefacts with no transfer strategy and are not run at $N\ne64$: this measures deterministic solve and certification costs, not neural or POD scaling. Times in seconds (median over queries of repetition medians); memory is the median peak traced allocation.}
\label{tab:scaling}
\resizebox{\textwidth}{!}{%
\begin{tabular}{rrlrrrrrrrr}
\toprule
$N$ & dof & method & verified & total & assembly & solve & gate & certify & PCG it. & MB\\
\midrule
64 & 66 & Trusted reference solve & 90/90 & $1.47$ & $0.56$ & $0.649$ & $0.005$ & $0.25$ & 2948 & $0.4$\\
64 & 66 & Free dense LU & 90/90 & $0.83$ & $0.56$ & $0.004$ & $0.005$ & $0.25$ & 2942 & $0.3$\\
64 & 66 & Free cold GMRES & 90/90 & $0.89$ & $0.56$ & $0.069$ & $0.005$ & $0.25$ & 2942 & $0.3$\\
128 & 130 & Trusted reference solve & 90/90 & $2.86$ & $1.14$ & $1.237$ & $0.013$ & $0.47$ & 4870 & $1.0$\\
128 & 130 & Free dense LU & 90/90 & $1.65$ & $1.14$ & $0.007$ & $0.013$ & $0.47$ & 4864 & $1.0$\\
128 & 130 & Free cold GMRES & 90/90 & $1.70$ & $1.14$ & $0.063$ & $0.013$ & $0.47$ & 4868 & $1.0$\\
256 & 258 & Trusted reference solve & 90/90 & $5.75$ & $2.27$ & $2.399$ & $0.040$ & $1.02$ & 8189 & $3.3$\\
256 & 258 & Free dense LU & 90/90 & $3.37$ & $2.28$ & $0.015$ & $0.040$ & $1.01$ & 8160 & $3.1$\\
256 & 258 & Free cold GMRES & 90/90 & $3.43$ & $2.28$ & $0.068$ & $0.040$ & $1.01$ & 8172 & $3.1$\\
512 & 514 & Trusted reference solve & 18/18 & $14.76$ & $4.59$ & $4.856$ & $0.284$ & $5.00$ & 12000 & $11.5$\\
512 & 514 & Free dense LU & 18/18 & $10.02$ & $4.58$ & $0.102$ & $0.282$ & $5.01$ & 12000 & $11.3$\\
512 & 514 & Free cold GMRES & 18/18 & $10.08$ & $4.58$ & $0.183$ & $0.281$ & $5.01$ & 12000 & $11.3$\\
\bottomrule
\end{tabular}}
\end{table}


\begin{figure}[ht]
\centering
\includegraphics[width=0.95\linewidth]{figures/v6_scaling.png}
\caption{Solver and certificate cost under spatial refinement (cold mode;
\texttt{Q\_scale}, 30 queries at $N\le256$ and the predeclared six at
$N=512$). Left: median total time per query. Right: phase medians for the
free dense LU solver. No neural or POD model is run at $N\ne64$. Replaces
the historical runtime-scaling figure (Appendix~\ref{app:historical-figures}).}
\label{fig:scaling}
\end{figure}

% -----------------------------------------------------------------
\subsection{Assembly convention, refinement and the corner discrepancy}
\label{ssec:refinement}

Table~\ref{tab:assembly} separates two comparisons for the assembled
operator (Remark~\ref{rem:assembly-convention}). Against an independent
40-digit quadrature of the \emph{same} principal-period integral, with
independently derived mass and derivative matrices, the production
operator differs by $10^{-3}$ in the largest entry at $N=64$, falling to
$10^{-5}$ at $N=512$, and by $0.7\%$ to $0.1\%$ in operator norm; the free
gate decision agrees on all 33 queries and $\lambda_{\min}(S_I)$ moves by at
most $4\times10^{-7}$ relative. This is the agreement of an implementation
with its convention. Against the exact real-line Galerkin coefficients the
operator-norm deviation is $33$--$41\%$, concentrated at high frequencies,
the largest entry differs by about $10\%$, and $\lambda_{\min}(S_I)$ moves by
at most $8.9\times10^{-4}$ relative. The certificates in this paper are
relative to the production discrete problem; the exact-Galerkin deviation
enters only the comparison with continuum prices.

% GENERATED by scripts/v6_build_manuscript_assets.py from raw CSVs -- do not edit by hand
\begin{table}[ht]
\centering\small
\caption{Assembly checks against the production operator (exact Toeplitz action of the principal-period, $L=2N_h$ IFFT coefficients). \emph{Same convention}: exact (40-digit) quadrature of the same principal-period integral, independently re-derived $M_h,D_h$; \emph{real-line Galerkin}: exact Galerkin coefficients (diagnostic). Relative errors; operator norms computed directly. Population: \texttt{Q\_scale} first 10 ($N\le256$), first 3 ($N=512$).}
\label{tab:assembly}
\resizebox{\textwidth}{!}{%
\begin{tabular}{lrrrrrrr}
\toprule
comparator & $N$ & $q$ & max entry $A$ & max entry $M$ & $\|\Delta A\|_2/\|A\|_2$ med./max & $|\Delta\lambda_{\min}(S_I)|$ rel. & free gate agrees\\
\midrule
same convention & 64 & 10 & $1.03\times10^{-3}$ & $4.0\times10^{-16}$ & $6.95\times10^{-3}$ / $0.0135$ & $4.1\times10^{-7}$ & 10/10\\
same convention & 128 & 10 & $2.59\times10^{-4}$ & $4.0\times10^{-16}$ & $3.26\times10^{-3}$ / $6.23\times10^{-3}$ & $6.0\times10^{-8}$ & 10/10\\
same convention & 256 & 10 & $6.34\times10^{-5}$ & $3.0\times10^{-15}$ & $1.55\times10^{-3}$ / $2.92\times10^{-3}$ & $8.3\times10^{-9}$ & 10/10\\
same convention & 512 & 3 & $1.22\times10^{-5}$ & $7.9\times10^{-15}$ & $1.00\times10^{-3}$ / $1.11\times10^{-3}$ & $9.4\times10^{-10}$ & 3/3\\
real-line Galerkin & 64 & 10 & $0.108$ & $4.0\times10^{-16}$ & $0.382$ / $0.412$ & $8.9\times10^{-4}$ & 10/10\\
real-line Galerkin & 128 & 10 & $0.105$ & $4.0\times10^{-16}$ & $0.371$ / $0.408$ & $4.4\times10^{-4}$ & 10/10\\
real-line Galerkin & 256 & 10 & $0.102$ & $3.0\times10^{-15}$ & $0.358$ / $0.405$ & $2.2\times10^{-4}$ & 10/10\\
real-line Galerkin & 512 & 3 & $0.101$ & $7.9\times10^{-15}$ & $0.333$ / $0.402$ & $1.1\times10^{-4}$ & 3/3\\
\bottomrule
\end{tabular}}
\end{table}


The refinement study (Table~\ref{tab:refinement}, first eight
\texttt{Q\_ref30} queries) measures how the production discrete price
approaches the continuum price under the production convention. Paired
differences between successive spatial levels decrease at a rate close to
second order, the temporal differences at fixed $N=1024$ decrease at second
order down to about $2\times10^{-6}$, and the COS gap flattens at about
$5\times10^{-5}$ once both are refined; at $N=1024$ the COS gap is
$5.0\times10^{-5}$. The domain study at fixed $h=0.125$ still shows a COS
gap that changes by $2.5\times10^{-3}$ between half-widths 6 and 12, so
domain sensitivity at that resolution is not resolved. These observations
are empirical convergence evidence for the production scheme towards the
independent price; they do not provide a computable bound for the
discretisation or truncation error.

% GENERATED by scripts/v6_build_manuscript_assets.py from raw CSVs -- do not edit by hand
\begin{table}[ht]
\centering\small
\caption{Refinement of the production discrete price (first 8 \texttt{Q\_ref30} queries; spots $0.8K..1.2K$, maximum within query, then median over queries). Paired difference $=\max_s|V_k-V_{k+1}|$ to the next level; COS gap is an independent continuum diagnostic. Space: $h=16/N$, 480 steps; time: $N=1024$; domain: $h=0.125$ fixed, half-width $L$.}
\label{tab:refinement}
\begin{tabular}{lrrrrr}
\toprule
study & $N$ & steps & $L$ & paired diff. (\$) & COS gap (\$)\\
\midrule
space & 64 & 480 & $8$ & $0.0233$ & $0.0215$\\
space & 128 & 480 & $8$ & $5.39\times10^{-3}$ & $5.37\times10^{-3}$\\
space & 256 & 480 & $8$ & $1.08\times10^{-3}$ & $1.15\times10^{-3}$\\
space & 512 & 480 & $8$ & $2.40\times10^{-4}$ & $1.94\times10^{-4}$\\
space & 1024 & 480 & $8$ & -- & $5.05\times10^{-5}$\\
time & 1024 & 15 & $8$ & $0.0411$ & $0.0413$\\
time & 1024 & 30 & $8$ & $5.71\times10^{-4}$ & $7.65\times10^{-4}$\\
time & 1024 & 60 & $8$ & $1.43\times10^{-4}$ & $2.01\times10^{-4}$\\
time & 1024 & 120 & $8$ & $3.59\times10^{-5}$ & $7.35\times10^{-5}$\\
time & 1024 & 240 & $8$ & $8.96\times10^{-6}$ & $4.84\times10^{-5}$\\
time & 1024 & 480 & $8$ & $2.13\times10^{-6}$ & $5.05\times10^{-5}$\\
time & 1024 & 960 & $8$ & -- & $5.04\times10^{-5}$\\
domain & 96 & 240 & $6$ & $1.87\times10^{-3}$ & $6.71\times10^{-3}$\\
domain & 128 & 240 & $8$ & $9.05\times10^{-4}$ & $5.37\times10^{-3}$\\
domain & 160 & 240 & $10$ & $6.01\times10^{-4}$ & $4.67\times10^{-3}$\\
domain & 192 & 240 & $12$ & -- & $4.15\times10^{-3}$\\
\bottomrule
\end{tabular}
\end{table}


At the two box corners with $C=0.05$, $Y=1.1$, $T=0.1$, $r=0.01$ the $N=64$
discrete price differs from COS by more than \$2 and is negative at some
spots (Table~\ref{tab:corner}). The gap falls to $0.77$, $0.20$, $0.028$ and
$0.0016$ under spatial refinement, is unchanged by time refinement, and the
$\tau=0$ projection error of the payoff decreases at first order in $h$. The
discrepancy is consistent with coarse spatial resolution and payoff
projection for these short-maturity, weakly diffusive parameters; the
propagated contribution of the initial projection alone was not isolated.
The algebraic certificate of these queries remains valid relative to the
discrete problem; the discrete problem itself is a poor continuum
approximation there.

% GENERATED by scripts/v6_build_manuscript_assets.py from raw CSVs -- do not edit by hand
\begin{table}[ht]
\centering\small
\caption{Corner discrepancy diagnostic (240 time steps; the $N=64$ values are unchanged with 60 steps). Q\_audit\_0401/0397: $C=0.05$, $Y=1.1$, $T=0.1$, $r=0.01$. Spots $0.8K..1.2K$.}
\label{tab:corner}
\begin{tabular}{lrrrr}
\toprule
query & $N$ & $\max_s|V_h(T)-\mathrm{COS}|$ & $\max_s|V_h(0)-\psi|$ & $\min_s V_h(T)$\\
\midrule
Q\_audit\_0401 & 64 & $2.22$ & $3.30$ & $-1.17$\\
Q\_audit\_0401 & 128 & $0.774$ & $1.66$ & $-0.357$\\
Q\_audit\_0401 & 256 & $0.195$ & $0.829$ & $0.0225$\\
Q\_audit\_0401 & 512 & $0.0275$ & $0.415$ & $0.0489$\\
Q\_audit\_0401 & 1024 & $1.61\times10^{-3}$ & $0.207$ & $0.0474$\\
Q\_audit\_0397 & 64 & $2.07$ & $3.30$ & $-1.11$\\
Q\_audit\_0397 & 128 & $0.693$ & $1.66$ & $-0.293$\\
Q\_audit\_0397 & 256 & $0.164$ & $0.829$ & $0.0379$\\
Q\_audit\_0397 & 512 & $0.0200$ & $0.415$ & $0.0517$\\
Q\_audit\_0397 & 1024 & $8.34\times10^{-4}$ & $0.207$ & $0.0510$\\
Q\_bench120\_0000 & 64 & $0.107$ & $3.30$ & $3.48$\\
Q\_bench120\_0000 & 128 & $8.46\times10^{-3}$ & $1.66$ & $3.58$\\
Q\_bench120\_0000 & 256 & $6.71\times10^{-4}$ & $0.829$ & $3.58$\\
Q\_bench120\_0000 & 512 & $1.65\times10^{-4}$ & $0.415$ & $3.58$\\
Q\_bench120\_0000 & 1024 & $2.82\times10^{-5}$ & $0.207$ & $3.58$\\
\bottomrule
\end{tabular}
\end{table}


% -----------------------------------------------------------------
\subsection{Independent Fourier prices and financial budgets}
\label{ssec:fourier-financial}

Independent continuum prices are diagnostics in this paper. The earlier
cross-check reported a maximum relative COS--Carr--Madan discrepancy of
$2.34\times10^{-8}$ on 27 queries, with its worst query (\texttt{box\_12})
reaching the Carr--Madan cutoff, while a later check on the first 20 \texttt{Q\_ref30} queries, a
disjoint population, reported $9.06\times10^{-10}$ with its own settings. We reran both populations, on the same
coordinates and spots, at the historical settings and at extended settings
(Table~\ref{tab:fourier}). The \texttt{box\_12} discrepancy is a COS-setting
effect: the historical COS value moves by $2.34\times10^{-8}$ when the
expansion length and range are increased, the Carr--Madan value is identical
at the historical and extended cutoffs, and the extended COS and
Carr--Madan prices agree to $2.1\times10^{-11}$; a 60-digit Carr--Madan evaluation
reproduces the 40-digit value. On all 150 \texttt{Q\_ref30} rows the two
methods agree to $8\times10^{-13}$ at the historical settings and
$4\times10^{-13}$ at the extended ones. Agreement between methods and settings is convergence evidence,
not an error enclosure.

% GENERATED by scripts/v6_build_manuscript_assets.py from raw CSVs -- do not edit by hand
\begin{table}[ht]
\centering\small
\caption{Same-query Fourier reruns (continuum-price diagnostics; relative differences scaled by $\max(|V|,1)$). Historical settings: COS $N\le1024$, $L\le20$; Carr--Madan cutoff cap 400. Extended: COS $N=8192$, $L=30$; cap 51{,}200. 60-digit Carr--Madan on the worst cases.}
\label{tab:fourier}
\resizebox{\textwidth}{!}{%
\begin{tabular}{lrrrrrrr}
\toprule
population & rows & max COS--CM hist. & COS hist. vs ext. & CM hist. vs ext. & COS--CM ext. & CM cap hits h./e. & 60-digit vs COS ext.\\
\midrule
historical 27 ($S_0=K$) & 27 & $2.34\times10^{-8}$ (box\_12) & $2.34\times10^{-8}$ & $0$ & $2.14\times10^{-11}$ & 1/0 & $2.14\times10^{-11}$ ($n=3$)\\
Q\_ref30 ($30\times5$ spots) & 150 & $7.87\times10^{-13}$ (Q\_ref30\_0021) & $4.17\times10^{-13}$ & $0$ & $3.71\times10^{-13}$ & 0/0 & $3.71\times10^{-13}$ ($n=3$)\\
\bottomrule
\end{tabular}}
\end{table}


The financial study (Table~\ref{tab:financial}) converts the trajectory
bound to spot prices through the exact point factor $\beta_I(x)$ at the
five spots of each \texttt{Q\_ref30} query. Raw DeepONet predictions meet no
budget from \$0.01 to \$0.10, with bound median \$7.4; after correction all
30 queries are verified at \$0.01 with bound median $\$2.2\times10^{-6}$.
Rank-48 POD is verified directly on all 30 queries. Both bounds are
relative to the exact production discrete price. The difference between the
discrete price and COS, median \$0.030 and maximum \$0.39 over queries, is
four orders of magnitude larger; it is the discretisation diagnostic of
Section~\ref{ssec:refinement}, not a certified quantity. A sampled maximum
of $\beta_I$ over the spot window and a cellwise floating-point maximum
agree; neither is a validated supremum.

% GENERATED by scripts/v6_build_manuscript_assets.py from raw CSVs -- do not edit by hand
\begin{table}[ht]
\centering\small
\caption{Financial budgets on all 30 \texttt{Q\_ref30} queries at five spots (maximum within query; medians over queries). Bounds are $\beta_I(x)\eta_{\rm aug}$ relative to the exact production discrete price (exact-arithmetic theorem, floating-point evaluation); the stored-reference and COS columns are diagnostics, the latter including discretisation error.}
\label{tab:financial}
\resizebox{\textwidth}{!}{%
\begin{tabular}{lrrrrrrrr}
\toprule
model & delivered & raw verif. 1/5/10\textcent & final verif. 1/5/10\textcent & raw bound (\$) & final bound med./max (\$) & ref. allowance (\$) & diff. stored ref. (\$) & diff. COS med./max (\$)\\
\midrule
DeepONet & 30/30 & 0 / 0 / 0 & 30 / 30 / 30 & $7.36$ & $2.23\times10^{-6}$ / $7.99\times10^{-6}$ & $2.17\times10^{-6}$ & $6.33\times10^{-8}$ & $0.0297$ / $0.391$\\
POD $r=48$ & 30/30 & 30 / 30 / 30 & 30 / 30 / 30 & $3.34\times10^{-5}$ & $3.34\times10^{-5}$ / $7.03\times10^{-5}$ & $2.17\times10^{-6}$ & $6.31\times10^{-8}$ & $0.0297$ / $0.391$\\
\bottomrule
\end{tabular}}
\par\smallskip\parbox{0.95\linewidth}{\footnotesize Window factors on $[\log0.8K,\log1.2K]$ (two mesh cells): sampled maximum of $\beta_I$ 2.199824 (diagnostic); cellwise piecewise-polynomial maximum 2.19982442 (floating-point evaluation, not a validated supremum); window-Gram $L^2$ factor 0.959 (integrated error, kept separate).}
\end{table}


% =================================================================
\subsection{Historical neural-model experiments}
\label{ssec:historical}

The studies below precede the integrated workflow. They used the legacy
full-space certificate, which requires the full matrix $\bar S$ to pass the
positivity gate (Section~\ref{ssec:audit-results}), and they were run on
their own populations with their own reference and setup rules. Their
accuracy, effectivity and stress observations are retained as historical
evidence about the learned models. Their speed ratios are withdrawn and
superseded by Section~\ref{ssec:benchmark}, and their POD results are
superseded by Section~\ref{ssec:pod-results}.

\subsubsection{Interpolation accuracy and seed stability}
\label{ssec:interpolation-results}

Table~\ref{tab:interpolation} compares the three neural models on $100$
unseen in-distribution queries.  DeepONet gives the lowest median and mean
relative $L^2$ and energy errors.  The historical POD row of this study is
omitted: that model evolved incorrect reduced equations, and the repaired
model is evaluated in Section~\ref{ssec:pod-results} on its own population.
The projected FNO is evaluated at changed resolutions in a separate study.

\begin{table}[ht]
\centering
\caption{Historical interpolation statistics on $100$ unseen queries (legacy POD row omitted, see text). Relative $L^2$ is the maximum-in-time mass-norm error divided by the
maximum-in-time reference mass norm. Price error is terminal ATM error
against the stored Galerkin reference. The full-space energy diagnostic
requires separate gate eligibility and is not certified pointwise accuracy.}
\label{tab:interpolation}
\small
\resizebox{\textwidth}{!}{%
\begin{tabular}{lrrrrr}
\toprule
Model & median rel. $L^2$ & mean rel. $L^2$ & median rel. energy & mean rel. energy & median price err.\\
\midrule
MLP & $3.80\!\times\!10^{-3}$ & $5.42\!\times\!10^{-3}$ & $2.71\!\times\!10^{-2}$ & $3.57\!\times\!10^{-2}$ & $2.84\!\times\!10^{-1}$\\
DeepONet & $2.02\!\times\!10^{-3}$ & $2.52\!\times\!10^{-3}$ & $1.22\!\times\!10^{-2}$ & $1.67\!\times\!10^{-2}$ & $1.64\!\times\!10^{-1}$\\
Projected FNO & $5.52\!\times\!10^{-3}$ & $6.53\!\times\!10^{-3}$ & $3.10\!\times\!10^{-2}$ & $3.63\!\times\!10^{-2}$ & $6.24\!\times\!10^{-1}$\\
\bottomrule
\end{tabular}}
\end{table}

\begin{figure}[ht]
\centering
\includegraphics[width=0.98\linewidth]{exp01_interpolation.png}
\caption{Historical interpolation errors on the common reported
discretisation. The POD curves show the legacy model with incorrect reduced
equations, superseded by Section~\ref{ssec:pod-results}; they are retained
only to preserve the original figure and support no comparison.}
\label{fig:interpolation}
\end{figure}

Retraining the main DeepONet with three independent seeds gives mean relative
$L^2$ error $(2.399\pm0.020)\times10^{-3}$ and mean relative energy error
$(1.634\pm0.034)\times10^{-2}$, where the uncertainty is the sample standard
deviation across seeds.  The mean effectivity is $5.68\pm0.16$.  This is limited evidence of initialisation stability on that test. It does
not replace parameter-sampling uncertainty or full training reproducibility.

\subsubsection{Certificate reliability and iterative Riesz solves}
\label{ssec:certificate-results}

For each certifiable query define
$E_i=\max_m\norm{U_h^m-U_{\theta,h}^m}_{L^2(\Omega)}$ and
$I_{\rm eff}=\eta/E_i$.  Three of the $100$ sampled queries fail the reported stability-bound
gate. The available records do not distinguish actual indefiniteness from
an inconclusive positive lower bound.  On the remaining $97$ queries, both the direct sharp estimator and the fully converged upper bracket
have zero reported violations. This is conditional algebraic evidence;
validated floating-point certification and reference-solve uncertainty are
not established by these counts.  Their effectivity statistics coincide to the reported precision:
minimum $2.46$, median $5.65$, mean $5.75$, $95$th percentile $8.43$, and
maximum $11.47$.  The mass-based diagnostic is substantially more
conservative and, because it is not used as a proven online bound in the
implementation, no violation count is attached to it.

\begin{table}[ht]
\centering
\caption{Effectivity statistics on the $97$ certifiable queries from the
$100$-query certificate test.}
\label{tab:effectivity}
\small
\begin{tabular}{lrrrrrr}
\toprule
Quantity & min & median & mean & p95 & max & violations\\
\midrule
Mass indicator (heuristic) & $8.00$ & $38.27$ & $40.98$ & $70.65$ & $99.46$ & --\\
$\etaSharp$ & $2.46$ & $5.65$ & $5.75$ & $8.43$ & $11.47$ & $0$\\
$\etaGuar$ & $2.46$ & $5.65$ & $5.75$ & $8.43$ & $11.47$ & $0$\\
\bottomrule
\end{tabular}
\end{table}

The same $97$ queries provide historical graph-norm diagnostics.
Here the test uses the full-space version of the norm
\eqref{eq:triple-norm}, including both the time-derivative dual term and the
midpoint energy term, and computes $\kappa_h$ from \eqref{eq:kappa-def} rather
than from $\operatorname{cond}(\Sh)$.  There are zero reported numerical inequality violations. Because the
residual masks prescribed rows while the graph norm uses the full matrices,
these records do not directly validate the constrained efficiency theorem.  The measured $\kappa_h$ ranges from $1.006$ to $2.636$ (median
$1.530$), while
\[
  \frac{\etaSharp^2}{|||\ee|||^2}
\]
ranges from $0.979$ to $1.196$ (median $1.018$).  For comparison, the
right-hand factor $2\max\{1,\kappa_h^2\}$ ranges from $2.022$ to $13.897$;
the largest observed fraction of the available bound is $0.506$.  The reported ratio is close to one on this set. The algebraic identity
$\rho=-M_h\partial e-A_h\overline e$ explains its relation to the graph
norm, but does not force a near-unit ratio for every nonsymmetric query.
Non-violation is a useful implementation check, not an independent
demonstration of estimator novelty.

\begin{table}[ht]
\centering
\caption{Reported consistency check of the discrete efficiency estimate on the $97$
gate-passing reported queries.}
\label{tab:efficiency-validation}
\small
\begin{tabular}{lrrr}
\toprule
Quantity & min & median & max\\
\midrule
$\kappa_h$ & $1.006$ & $1.530$ & $2.636$\\
$\etaSharp^2/|||\ee|||^2$ & $0.979$ & $1.018$ & $1.196$\\
$2\max\{1,\kappa_h^2\}$ & $2.022$ & $4.681$ & $13.897$\\
\bottomrule
\end{tabular}
\end{table}

\begin{figure}[ht]
\centering
\includegraphics[width=0.72\linewidth]{exp02_effectivity.png}
\caption{Reported assembled-reference error versus the sharp and upper-bracket
quantities; their numerical validation assumptions are discussed in the text.
The dashed line is $\eta=E$.  The mass quantity is shown only as a heuristic
diagnostic.}
\label{fig:error-estimator}
\end{figure}

The inexact-solve experiment confirms the bracket mechanism directly.  With a
PCG cap of $1,2,5,10,20,50,100,200$ iterations, the mean guaranteed
effectivity decreases from $14.60$ to $6.04$ and reaches the sharp value by
about $50$ iterations at this resolution.  Hence computational effort changes
sharpness but not the direction of the guarantee.  At larger meshes a fixed
$200$-iteration cap leaves a relative bracket gap of about
$5.4\times10^{-6}$ at $N=256$ and $6.7\times10^{-3}$ at $N=512$, which is a
sharpness limitation under the bracket hypotheses. The separate cost
figure shows a much smaller gap near $10^{-11}$ at $N=512$; the right-hand
sides, arithmetic settings and aggregation needed to reconcile these two
reported gaps are not supplied. Neither value is used here to infer a
universal convergence rate.

\begin{figure}[ht]
\centering
\includegraphics[width=0.92\linewidth]{exp03_certificate_levels.png}
\caption{Guaranteed-certificate convergence as the PCG iteration budget is
increased.  Additional iterations reduce the bracket width and drive
$\etaGuar$ toward $\etaSharp$; soundness does not depend on convergence.}
\label{fig:certificate-levels}
\end{figure}

\subsubsection{Stress and out-of-distribution behaviour}
\label{ssec:ood-results}

The certificate should increase when the learned approximation deteriorates.
The reported four-tier stress test has this trend in its mean statistics.  In the latter, the mean true error rises from $4.05$ in the interior to
$5.04$ near the boundary, $13.42$ at the edge and $28.18$ beyond the nominal
domain; the mean guaranteed certificate rises in the same order from $13.57$
to $16.92$, $43.62$ and $86.81$.  The corresponding certification-gate
failure counts are $0/50$, $0/50$, $3/50$ and $9/50$.  These are aggregate observations. The figure refers to a shrunken training
region (fraction $0.35$), whereas the main dataset is described on the full
box; the stress checkpoint and tier ranges must be identified before treating
these as comparable in-distribution/OOD statistics. Means of true error and
certificate should use the same gate-passing subset, with all failures
reported separately.

\begin{figure}[ht]
\centering
\includegraphics[width=0.78\linewidth]{exp05_ood.png}
\caption{True error and guaranteed certificate across increasingly difficult
out-of-distribution tiers.}
\label{fig:ood}
\end{figure}

A separate decision-table experiment evaluates acceptance against the true
error over a sweep of thresholds and records zero false accepts for the
guaranteed certificate.  False rejects are expected because the estimator is
conservative; they represent extra computation rather than an incorrect
accepted solution.

\subsubsection{Learning curve and residual-loss ablation}
\label{ssec:learning-ablation}

The DeepONet learning curve uses $N_{\rm train}\in\{100,200,400,800\}$ and a
fixed held-out test set.  The reported test relative $L^2$ error decreases from
$4.12\times10^{-3}$ to $2.70\times10^{-3}$, while the mean guaranteed
certificate decreases from $4.24$ to $3.49$.  A log--log fit gives the rough
empirical exponent $\beta=0.189$; with only four points we treat this as a
descriptive slope, not a scaling law.

\begin{figure}[ht]
\centering
\includegraphics[width=0.92\linewidth]{exp13_learning_curve.png}
\caption{Training-data scaling for the main DeepONet.}
\label{fig:learning-curve}
\end{figure}

Residual regularisation does not provide a robust monotone improvement.  In a
controlled MLP ablation, $\lambda_R=10^{-9}$ gives a small improvement over
data-only training (prediction error $1.31\times10^{-2}$ versus
$1.36\times10^{-2}$ and certificate $14.10$ versus $14.95$), but every tested
$\lambda_R\ge10^{-8}$ eventually degrades the prediction and the certificate;
at $10^{-5}$ the prediction error is $4.15\times10^{-2}$ and the certificate
is $33.14$.  We therefore regard residual regularisation as a sensitive
hyperparameter rather than a mechanism that automatically improves
certification.

\subsubsection{Resolution transfer}
\label{ssec:resolution-results}

The coefficient MLP, coefficient DeepONet and POD basis are structurally tied
to the training mesh and cannot natively accept a changed coefficient
dimension.  This is an architectural incompatibility, not an empirical
``failure to generalise.''  Projected FNO is discretisation-flexible and can be
evaluated natively at $N=32,64,128,256$.  Its median relative $L^2$ errors are
$1.45\times10^{-2}$, $6.00\times10^{-3}$,
$1.06\times10^{-2}$ and $1.99\times10^{-2}$, respectively, with median
guaranteed certificates $19.13$, $8.84$, $12.02$ and $17.79$.  At $N=32$,
$12/50$ queries fail the certification gate; the finer three resolutions in
this experiment do not show that issue.

\begin{figure}[ht]
\centering
\includegraphics[width=0.98\linewidth]{exp10_cross_resolution.png}
\caption{Cross-resolution study.  The projected FNO is evaluated natively at each resolution. The right
panel additionally compares projected DeepONet terminal states with FNO
trajectories; that comparison is descriptive and needs a common temporal
metric. Native evaluation is not equivalent to mesh-independent accuracy.}
\label{fig:cross-resolution}
\end{figure}

\subsubsection{Varying-strike operator generalisation}
\label{ssec:operator-generalization}

The main reported dataset fixes $K=100$, so by itself it tests a
parameter-to-solution map.  To examine a restricted payoff family, the reported study generated a
separate $400$-trajectory dataset with strike sampled jointly over
$K\in[70,150]$, excluding the held-out interval $[95,105]$, and train a new
DeepONet with the same architecture.  Table~\ref{tab:operator-generalization}
compares this varying-strike model with the original $K=100$ parameter map on
three test tiers.

\begin{table}[ht]
\centering
\caption{Operator-generalisation experiment.  Entries are means over
successfully certified queries; the final column records certification-gate
failures.}
\label{tab:operator-generalization}
\small
\resizebox{\textwidth}{!}{%
\begin{tabular}{llrrrrr}
\toprule
Tier & Model & rel. $L^2$ & rel. energy & $\etaGuar$ & effectivity & gate failures\\
\midrule
unseen $\mu$ & varying-$K$ operator & $2.82\!\times10^{-2}$ & $6.80\!\times10^{-2}$ & $29.31$ & $3.97$ & $2/30$\\
 & $K=100$ map & $3.13\!\times10^{-2}$ & $1.25\!\times10^{-1}$ & $27.22$ & $3.05$ & $2/30$\\
unseen strike & varying-$K$ operator & $1.96\!\times10^{-2}$ & $6.02\!\times10^{-2}$ & $17.48$ & $3.87$ & $0/30$\\
 & $K=100$ map & $4.55\!\times10^{-3}$ & $2.63\!\times10^{-2}$ & $4.55$ & $4.41$ & $0/30$\\
unseen combination & varying-$K$ operator & $4.32\!\times10^{-2}$ & $8.84\!\times10^{-2}$ & $50.18$ & $4.14$ & $1/30$\\
 & $K=100$ map & $3.81\!\times10^{-2}$ & $1.47\!\times10^{-1}$ & $36.03$ & $3.25$ & $1/30$\\
\bottomrule
\end{tabular}}
\end{table}

The table reports a restricted payoff-family experiment, not generalisation
to arbitrary functions. On the full line the multiplicative CGMY model gives
\begin{equation}
 P(S,K;\mu,T)=\frac{K}{K_0}
             P\!\left(\frac{K_0}{K}S,K_0;\mu,T\right),
 \label{eq:strike-homogeneity}
\end{equation}
since $(K-SZ_T)^+=(K/K_0)(K_0-(K_0/K)SZ_T)^+$.
Equivalently, log-price surfaces shift by $\log(K_0/K)$ and scale by
$K/K_0$. Truncation and projection modify this equivalence, so a shifted,
scaled fixed-strike baseline should include its domain and projection errors.
The varying-strike payoff is a functional input on a one-dimensional manifold;
it is not a demanding test of general payoff-shape learning.

The original fixed-strike model has no trained strike variation. Its use on
other strikes (including any altered payoff encoding, projection or shift)
must be documented. Its reported mean spatial error $3.13\times10^{-2}$ in
this table also differs substantially from the main interpolation mean
$2.52\times10^{-3}$; the test distributions and evaluation rules are not
specified enough to infer the cause. Training budgets differ ($400$ versus
$800$ trajectories), so the table does not isolate a breadth--specialisation
trade-off. A matched-budget study and the homogeneity baseline are needed
before claiming an operator-learning advantage.

\subsubsection{Predict--certify--correct tolerance sweep}
\label{ssec:tolerance-results}

The central end-to-end experiment runs $120$ common queries at eight
acceptance tolerances, for $960$ workflow executions in total.  One query per
tolerance fails the certification gate, so the final certified workflow
success rate is $99.2\%$ throughout.  For $\varepsilon\le2$, direct
acceptance is zero and correction is required for $99.2\%$ of queries.
At $\varepsilon=5$, $53.3\%$ of queries are accepted directly, and at
$\varepsilon=10$ direct acceptance reaches $87.5\%$. The speed ratios
originally reported for this sweep compared against an unoptimised
reference with different setup accounting and are withdrawn; at the
production tolerance $10^{-3}$ the matched benchmark of
Section~\ref{ssec:benchmark} supersedes them.  The $960$ executions reuse $120$ parameter queries, rather than represent
$960$ independent tests. Eight executions fail the gate, leaving $952$
gate-passing executions with zero reported violations. The corrected output
must be rechecked against its requested tolerance, as in
Figure~\ref{alg:pcc}; no tolerance guarantee follows from the correction
budget alone.

\begin{table}[ht]
\centering
\caption{Historical absolute-tolerance sweep (legacy certificate, 120
queries). The original runtime and speedup columns are withdrawn
(Section~\ref{ssec:benchmark}).}
\label{tab:tolerance}
\small
\begin{tabular}{rrrrr}
\toprule
$\varepsilon$ & direct accept & correction & final certified & max true error\\
\midrule
$10^{-3}$ & $0.0\%$ & $99.2\%$ & $99.2\%$ & $2.32\times10^{-6}$\\
$10^{-2}$ & $0.0\%$ & $99.2\%$ & $99.2\%$ & $2.32\times10^{-6}$\\
$10^{-1}$ & $0.0\%$ & $99.2\%$ & $99.2\%$ & $2.32\times10^{-6}$\\
$0.5$ & $0.0\%$ & $99.2\%$ & $99.2\%$ & $2.32\times10^{-6}$\\
$1$ & $0.0\%$ & $99.2\%$ & $99.2\%$ & $2.32\times10^{-6}$\\
$2$ & $0.0\%$ & $99.2\%$ & $99.2\%$ & $2.32\times10^{-6}$\\
$5$ & $53.3\%$ & $45.8\%$ & $99.2\%$ & $9.04\times10^{-1}$\\
$10$ & $87.5\%$ & $11.7\%$ & $99.2\%$ & $1.67$\\
\bottomrule
\end{tabular}
\end{table}

The original tolerance-sweep figure is retained, with its superseded speed
panel marked, as Figure~\ref{fig:tolerance-sweep} in
Appendix~\ref{app:historical-figures}.

\paragraph{Financial interpretation of the tolerance.}
The tolerance in this experiment is an absolute global log-space $L^2$ budget,
not a pointwise price budget. With $|\Omega|=16$, an error norm $1.67$
corresponds to spatial RMS $1.67/4\approx0.418$, but says nothing by itself
about error at a specific spot. No raw predictions pass for
$\varepsilon\le2$ in this sample. The withdrawn $5.43\times$ mean ratio
occurred at a loose absolute global tolerance, not at a demonstrated
financially relevant relative tolerance. The discrete evaluation bound in
Section~\ref{ssec:greeks-cert} can convert a mass-norm budget to a spotwise
price bound without changing the stability proof. Relative budgets need an
explicit online-computable normalisation; near-zero prices require an
absolute floor. The v5 financial study of Section~\ref{ssec:fourier-financial} uses
exactly this conversion at fixed spots.

\subsubsection{Historical runtime scaling and certificate cost}
\label{ssec:runtime-scaling}

The historical scaling experiment timed the projected FNO with forced
correction against the unoptimised trusted reference solve at
$N=32,\dots,512$ on a common set of 30 queries certifiable at every
resolution. Its reference/workflow ratios ($1.9$--$8.3\times$) compared a
workflow with shared operator caches against a reference that built its own
setup, at different accuracy criteria, and are withdrawn; the original
figure is retained as Figure~\ref{fig:runtime-scaling} in
Appendix~\ref{app:historical-figures}. Section~\ref{ssec:scaling-results}
replaces it with a cold-mode measurement of the deterministic solve and
certificate costs on a declared population.

A separate apples-to-apples certification-cost benchmark charges the
stability-constant setup to both direct and PCG Riesz evaluation.  Under this
comparison PCG is not faster at any tested resolution: total direct versus PCG
times are $2.53/5.02$ ms at $N=32$, $6.16/10.96$ ms at $N=64$,
$21.63/29.12$ ms at $N=128$, $182.48/197.49$ ms at $N=256$, and
$344.18/421.47$ ms at $N=512$.  The reason is that the dominant cost at large
$N$ is shared setup rather than the incremental direct Riesz solve.  At
$N=512$, dense symmetric-operator assembly alone costs about $191.6$ ms,
whereas the direct factorisation plus solve costs only about $0.60$ ms and the
PCG solve about $77.9$ ms.

\begin{figure}[ht]
\centering
\includegraphics[width=0.78\linewidth]{figures/cert_cost_benchmark.png}
\caption{Historical certification-cost benchmark (legacy full-space
certificate) with the same stability-constant setup charged to both direct
and PCG paths.  No PCG crossover is reported on the tested resolutions. The plotted
near-$10^{-11}$ bracket gap needs reconciliation with the larger gaps
reported in the iterative trajectory study.}
\label{fig:cert-cost}
\end{figure}

The preconditioner study gives complementary information.  At the reported
resolution, no preconditioner requires $51.5$ PCG iterations and about
$4.08$ ms per amortised Riesz solve; the mass preconditioner increases this
to $79.9$ iterations and $6.56$ ms.  A banded approximation reduces the mean
iteration count to $6.0$ and the preconditioned condition number to about
$1.02$, but a roughly $0.53$ s setup cost leaves its amortised time at
$9.50$ ms for a $60$-step trajectory.  The reported $4.08$ ms per Riesz solve would correspond to about $245$ ms
for $60$ solves, whereas the separate cost benchmark reports approximately
$11$ ms total PCG cost at $N=64$. Its number of right-hand sides and
amortisation convention must be specified to reconcile these workloads.
The comparison shows that setup and iteration counts must be measured
together; it does not establish the fastest certification configuration.

These historical cost measurements used the legacy full-space certificate
and their own query sets; the v5 matched benchmark
(Section~\ref{ssec:benchmark}) and resolution study
(Section~\ref{ssec:scaling-results}) supersede them for all speed and cost
claims.

\section{Discussion}
\label{sec:discussion}

\subsection{Three levels of claim}
The paper makes claims at three levels that must not be conflated.
(i) \emph{Production discrete certification.} Every certificate bounds the
error relative to the exact solution of the assembled constrained discrete
problem: the principal-period operator of
Remark~\ref{rem:assembly-convention}, the stored forcing and the
prescribed lifting. These are exact-arithmetic theorems evaluated in
ordinary floating point; the v5 evidence (no detected bracket violations,
augmented bounds holding in every tested case) is finite-sample validation
of their implementation, not a machine-validated enclosure. (ii)
\emph{Exact Galerkin.} The assembled operator is not the exact Galerkin
matrix written in Section~\ref{sec:galerkin}; it deviates by 33--41\% in
operator norm at high frequencies (Table~\ref{tab:assembly}). Statements
proved for the exact Galerkin matrix (Lemmas~\ref{lem:Sh-spd},
\ref{lem:mass-lower} and~\ref{lem:kappa}) are not statements about the
production matrix, whose positivity the online gate checks directly.
(iii) \emph{Continuum prices.} Independent COS and Carr--Madan prices are
diagnostics. The production discrete price differs from them by a median
\$0.03 at $N=64$ and by more than \$2 at two box corners; refinement
reduces the gap at close to second order, but no computable continuum
error bound is supplied (Corollary~\ref{cor:total-error} is not
evaluated).

\subsection{The Galerkin solver remains essential}
The learned models do not replace the deterministic method. The Galerkin
solver defines the trial space and the discrete problem, supplies the
training data, operators and stored forcing, and performs correction and
fallback. A certificate can be attached to each prediction without a
second, learned error model because all of these objects are shared.

\subsection{Speed and certification cost}
At the production resolution, the neural predict--certify--correct
workflow is slower than an optimised free dense direct solve of the same
constrained problem: $1.18\times$ with cold setup and $1.58\times$ with
reusable setup, with every method verified on all 120 queries. The raw
DeepONet prediction never meets the $10^{-3}$ target, so every query is
corrected, and warm-starting the correction saves no time. The withdrawn
historical speed ratios compared against an unoptimised reference with
different setup accounting. The repaired rank-48 POD model is the only
method faster than dense LU ($0.78\times$ cold, $0.30\times$ with reuse) and
is accepted directly on every query; at rank 32 its initial error is
slightly too large and correction removes the advantage.

Certification and operator assembly, not the solve, dominate the cost: at
$N=64$ they account for most of the $0.84$ s cold dense-LU time, and up to
$N=512$ the time-stepping solve is a few percent of the total while the
certificate grows faster than linearly once the PCG cap is reached. Any
practical acceleration must therefore reduce assembly and certification
costs, for example by reusing parameter-independent operator data or by
cheaper lower bounds on fine meshes, before the speed of the predictor can
matter.

\subsection{Parameter maps and the restricted payoff family}
The fixed-strike model learns a parameter-to-trajectory map in
$(C,G,M,Y,r,T)$. The varying-strike model adds a one-dimensional manifold of
put payoffs. This is legitimate functional-input variation, but homogeneity
\eqref{eq:strike-homogeneity} makes it a weak test of payoff-shape learning.
Broader claims require matched training budgets and richer payoff inputs,
with consistent exterior data and initial-error treatment.

\subsection{Limitations and extensions}
The study is restricted to European puts, $Y\in(1,2)$, positive rates in
the reported experiments, a fixed family of truncated domains and the
parameter box of Table~\ref{tab:proposed-box}. The neural and POD models
are $N=64$ artefacts; resolution transfer of the certified workflow was not
studied. The domain sensitivity of the discrete price at $h=0.125$ is not
resolved, and the corner discrepancy shows that the production mesh is too
coarse for short-maturity, weakly diffusive parameters if continuum prices
are the goal. Validated arithmetic would be required for machine-enclosed
guarantees. Assembling the exact real-line Galerkin operator, or proving
the certificate hypotheses for the principal-period convention, would
connect levels (i) and (ii). The inverse-inequality bounds for Delta and
Gamma, certificate-driven enrichment and synthetic calibration are not
tested numerically. Further extensions include broader payoff families,
American-option obstacle problems and inexpensive certified lower bounds on
fine meshes.

% =================================================================
\section{Conclusion}
\label{sec:conclusion}

We formulated and implemented algebraic residual certification for learned
and reduced approximations of CGMY pricing trajectories, posed in the
constrained free coordinates of the production discretisation. The
Crank--Nicolson bound, augmented by the initial error, and the inexact-Riesz
upper bracket certify the exact coefficients that are returned, relative to
the exact solution of the assembled discrete problem. On the v5 evidence,
the free positivity gate passes at all 459 audit coordinates, where the
legacy full-space gate fails at 154; no bracket violation is detected
among 11{,}760 comparisons with a 40-digit reference (8{,}065 resolved,
3{,}695 unresolved at roundoff); and the augmented bound holds in all 144
tested candidates, including the 48 for which the residual alone
underestimates the error.

The practical conclusions are more limited. Certification and operator
assembly dominate the cost of every method. The DeepONet workflow is
correct but slower than a dense direct solve at $N=64$ ($1.18\times$ cold,
$1.58\times$ with reusable setup), whereas a correctly constrained
mass-weighted POD model at rank 48 is accepted directly and is faster
($0.78\times$ and $0.30\times$). Price bounds relative to the
discrete problem are of order $10^{-6}$ dollars after correction, while the
discrete price itself differs from independent continuum prices by a median
of \$0.03 at the production resolution. Certification therefore controls
the surrogate error, not the discretisation error, and the latter is
assessed here only empirically through refinement.

% =================================================================
\appendix
\section{Proofs}
\label{app:proofs}
% =================================================================

\subsection{Proof of Theorem~\ref{thm:uniform-coercivity}}
\label{app:proof-uniform}

Throughout, $\mu=(C,G,M,Y,\rf)\in\Pset$ and $s_\mu=Y/2\in(1/2,1)$.  We work on
the full line and transfer to $\Omega$ at the end using
Lemma~\ref{lem:conformity} and the zero-extension convention
\eqref{eq:energy-space}, which preserve all constants.

\medskip\noindent
\textbf{Step 1: uniform equivalence of the parameter-dependent norm.}
Let $s\in[Y_{\min}/2,Y_{\max}/2]\subset(1/2,1)$ and let
$w_\mu(\xi):=1+|M-i\xi|^{2s}+|G+i\xi|^{2s}
=1+(M^2+\xi^2)^{s}+(G^2+\xi^2)^{s}$ be the weight in
\eqref{eq:HMG-norm}.  Since $s\le1$ and $(a+b)^s\le a^s+b^s$ for $s\in(0,1]$,
\[
  w_\mu(\xi)\ \le\ 1+2\max\{1,M_{\max}^2,G_{\max}^2\}^{s}(1+\xi^2)^{s}
  \ \le\ (1+2\Theta)(1+\xi^2)^{s},
\]
with $\Theta=\max\{1,M_{\max}^2,G_{\max}^2\}^{Y_{\max}/2}$, using
$(1+\xi^2)^s\ge1$.  For the lower bound, $w_\mu(\xi)\ge1+|\xi|^{2s}$ and,
since $(1+\xi^2)^s\le2^s\max\{1,|\xi|^{2s}\}\le2(1+|\xi|^{2s})$ for
$s\le1$,
\[
  w_\mu(\xi)\ \ge\ \tfrac12(1+\xi^2)^{s}.
\]
Hence, uniformly in $\mu\in\Pset$,
\begin{equation}
  \tfrac12\norm{u}_{H^{s_\mu}(\R)}^2
  \ \le\ \norm{u}_{H^{s_\mu}_{M,G}(\R)}^2
  \ \le\ (1+2\Theta)\norm{u}_{H^{s_\mu}(\R)}^2 ,
  \label{eq:uniform-norm-equivalence}
\end{equation}
where $\norm{u}_{H^{s}(\R)}^2=\tfrac1{2\pi}\int(1+\xi^2)^{s}|\widehat u|^2$.

\medskip\noindent
\textbf{Step 2: uniform continuity.}
From \eqref{eq:symbol}, for every $\xi$,
\[
  |A_\mu(\xi)|
  \le|\rf-\omega_S||\xi|
   +C\Gamma(-Y)\bigl[(M^2+\xi^2)^{Y/2}+(G^2+\xi^2)^{Y/2}\bigr]
   +\lambda_0 .
\]
The map $(C,G,M,Y)\mapsto\omega_S$ of \eqref{eq:martingale-correction} is
continuous on $\Pset$ because $(M-1)^Y$ is continuous and finite on the selected box
(and remains so at $M=1$), and $\Gamma(-Y)$ is continuous on $[Y_{\min},Y_{\max}]\subset(1,2)$;
hence $\Omega_\ast<\infty$ and $\lambda_\ast<\infty$.  Since $Y>1$ we have
$|\xi|\le(1+\xi^2)^{Y/2}$, and
$(a^2+\xi^2)^{Y/2}\le\max\{1,a^2\}^{Y/2}(1+\xi^2)^{Y/2}$.  Therefore
\begin{equation}
  |A_\mu(\xi)|
  \le\Bigl(\Omega_\ast+2C_{\max}\gamma_{\max}\max\{1,a_\ast^2\}^{Y_{\max}/2}
   +\lambda_\ast\Bigr)(1+\xi^2)^{s_\mu}.
  \label{eq:uniform-symbol-upper}
\end{equation}
By Plancherel's theorem and the Cauchy--Schwarz inequality,
\[
  |a_\R(u,w)|
  \le\Bigl(\sup_{\xi\in\R}\frac{|A_\mu(\xi)|}{(1+\xi^2)^{s_\mu}}\Bigr)
  \norm{u}_{H^{s_\mu}(\R)}\norm{w}_{H^{s_\mu}(\R)},
\]
and
\eqref{eq:uniform-norm-equivalence} converts the $H^{s_\mu}$ norms into
$H^{s_\mu}_{M,G}$ norms at the cost of a factor $2$.  This gives
\eqref{eq:Ccont-explicit}.

\medskip\noindent
\textbf{Step 3: the zero-frequency and low-frequency region.}
By the exact identity \eqref{eq:symbol-real-exact}, $\Gamma(-Y)>0$ for
$Y\in(1,2)$, and the monotonicity $\Real(a\pm i\xi)^Y\le a^Y$ of
\cite{DamircheliCGMY2026}, both bracketed terms in
\eqref{eq:symbol-real-exact} are nonnegative, so
\begin{equation}
  \Real A_\mu(\xi)\ \ge\ \rf\ \ge\ r_{\min}>0
  \qquad\text{for all }\xi\in\R,\ \mu\in\Pset .
  \label{eq:uniform-low-freq}
\end{equation}
This is where $r_{\min}>0$ is used, and it is the only place: at $\xi=0$ the
fractional contribution vanishes identically.

\medskip\noindent
\textbf{Step 4: a uniform high-frequency lower bound.}
Write $\Real(a\pm i\xi)^Y=(a^2+\xi^2)^{Y/2}\cos\bigl(Y\vartheta_a(\xi)\bigr)$
with $\vartheta_a(\xi)=\arctan(|\xi|/a)\in[0,\pi/2)$.  Recall
$\delta_\ast=\tfrac\pi4(Y_{\min}-1)>0$ and
$\vartheta_\ast:=(\pi/2+\delta_\ast)/Y_{\min}$.  Since
$Y_{\min}>1$ we have
\[
  \vartheta_\ast
  =\frac{\pi/2+\tfrac\pi4(Y_{\min}-1)}{Y_{\min}}
  =\frac{\pi}{4}\cdot\frac{Y_{\min}+1}{Y_{\min}}
  <\frac{\pi}{2},
\]
because $Y_{\min}+1<2Y_{\min}$.  Hence $R_\ast=a_\ast\tan\vartheta_\ast$ is
finite and positive.  Let $|\xi|\ge R_\ast$ and let $a\in\{M,G\}$, so that
$a\le a_\ast$ and therefore
$\vartheta_a(\xi)=\arctan(|\xi|/a)\ge\arctan(R_\ast/a_\ast)=\vartheta_\ast$.
Consequently
\[
  Y\vartheta_a(\xi)\ \ge\ Y_{\min}\vartheta_\ast=\frac\pi2+\delta_\ast,
  \qquad
  Y\vartheta_a(\xi)\ <\ Y\cdot\frac\pi2\ <\ \pi ,
\]
the last inequality because $Y<2$.  As $\cos$ is decreasing on
$[\pi/2,\pi]$,
\[
  \cos\bigl(Y\vartheta_a(\xi)\bigr)\ \le\ \cos\Bigl(\frac\pi2+\delta_\ast\Bigr)
  =-\sin\delta_\ast\ <\ 0 .
\]
Therefore, for $|\xi|\ge R_\ast$,
$\Real\bigl[(M-i\xi)^Y+(G+i\xi)^Y\bigr]\le-2\sin(\delta_\ast)|\xi|^Y$, using
$(a^2+\xi^2)^{Y/2}\ge|\xi|^Y$.  Since the drift term in \eqref{eq:symbol} is
purely imaginary and $\lambda_0>0$,
\begin{equation}
  \Real A_\mu(\xi)
  \ \ge\ 2C\Gamma(-Y)\sin(\delta_\ast)|\xi|^Y
  \ \ge\ \kappa_\ast|\xi|^{Y},
  \qquad |\xi|\ge R_\ast,
  \label{eq:uniform-high-freq}
\end{equation}
with $\kappa_\ast:=2C_{\min}\gamma_{\min}\sin\delta_\ast>0$, where
$\gamma_{\min}=\min_{[Y_{\min},Y_{\max}]}\Gamma(-Y)>0$ by continuity and
positivity of $\Gamma(-\cdot)$ on $(1,2)$.  Note that $\delta_\ast$, and hence
$\kappa_\ast$, degenerate as $Y_{\min}\downarrow1$; this is the analytic
content of the requirement $Y_{\min}>1$.

\medskip\noindent
\textbf{Step 5: combination.}
For $|\xi|\le R_\ast$, \eqref{eq:uniform-low-freq} and
$(1+\xi^2)^{s_\mu}\le(1+R_\ast^2)^{Y_{\max}/2}$ give
$\Real A_\mu(\xi)\ge r_{\min}(1+R_\ast^2)^{-Y_{\max}/2}(1+\xi^2)^{s_\mu}$.
For $|\xi|\ge R_\ast$, since $t\mapsto t^2/(1+t^2)$ is increasing and
$\bigl(R_\ast^2/(1+R_\ast^2)\bigr)^{Y/2}
\ge\bigl(R_\ast^2/(1+R_\ast^2)\bigr)^{Y_{\max}/2}$,
\[
  \frac{|\xi|^{Y}}{(1+\xi^2)^{Y/2}}
  =\Bigl(\frac{\xi^2}{1+\xi^2}\Bigr)^{Y/2}
  \ \ge\ \Bigl(\frac{R_\ast^2}{1+R_\ast^2}\Bigr)^{Y_{\max}/2},
\]
so \eqref{eq:uniform-high-freq} gives
$\Real A_\mu(\xi)\ge\kappa_\ast(R_\ast^2/(1+R_\ast^2))^{Y_{\max}/2}(1+\xi^2)^{s_\mu}$.
Taking the minimum of the two prefactors and applying Plancherel's theorem,
\[
  \Real a_\R(u,u)\ \ge\ \tilde c_1\norm{u}_{H^{s_\mu}(\R)}^2,
  \qquad
  \tilde c_1
  =\min\Bigl\{\tfrac{r_{\min}}{(1+R_\ast^2)^{Y_{\max}/2}},\
   \kappa_\ast\bigl(\tfrac{R_\ast^2}{1+R_\ast^2}\bigr)^{Y_{\max}/2}\Bigr\},
\]
and \eqref{eq:uniform-norm-equivalence} converts this into the
$H^{s_\mu}_{M,G}$ norm at the cost of the factor $(1+2\Theta)^{-1}$, giving
\eqref{eq:c1-explicit}.

\medskip\noindent
\textbf{Step 6: transfer to $\Omega$.}
For $v,w\in\Vmu$ the zero extensions satisfy
$\widetilde v,\widetilde w\in H^{s_\mu}(\R)$ with
$\norm{\widetilde v}_{L^2(\R)}=\norm{v}_{L^2(\Omega)}$, and by definition
$a_\mu(v,w)=a_\R(\widetilde v,\widetilde w)$ and
$\norm{v}_{\Vmu}=\norm{\widetilde v}_{H^{s_\mu}_{M,G}(\R)}$.  Both estimates
therefore transfer with identical constants, independent of $h$ and of
$\Omega$.  Finally, for real-valued $v$ the form $a_\mu(v,v)$ is real, so
$a_\mu(v,v)=\Real a_\mu(v,v)$.  \qed

\begin{remark}[Sharpness]
\label{rem:sharpness}
No claim of sharpness is made for
\eqref{eq:c1-explicit}--\eqref{eq:Ccont-explicit}.  The choice
$\delta_\ast=\tfrac\pi4(Y_{\min}-1)$ is one convenient admissible value;
optimising $\delta_\ast$ trades the size of $R_\ast$ against $\sin\delta_\ast$
and can be done numerically for a given box.  Since the constants enter the
paper only through the efficiency bound
\eqref{eq:efficiency}, and since $\kappa(\mu)$ is directly computable
(Remark~\ref{rem:kappa-computable}), the explicit constants serve to
establish uniformity rather than to be used numerically.
\end{remark}


% =================================================================
\section{Historical figures}
\label{app:historical-figures}

The two figures below are retained unchanged from the earlier version of
this paper so that all nine original figures remain available. Their speed
claims are superseded and must not be cited: they compare a learned
workflow with shared operator caches against an unoptimised reference that
built its own setup, at different accuracy criteria, under the legacy
full-space certificate. Figures~\ref{fig:benchmark-ratios}
and~\ref{fig:scaling} replace them.

\begin{figure}[ht]
\centering
\includegraphics[width=0.85\linewidth]{exp11_tolerance_sweep.png}
\caption{Historical tolerance sweep (legacy certificate, 120 queries per
tolerance). Superseded: the runtime and speedup panels (top right, bottom
left) are withdrawn; see Section~\ref{ssec:benchmark}. The decision-branch
and error/certificate panels correspond to Table~\ref{tab:tolerance}.}
\label{fig:tolerance-sweep}
\end{figure}

\begin{figure}[ht]
\centering
\includegraphics[width=0.75\linewidth]{exp15_runtime_scaling_matrix.png}
\caption{Historical runtime scaling of the reference solve and the
forced-correction projected-FNO workflow on 30 common queries certifiable
at every resolution. Superseded: the gap between the reference and the
workflow reflects reference setup costs, not neural acceleration; see
Section~\ref{ssec:scaling-results}.}
\label{fig:runtime-scaling}
\end{figure}

% =================================================================
\section{Evidence and reproducibility}
\label{app:evidence}

All v5 results are stored under \texttt{results/free\_integrated\_v5/}
with a manifest per run recording source and configuration hashes, git
state, model checkpoints and bases, query-population hashes, software,
hardware, thread settings and seeds. Table~\ref{tab:runs} lists the runs
used here. Every table in
Sections~\ref{ssec:audit-results}--\ref{ssec:fourier-financial} and
Figures~\ref{fig:benchmark-ratios} and~\ref{fig:scaling} are generated by
\texttt{scripts/v6\_build\_manuscript\_assets.py} from the raw records of
these runs; the generator writes \texttt{evidence\_v6.json} with the SHA-256
hash of each of its 32 raw inputs and of every generated asset. The complete
test suite (290 tests) passes; its output is archived with the M9 runs.

\begin{table}[ht]
\centering\small
\caption{Runs behind the v5 evidence.}
\label{tab:runs}
\begin{tabular}{ll}
\toprule
Study & run identifier\\
\midrule
POD rank sweep (M3) & \texttt{20261006T215800Z\_m3\_pod\_study}\\
Gate G3 checks (M4) & \texttt{20261006T220619Z\_m4\_g3}\\
Positivity audit, 459 queries & \texttt{20261007T175335Z\_m6\_audit459}\\
Riesz-bracket validation & \texttt{20261007T175335Z\_m6\_riesz\_bracket}\\
Benchmark, cold mode & \texttt{20261007T175406Z\_gate\_d\_v5\_full\_cold}\\
Benchmark, reuse mode & \texttt{20261007T175334Z\_gate\_d\_v5\_full}\\
Resolution study & \texttt{20261007T175334Z\_scaling\_v5}\\
Corner refinement & \texttt{20261007T175334Z\_corner\_refinement}\\
Assembly convention check & \texttt{20261008T141446Z\_m8\_coefficient\_convention}\\
Independent assembly & \texttt{20261008T150219Z\_m8\_independent\_assembly}\\
Price refinement & \texttt{20261008T150220Z\_m8\_refinement}\\
Fourier reruns & \texttt{20261008T150600Z\_m9\_fourier}\\
Financial budgets & \texttt{20261008T152124Z\_m9\_financial}\\
\bottomrule
\end{tabular}
\end{table}

% =================================================================
% =================================================================
% Bibliography
% =================================================================

\begin{thebibliography}{99}

\bibitem{BhattacharyaEtAl2021}
K.~Bhattacharya, B.~Hosseini, N.~B. Kovachki, and A.~M. Stuart,
\newblock Model reduction and neural networks for parametric PDEs,
\newblock \emph{SMAI Journal of Computational Mathematics}, 7 (2021),
pp.~121--157.

\bibitem{CarrGemanMadanYor2002}
P.~Carr, H.~Geman, D.~B. Madan, and M.~Yor,
\newblock The fine structure of asset returns: an empirical investigation,
\newblock \emph{Journal of Business}, 75 (2002), pp.~305--332.

\bibitem{CarteadelCastilloNegrete2007}
\'A.~Cartea and D.~del Castillo-Negrete,
\newblock Fractional diffusion models of option prices in markets with jumps,
\newblock \emph{Physica A}, 374 (2007), pp.~749--763.

\bibitem{ContLantosPironneau2011}
R.~Cont, N.~Lantos, and O.~Pironneau,
\newblock A reduced basis for option pricing,
\newblock \emph{SIAM Journal on Financial Mathematics}, 2 (2011),
pp.~287--316. DOI: \url{https://doi.org/10.1137/10079851X}.

\bibitem{ContVoltchkova2005}
R.~Cont and E.~Voltchkova,
\newblock A finite difference scheme for option pricing in jump diffusion and
exponential L\'evy models,
\newblock \emph{SIAM Journal on Numerical Analysis}, 43 (2005),
pp.~1596--1626.

\bibitem{DamircheliCGMY2026}
D.~Damircheli,
\newblock A B-spline Galerkin method for option pricing under the CGMY
tempered stable process,
\newblock submitted manuscript, 2026.

\bibitem{FangOosterlee2008}
F.~Fang and C.~W. Oosterlee,
\newblock A novel pricing method for European options based on Fourier-cosine
series expansions,
\newblock \emph{SIAM Journal on Scientific Computing}, 31 (2008),
pp.~826--848.

\bibitem{GreplPatera2005}
M.~A. Grepl and A.~T. Patera,
\newblock A posteriori error bounds for reduced-basis approximations of
parametrized parabolic partial differential equations,
\newblock \emph{ESAIM: Mathematical Modelling and Numerical Analysis},
39 (2005), pp.~157--181.

\bibitem{HaasdonkOhlberger2008}
B.~Haasdonk and M.~Ohlberger,
\newblock Reduced basis method for finite volume approximations of
parametrized linear evolution equations,
\newblock \emph{ESAIM: Mathematical Modelling and Numerical Analysis},
42 (2008), pp.~277--302.

\bibitem{HesthavenUbbiali2018}
J.~S. Hesthaven and S.~Ubbiali,
\newblock Non-intrusive reduced order modeling of nonlinear problems using
neural networks,
\newblock \emph{Journal of Computational Physics}, 363 (2018), pp.~55--78.

\bibitem{HilberReichmannSchwabWinter2013}
N.~Hilber, O.~Reichmann, C.~Schwab, and C.~Winter,
\newblock \emph{Computational Methods for Quantitative Finance: Finite Element
Methods for Derivative Pricing},
\newblock Springer, Berlin, 2013.

\bibitem{JunikePankrashkin2022}
G.~Junike and K.~Pankrashkin,
\newblock Precise option pricing by the COS method---how to choose the
truncation range,
\newblock \emph{Applied Mathematics and Computation}, 421 (2022), 126935.
DOI: \url{https://doi.org/10.1016/j.amc.2022.126935}.

\bibitem{KovachkiEtAl2023}
N.~Kovachki, Z.~Li, B.~Liu, K.~Azizzadenesheli, K.~Bhattacharya, A.~Stuart,
and A.~Anandkumar,
\newblock Neural operator: learning maps between function spaces with
applications to PDEs,
\newblock \emph{Journal of Machine Learning Research}, 24 (2023), pp.~1--97.

\bibitem{LiKovachkiEtAl2021}
Z.~Li, N.~Kovachki, K.~Azizzadenesheli, B.~Liu, K.~Bhattacharya, A.~Stuart,
and A.~Anandkumar,
\newblock Fourier neural operator for parametric partial differential
equations,
\newblock International Conference on Learning Representations, 2021.

\bibitem{LiZhengEtAl2021}
Z.~Li, H.~Zheng, N.~Kovachki, D.~Jin, H.~Chen, B.~Liu, K.~Azizzadenesheli,
and A.~Anandkumar,
\newblock Physics-informed neural operator for learning partial differential
equations,
\newblock arXiv:2111.03794, 2021.

\bibitem{LuJinPangZhangKarniadakis2021}
L.~Lu, P.~Jin, G.~Pang, Z.~Zhang, and G.~E. Karniadakis,
\newblock Learning nonlinear operators via DeepONet based on the universal
approximation theorem of operators,
\newblock \emph{Nature Machine Intelligence}, 3 (2021), pp.~218--229.

\bibitem{MatacheVonPetersdorffSchwab2004}
A.-M. Matache, T.~von Petersdorff, and C.~Schwab,
\newblock Fast deterministic pricing of options on L\'evy driven assets,
\newblock \emph{ESAIM: Mathematical Modelling and Numerical Analysis},
38 (2004), pp.~37--71.

\bibitem{QiuDahmenChen2026}
Y.~Qiu, W.~Dahmen, and P.~Chen,
\newblock Variationally correct operator learning: reduced basis neural
operator with a posteriori error estimation,
\newblock \emph{Computer Methods in Applied Mechanics and Engineering},
461, Part B (2026), 119228.
DOI: \url{https://doi.org/10.1016/j.cma.2026.119228}.

\bibitem{Rannacher1984}
R.~Rannacher,
\newblock Finite element solution of diffusion problems with irregular data,
\newblock \emph{Numerische Mathematik}, 43 (1984), pp.~309--327.

\bibitem{RozzaHuynhPatera2008}
G.~Rozza, D.~B.~P. Huynh, and A.~T. Patera,
\newblock Reduced basis approximation and a posteriori error estimation for
affinely parametrized elliptic coercive partial differential equations,
\newblock \emph{Archives of Computational Methods in Engineering},
15 (2008), pp.~229--275.

\bibitem{YouYuDEliaGaoSilling2022}
H.~You, Y.~Yu, M.~D'Elia, T.~Gao, and S.~Silling,
\newblock Nonlocal kernel network (NKN): a stable and resolution-independent
deep neural network,
\newblock \emph{Journal of Computational Physics}, 469 (2022), 111536.

\end{thebibliography}

\end{document}
